\documentclass[11pt]{report}

\usepackage[T1]{fontenc}
\usepackage[utf8]{inputenc}
\usepackage{lmodern}
\usepackage[margin=1in]{geometry}
\usepackage{amsmath,amssymb,amsthm,mathtools}
\usepackage{xcolor}
\usepackage{booktabs}
\usepackage{longtable}
\usepackage{enumitem}
\usepackage{hyperref}
\usepackage{listings}

\setlength{\emergencystretch}{3em}
\setlist[itemize]{leftmargin=2em}
\setlist[enumerate]{leftmargin=2em}

\hypersetup{
  colorlinks=true,
  linkcolor=blue!55!black,
  urlcolor=blue!55!black,
  citecolor=blue!55!black
}

\lstdefinelanguage{EasyCrypt}{
  morekeywords={
    require,import,theory,end,section,lemma,proof,qed,op,pred,type,module,
    module,include,declare,clone,proc,var,if,then,else,while,return,
    forall,exists,real,int,bool,unit,phoare,hoare,equiv,byequiv,byphoare,
    by,apply,rewrite,move,case,smt,lossless,islossless
  },
  sensitive=true,
  morecomment=[l]{//},
  morecomment=[s]{(*}{*)},
  morestring=[b]"
}

\lstdefinelanguage{bash}{
  morekeywords={cd,sh,make,test,for,do,done,if,then,else,fi,echo,sed,tail,grep,rg},
  sensitive=true,
  morecomment=[l]{\#},
  morestring=[b]"
}

\lstset{
  basicstyle=\ttfamily\small,
  breaklines=true,
  columns=fullflexible,
  frame=single,
  framerule=0.3pt,
  rulecolor=\color{black!30},
  xleftmargin=0.5em,
  xrightmargin=0.5em,
  aboveskip=0.7em,
  belowskip=0.7em
}

\DeclareRobustCommand{\file}[1]{\ifmmode\text{\nolinkurl{#1}}\else\nolinkurl{#1}\fi}
\DeclareRobustCommand{\ec}[1]{\ifmmode\text{\nolinkurl{#1}}\else\nolinkurl{#1}\fi}
\newcommand{\HAETAE}{\mathsf{HAETAE}}
\newcommand{\ROM}{\mathsf{ROM}}
\newcommand{\ECModel}{\mathsf{ECModel}}
\newcommand{\Spec}{\mathsf{Spec}_{260204}}
\newcommand{\EUFCMA}{\mathsf{EUF\mbox{-}CMA}}
\newcommand{\UFNMA}{\mathsf{UF\mbox{-}NMA}}
\newcommand{\MLWE}{\mathsf{MLWE}}
\newcommand{\MSIS}{\mathsf{Module\mbox{-}SIS}}
\newcommand{\BMSIS}{\mathsf{BimodalSelfTargetMSIS}}
\newcommand{\Adv}{\operatorname{Adv}}
\newcommand{\Prb}{\Pr}
\newcommand{\SubDist}{\mathsf{SubDist}}
\newcommand{\Dist}{\mathsf{Dist}}

\theoremstyle{definition}
\newtheorem{definition}{Definition}[chapter]
\newtheorem{invariant}{Invariant}[chapter]
\newtheorem{obligation}{Obligation}[chapter]
\newtheorem{checkpoint}{Checkpoint}[chapter]
\theoremstyle{plain}
\newtheorem{theorem}{Theorem}[chapter]
\newtheorem{lemma}{Lemma}[chapter]

\title{Source-Guided Roadmap for Reimplementing and Machine-Checking the HAETAE Provable-Security Proof}
\author{Technical report for the HAETAE EasyCrypt proof surface}
\date{May 2026}

\begin{document}
\pagenumbering{roman}
\maketitle

\begin{abstract}
This report is a repository-guided roadmap for reimplementing the HAETAE
provable-security proof surface.  It is written for a reader who has not seen
the repository before and wants to understand the pen-and-paper proof boundary,
the EasyCrypt file architecture, and the verification workflow.  It is not, by
itself, a literal replacement for the \file{.ec} source files: several blueprint
snippets are schematic, and a faithful machine-checked reimplementation still
requires the existing EasyCrypt files, generated declaration inventory, and
EasyCrypt libraries as companion artifacts.  A claim of report-only
machine-checkable reproduction should therefore be rejected unless this roadmap
is expanded into a complete source-and-proof archive.  The report identifies the
exact theorem being checked, separates checked facts from external assumptions, gives
the source-file dependency order, lists the mathematical object families that
must be introduced, and records the verification commands that certify the proof
gate.  The document-only reproducibility question has a negative answer: the
report can explain the checked formal proof boundary, but it cannot by itself
establish HAETAE algorithm security or reproduce the EasyCrypt verification.
The checked theorem is a conditional random-oracle-model statement for the
formal EasyCrypt model.  It is not a standard-model SHAKE theorem, not an
implementation-correctness theorem, and not a non-vacuous concrete security
instantiation until the currently oversized rejection-sampling antecedent is
repaired.  When full \file{.ec} proof scripts are provided as companion
artifacts, this report also states the documentation standard needed for a
mathematician or cryptologist to understand the mathematical content of those
scripts without relying on unstated repository lore.  This version records that
standard and an initial coverage ledger; it should not be described as complete
script-companion documentation until every required coverage row is discharged.
\end{abstract}

\tableofcontents
\clearpage
\pagenumbering{arabic}

\chapter{Scope and End Product}

\section{What the reader should be able to reproduce}

After following this report together with the companion EasyCrypt sources, a
reader should be able to recreate a proof surface with the following
properties.

\begin{enumerate}
\item The EasyCrypt files compile as a manifest-driven proof gate.
\item The generic random-oracle signature game is instantiated by a formal
HAETAE scheme module.
\item Correctness is proved for the formal model.
\item The EUF-CMA game is related to a sequence of simulated, transcript, NMA,
sampling, and extraction games.
\item The final checked security implication has the same dependency boundary
as the current artifact: it is conditional on explicit hop/reduction premises
and abstract MLWE/Module-SIS hardness bounds.
\item The counted-ROM loss ledger is normalized to the concrete expression
\[
  \frac{37}{2^{58}}
\]
for the fixed wrapper constants \(q_H=q_S=2\).
\item The proof explicitly records the current vacuity issue:
\(\ec{rejection\_sampling\_loss\_term}\ge 1\), and in fact the
min-entropy part alone is \(12\), so a final premise that adds this term to a
probability cannot be treated as a useful concrete-security premise.
\end{enumerate}

\section{Sufficiency status}

This report should be read as a technical map and acceptance checklist, not as a
complete standalone source of the proof.  It is sufficient for understanding the
shape and limits of the checked theorem, for navigating the proof files in the
right order, and for checking that a reconstructed development has the expected
manifest, theorem statement, and loss ledger.  It is not sufficient, on its own,
to type all EasyCrypt code from a blank file and recover the current proof.

The reason is structural rather than cosmetic.  A machine-checked EasyCrypt
development depends on exact declarations, module restriction sets, local helper
lemmas, proof tactics, import paths, and library behavior.  This report records
the required proof architecture and representative source forms, but it does not
inline all 22046 lines of checked EasyCrypt.  Whenever a snippet contains
\ec{...}, it is a blueprint placeholder, not compilable code.  A reader
attempting a literal reimplementation must replace every such placeholder by the
corresponding complete definition or proof from the companion artifacts.

\section{Document-only verdict}

The question ``Can a reader understand the provable security of the HAETAE
digital signature algorithm and reproduce the machine-checked verification
based solely on this technical document?'' has the following answer.

\begin{center}
\fbox{\begin{minipage}{0.92\linewidth}
\textbf{No.}  The document alone is enough to orient a reader to the checked
formal-model theorem, but it is not enough to certify the HAETAE algorithm as a
deployed digital signature scheme and it is not enough to recreate the
machine-checked EasyCrypt artifact from a blank workspace.
\end{minipage}}
\end{center}

There are three distinct claims, and only the weakest one is supported by the
report alone.

\begin{enumerate}
\item \textbf{Understanding the checked theorem boundary: yes.}  The report
states the final theorem shape, the random-oracle setting, the fixed counted-ROM
ledger, the declared hardness premises, and the rejection-sampling vacuity
warning.
\item \textbf{Understanding full HAETAE algorithm security: no.}  The report
does not supply all bridges from the formal EasyCrypt model to the specification,
SHAKE instantiation, concrete lattice estimates, implementation behavior, or
side-channel properties.
\item \textbf{Reproducing machine-checked verification from the document alone:
no.}  The report omits literal EasyCrypt source, complete proof scripts, exact
module restrictions, full helper-lemma bodies, and a full import/library
contract.  Those are exactly the objects that the checker consumes.
\end{enumerate}

The strongest truthful claim is therefore conditional: with the companion
source tree and libraries, the report guides a reader through the proof surface
and the verification gate.  Without those artifacts, it is an audit roadmap, not
a standalone proof.

\section{Can this document stand alone?}

The criticism that motivated this revision is correct: the report, by itself,
cannot be both a readable roadmap and a complete machine-checkable source
archive.  The following table states exactly what can and cannot be recovered
from the document alone.

\begin{longtable}{p{0.27\linewidth}p{0.16\linewidth}p{0.47\linewidth}}
\toprule
\textbf{Goal} & \textbf{From this report alone?} & \textbf{Reason} \\
\midrule
Understand the final theorem boundary & Yes & The theorem shape, ROM ledger,
declared hardness premises, and vacuity warning are stated explicitly. \\
Reconstruct the pen-and-paper proof outline & Partly & The hybrid order and
loss accounting are given, but each step still has to be checked against the
source lemmas that implement it. \\
Type a fresh EasyCrypt development from a blank directory & No & The report
does not inline every declaration, helper lemma, proof script, import, and
module restriction set. \\
Machine-check equivalence to the current artifact & No & Equivalence requires
the companion source tree or an independently expanded document containing
literal source and proof scripts. \\
Conclude end-to-end HAETAE implementation security & No & The checked theorem
is formal-model, random-oracle-model, fixed-ledger, and conditional; byte-level,
constant-time, SHAKE-instantiation, and concrete hardness claims are outside
the checked result. \\
\bottomrule
\end{longtable}

Consequently, any phrase such as ``reproduce the machine-checked verification
based solely on this report'' must be read as false for the present artifact.
The accurate claim is narrower: a reader can use the report to understand the
proof boundary and to reproduce the verification when the report is paired with
the listed source artifacts.

\section{Formal-model understanding versus algorithm security}

The phrase ``the provable security of HAETAE'' is ambiguous.  In this report it
means the checked EasyCrypt theorem about the formal HAETAE model, not a full
security certification of every object that a cryptographer or implementer
might call the HAETAE digital signature algorithm.  A reader can learn the
formal proof boundary from this document: the game being bounded, the
random-oracle abstraction, the declared hardness games, the counted-ROM loss
ledger, and the final vacuity warning.  That is a real but narrower achievement
than understanding an end-to-end algorithm-security theorem.

To turn formal-model understanding into algorithm-security understanding, the
following bridges would also have to be supplied and checked or separately
audited.

\begin{itemize}
\item A specification-to-model bridge showing that the EasyCrypt operations
faithfully encode the HAETAE specification being claimed.
\item A hash-instantiation argument explaining why the random-oracle theorem is
applicable to the concrete SHAKE-based algorithm, or an explicit decision not
to make that claim.
\item Concrete parameter analysis for the MLWE and Module-SIS assumptions,
including any reduction losses not internalized by the EasyCrypt theorem.
\item Implementation-correctness and side-channel arguments for concrete code,
if the claim is about deployed signing and verification programs.
\item A repaired or reinterpreted rejection-sampling premise if the theorem is
to be cited as a non-vacuous concrete bound.
\end{itemize}

Because those bridges are not all present as checked theorems, the document
should be used to understand the formal verification surface, not as sole
evidence that the deployed HAETAE algorithm has been proved secure.

\section{Companion artifacts required}

A faithful reimplementation therefore requires at least the following companion
materials.

\begin{itemize}
\item The manifest files under \file{provable-security/easycrypt/}, or an
equivalent complete source tree.
\item The generic EasyCrypt libraries imported by those files.
\item The \file{kyber-security} include path, unless every imported dependency
from that path is audited and replaced.
\item The declaration inventory generated by
\file{dissertation/etc/count-easycrypt-declarations.sh}.
\item The HAETAE specification PDF and model/specification audit material if the
formal EasyCrypt model is to be interpreted as the February 2026 HAETAE
specification.
\end{itemize}

Without those artifacts, the report can support a mathematical review of the
proof boundary, but it cannot certify that a newly written EasyCrypt
development is definitionally identical to the checked one.

\section{Source-truth contract}

A faithful reproduction must declare its source of truth before any theorem is
accepted.

\begin{enumerate}
\item In \emph{audit mode}, the source of truth is the existing
\file{provable-security/easycrypt/} tree plus the manifest.  The report is a
guide for reading, checking, and rebuilding that tree.
\item In \emph{literal standalone mode}, the source of truth would have to be a
larger document that contains every active EasyCrypt file in full, every proof
script, every module restriction set, and every imported helper dependency that
affects type checking or proof checking.
\item In \emph{independent reimplementation mode}, the new development may
choose different local names or proof scripts, but it must provide a
traceability matrix from each manifest file, declaration family, theorem, and
loss term in this report to a checked EasyCrypt object in the new tree.
\end{enumerate}

The current report intentionally supports audit mode and independent
reimplementation mode with companion sources.  It does not satisfy literal
standalone mode.  Treating it as if it did would erase exactly the facts that
EasyCrypt checks: definitional equality, dependency resolution, tactic success,
and module-separation restrictions.

\section{Script companion mode}

There is a fourth, weaker-than-standalone but stronger-than-roadmap use case:
the \file{.ec} proof scripts are provided alongside this document, and the
document is responsible for making their mathematical content understandable to
mathematicians and cryptologists who have not previously studied the project.
This is \emph{script companion mode}.

Script companion mode does not mean that the prose replaces EasyCrypt.  The
scripts remain the machine-checkable artifact.  It means that a reader should be
able to open a script, look up every important declaration and proof block in
this report, and understand what mathematical claim is being checked, why the
claim is needed, which assumptions it uses, and what each proof step contributes.
The reader may still need EasyCrypt to check the proof, but should not need
unstated repository lore to understand the proof.

In script companion mode, the documentation must provide the following for each
manifest file.

\begin{longtable}{p{0.26\linewidth}p{0.64\linewidth}}
\toprule
\textbf{Documentation item} & \textbf{Required content} \\
\midrule
File role & The mathematical purpose of the file, its upstream dependencies, and
the later files that rely on it. \\
Declaration dictionary & Every type, operation, module type, module, predicate,
event, and distribution that is not standard EasyCrypt vocabulary, with a
mathematical reading. \\
State interpretation & For every game or oracle module, the meaning of each
state variable, log, counter, bad flag, and sampled value. \\
Restriction explanation & For every security-critical \ec{declare module}
restriction set, the information-flow reason for the exclusions. \\
Lemma ledger & A table of important lemmas and theorems giving the EasyCrypt
name, the informal mathematical statement, dependencies, and the later result
that consumes the lemma. \\
Proof-script commentary & A block-by-block explanation of the tactic script:
which invariant is opened, which equivalence or probability rule is applied,
which algebraic rewrite is used, and which side condition is delegated to
\ec{smt}. \\
Assumption boundary & A statement of whether a fact is checked locally,
imported from a library, declared as a hardness interface, or left as an
external cryptographic assumption. \\
\bottomrule
\end{longtable}

This requirement is intentionally stronger than merely listing theorem names.
For example, a line such as \ec{byphoare} or \ec{byequiv} must be explained by
the game transformation it establishes; a call to \ec{smt} must be associated
with the arithmetic, set-theoretic, or losslessness side condition it discharges;
and a module restriction list must be explained as an adversarial-separation
condition rather than as syntax.

\section{Per-lemma explanation template}

Every nontrivial EasyCrypt lemma used by the final theorem should have an entry
with the following fields.

\begin{enumerate}
\item \textbf{EasyCrypt anchor:} file name, lemma name, and whether the proof is
local, imported, or a direct composition of earlier lemmas.
\item \textbf{Mathematical statement:} a prose or displayed-math rendering that
does not require EasyCrypt syntax to parse.
\item \textbf{Objects involved:} the games, modules, distributions, events,
state variables, and loss terms appearing in the statement.
\item \textbf{Assumptions:} module restrictions, losslessness facts, positivity
facts, hardness premises, or previous hop lemmas consumed by the proof.
\item \textbf{Proof idea:} the mathematical argument before tactic details.
\item \textbf{Script reading:} an explanation of the proof script grouped into
meaningful blocks, including what is proved by rewriting, by probabilistic
program logic, by equivalence reasoning, and by SMT.
\item \textbf{Security role:} the place of the lemma in the global reduction:
exact hop, bad-event bound, sampling-distance bound, extraction step, hardness
adapter, or arithmetic normalization.
\item \textbf{Failure mode:} what would go wrong if the lemma were misstated,
weakened, or used outside its hypotheses.
\end{enumerate}

The goal is not to duplicate every tactic token in prose.  The goal is that a
mathematician or cryptologist can read the \file{.ec} script and this entry
together and recover the proof argument without asking what an identifier, game
state, or proof step was intended to mean.

\section{Minimum EasyCrypt vocabulary for readers}

When proof scripts are supplied, this report must also define the EasyCrypt
vocabulary used by the scripts at the level needed to read the proof.  At
minimum, it must explain:

\begin{itemize}
\item \ec{op}, \ec{pred}, \ec{type}, \ec{module type}, \ec{module}, and
\ec{declare module};
\item \ec{proc}, program state, sampling, procedure calls, and return values;
\item probability expressions \(\ec{Pr[... @ \&m : ...]}\);
\item Hoare, probabilistic Hoare, and equivalence judgments;
\item \ec{byphoare}, \ec{byequiv}, \ec{islossless}, \ec{lossless}, \ec{rewrite},
\ec{case}, \ec{move}, \ec{apply}, and \ec{smt};
\item the meaning of module restriction sets such as \ec{{-H, -A, -HAETAE}};
\item the difference between a checked lemma, a declared module interface, and
an external cryptographic assumption.
\end{itemize}

These explanations should be tied to the HAETAE scripts, not presented as a
generic EasyCrypt tutorial.  Each tactic or syntactic form should be introduced
where it first matters for a HAETAE proof step.

\section{Current script-companion coverage status}

The criticism that this document still only states the script-companion
standard is correct.  The current report contains a useful roadmap and several
mathematical readings, but it is not yet a complete commentary for every
security-relevant block of the supplied \file{.ec} scripts.  The table below is
therefore part of the technical specification: it records what remains before
the document can honestly claim that mathematicians and cryptologists can
understand the scripts' proof content solely through this report plus the
scripts themselves.

\begin{longtable}{p{0.25\linewidth}p{0.18\linewidth}p{0.47\linewidth}}
\toprule
\textbf{Coverage item} & \textbf{Current status} & \textbf{Required to close} \\
\midrule
File roles for manifest files & Partial & The file-by-file blueprint gives the
main role of each file, but it must be cross-checked against the actual
\file{.ec} scripts and expanded with upstream and downstream dependencies. \\
Declaration dictionary & Partial & Every nonstandard type, operation,
predicate, distribution, module type, module, and event in the scripts must have
a mathematical reading and script anchor. \\
State interpretation & Partial & Game modules, signing oracles, random-oracle
modules, transcript simulators, and samplers need complete state-variable
tables. \\
Restriction explanations & Partial & Representative restriction sets are
shown, but every security-critical \ec{declare module} restriction must be
explained by its information-flow purpose. \\
Lemma ledger & Partial & The report lists important lemma names, but each
nontrivial lemma used in the final theorem needs a full entry using the
per-lemma template. \\
Proof-script commentary & Not complete & Tactic scripts must be grouped into
mathematical proof blocks explaining \ec{byphoare}, \ec{byequiv}, rewrites,
losslessness arguments, and \ec{smt} side conditions. \\
EasyCrypt vocabulary tied to HAETAE & Partial & The report states the vocabulary
that must be explained, but the explanations must be integrated at the first
HAETAE proof point where each construct matters. \\
Assumption boundary & Partial & The major boundaries are recorded, but each
imported library fact, declared hardness interface, and external cryptographic
assumption must be labelled in the lemma ledger. \\
\bottomrule
\end{longtable}

Until the rows above are complete, the report should be cited as a
script-companion specification and partial reading guide, not as a completed
script-companion document.  The completion test is concrete: for every
security-relevant declaration and proof block in the supplied \file{.ec}
scripts, a reader must be able to find a nearby entry in this report explaining
the mathematical statement, the proof idea, the tactic role, and the place of
the block in the global security reduction.

\section{Coverage ledger template}

A complete script-companion version should maintain a ledger with one row per
security-relevant script block.  The ledger may be embedded in this report or
generated as an appendix, but it must be reviewable as prose.  Each row should
have the following shape.

\begin{longtable}{p{0.22\linewidth}p{0.68\linewidth}}
\toprule
\textbf{Field} & \textbf{Meaning} \\
\midrule
Script anchor & File name, line or lemma name, and whether the block is a
definition, module, declaration, lemma, proof, or tactic sub-block. \\
Mathematical reading & The human mathematical object or claim represented by
the script block. \\
Dependencies & Prior declarations, imported theories, assumptions, module
restrictions, or lemmas needed to read the block. \\
Proof explanation & For proof blocks, the argument split into exact
equivalence, probabilistic bound, invariant preservation, algebraic rewrite,
losslessness proof, extraction step, or arithmetic normalization. \\
Checker role & What EasyCrypt verifies at this block: type checking,
losslessness, Hoare judgment, equivalence judgment, probability inequality, or
SMT side condition. \\
Security role & Why the block matters for the HAETAE proof: correctness,
hybrid hop, random-oracle programming, transcript validity, rejection bound,
hardness adapter, or final theorem composition. \\
Status & Complete only when the prose is sufficient for a cryptographer to read
the script block without guessing the intended mathematics. \\
\bottomrule
\end{longtable}

This ledger is the mechanism that turns the current roadmap into a genuine
script companion.  A build of the \LaTeX{} document is not enough; the ledger
must also be substantively complete against the supplied scripts.

\section{Initial filled coverage entries}

The entries below are the first concrete script-companion rows.  They are
included to demonstrate the required level of explanation and to prevent the
coverage ledger from remaining a purely aspirational template.  The status is
``partial'' unless the row covers every dependent declaration and proof block in
the surrounding script region.

\begin{longtable}{p{0.20\linewidth}p{0.28\linewidth}p{0.42\linewidth}}
\toprule
\textbf{Script anchor} & \textbf{Mathematical reading} & \textbf{Current explanation and remaining work} \\
\midrule
\file{Sig\_ROM.ec}, types and freshness operations &
Generic signature-game vocabulary.  The abstract types
\ec{pkey}, \ec{skey}, \ec{message}, \ec{context}, \ec{signature},
\ec{ro\_query}, and \ec{ro\_output} define the carrier sets of a random-oracle
signature experiment.  The predicates \ec{fresh\_msg} and \ec{fresh\_sig}
express whether a forgery target was absent from the signing log. &
The report explains the EUF-CMA/NMA game boundary and freshness condition.
To close the row, add a declaration dictionary entry for every generic type,
the signed-query variant used by SUF-CMA, and the exact list operations used in
the freshness lemmas. \\
\midrule
\file{Sig\_ROM.ec}, \ec{EUF\_CMA} and \ec{UF\_NMA} modules &
\ec{EUF\_CMA} is the adaptive chosen-message experiment with a signing oracle
that records queried message-context pairs.  \ec{UF\_NMA} is the no-message
experiment in which the adversary only has random-oracle access. &
The mathematical game notation table identifies both games.  To close the row,
add state-variable tables for \ec{EUF\_CMA.sk}, \ec{EUF\_CMA.queries}, and the
local variables \ec{pk}, \ec{sk}, \ec{m}, \ec{ctx}, \ec{sig}, and \ec{ok}. \\
\midrule
\file{Sig\_ROM.ec}, \ec{nma\_as\_cma\_exact} &
The NMA adversary adapter is distribution-preserving when embedded into the
CMA game with an unused signing oracle: the two success probabilities are
equal. &
The proof is an EasyCrypt equivalence proof using \ec{byequiv}, inlining the
adapter and simulating matching calls to the adversary, scheme, and oracle.  To
close the row, add a block-by-block explanation of the \ec{call}, \ec{wp},
\ec{sim}, and final \ec{auto} steps. \\
\midrule
\file{HAETAE\_Reductions.ec}, hardness and loss operations &
The operations \ec{mlwe\_hardness\_term},
\ec{module\_sis\_hardness\_term}, and
\ec{bimodal\_to\_msis\_reduction\_loss\_term} are abstract real-valued
assumption/loss slots.  The counted-ROM terms encode collision, prequery, and
min-entropy losses for \(q_H=q_S=2\) and challenge support \(2^{58}\). &
The report gives the counted-ROM arithmetic and the \(37/2^{58}\) calculation.
To close the row, add dictionary entries for every real-valued operation in
the file and identify which are abstract assumptions, definitions, or checked
arithmetic normalizations. \\
\midrule
\file{HAETAE\_Reductions.ec}, \ec{haetae\_euf\_counted\_rom\_bound\_groupedE} &
This lemma rewrites the final counted-ROM bound into the sum of MLWE hardness,
Module-SIS hardness, bimodal-to-MSIS loss, and counted random-oracle
programming loss. &
The proof is pure real algebra: unfold the grouped definitions and discharge
the equality by the ring tactic.  To close the row, add similar explanation
entries for the nonnegativity, positivity, subunit, concrete-cardinality, and
vacuity lemmas in the same file. \\
\midrule
\file{HAETAE\_HopGames.ec}, section
\ec{ROMInternalTranscriptROSigningAttemptExactHyperballPaperSimNMA} &
This region relates an NMA game using a random-oracle signing-attempt sampler
to one using the exact-hyperball sampler, while tracking rejection and
sampler-prequery losses. &
The report describes the layer as NMA and transcript game lifting.  To close
the row, add proof-block commentary for the paper-sample lifting lemma, the
budgeted lifting lemmas, the clean-conditioned lifting lemma, and the bad
prequery bound that is imported into the inequality chain. \\
\midrule
\file{HAETAE\_HopGames.ec},
\ec{rom\_internal\_budgeted\_nma\_ro\_signing\_attempt\_ro\_exact\_hyperball\_loss\_from\_clean\_lifting} &
The lemma decomposes the success event into a clean branch and a sampler
prequery bad event.  It bounds the clean branch by the exact-hyperball game plus
rejection loss, and charges the bad branch to a counted-ROM prequery term. &
The script uses union/subdistribution inequalities, a clean-lifting premise,
and the lemma
\ec{budgeted\_paper\_sim\_sampler\_bad\_prequery\_nma\_fel\_bound}.  To close
the row, explain each \ec{ler\_trans}, \ec{Pr[mu\_sub]}, \ec{Pr[mu\_or]},
\ec{ler\_add}, and final bad-event-bound application. \\
\midrule
\file{HAETAE\_Security.ec}, section
\ec{MechanizedROMTranscriptConstructedNMALift} &
This section is the final composition lane.  It progressively replaces the
real EUF-CMA view by structural, public-simulation, paper-simulation,
rejection-simulation, exact-hyperball, and random-oracle exact-hyperball NMA
views before applying MLWE and bimodal/Module-SIS assumptions. &
The report states the final theorem shape and representative restriction sets.
To close the row, every lemma in this section must be entered in the ledger,
including its predecessor theorem, hop premise, hardness premise, arithmetic
bound, and whether the proof is a rewrite, transitivity step, or final
composition. \\
\midrule
\file{HAETAE\_Security.ec}, \ec{declare module H}, \ec{A}, \ec{Samp},
\ec{Bmlwe}, \ec{Bbimodal} in the final composition section &
The restrictions enforce adversarial separation: the random oracle and
adversary cannot call hidden game internals, the sampler cannot inspect the
adversary or oracle beyond its interface, and the hardness solvers remain
abstract modules. &
The report gives a representative restriction block and explains why
restriction sets are proof obligations.  To close the row, every excluded module
name in the final block must be mapped to the hidden state or procedure it
protects. \\
\bottomrule
\end{longtable}

\section{Script companion coverage for Sig ROM}

The generic signature-game file is small enough that its script-companion
coverage can be closed in this report.  The file is an abstract theory: it does
not know HAETAE-specific algebra, rejection sampling, or lattice assumptions.
It defines the reusable game vocabulary that later files instantiate with the
HAETAE scheme and random oracle.

\begin{longtable}{p{0.24\linewidth}p{0.66\linewidth}}
\toprule
\textbf{Script item} & \textbf{Mathematical reading} \\
\midrule
\ec{pkey}, \ec{skey}, \ec{message}, \ec{context}, \ec{signature} &
Abstract carrier sets for public keys, secret keys, signed messages, auxiliary
contexts, and signatures.  No algebraic structure is assumed here. \\
\ec{ro\_query}, \ec{ro\_output}, \ec{dro\_output} &
The input and output domains of the random oracle, together with a lossless
output distribution used by the cloned full-random-oracle library. \\
\ec{query}, \ec{signed\_query} &
The signing-log domains.  A \ec{query} records a message-context pair; a
\ec{signed\_query} records the same pair together with the signature returned by
the signing oracle. \\
\ec{fresh\_msg} &
EUF-CMA freshness: the message-context pair in the alleged forgery was not
previously submitted to the signing oracle. \\
\ec{fresh\_sig} &
Strong unforgeability freshness: the full message-context-signature triple was
not previously returned by the signing oracle. \\
\ec{Oracle} and \ec{POracle} &
\ec{Oracle} exposes random-oracle initialization and query access; \ec{POracle}
exposes query access only.  Adversaries receive \ec{POracle} so they cannot
reset the oracle. \\
\ec{Scheme(O)} &
The signature-scheme interface parameterized by public oracle access: key
generation, signing, and verification. \\
\ec{CORR\_ADV(O)} &
The correctness-game adversary, which chooses the message and context after
seeing both public and secret keys. \\
\ec{Adversary(H,O)} &
The EUF-CMA adversary.  Its procedure may call the random oracle and the signing
oracle. \\
\ec{NMA\_Adversary(H)} &
The no-message adversary.  It may call the random oracle but has no signing
oracle. \\
\ec{NMA\_As\_CMA} &
An adapter that lets an NMA adversary inhabit the CMA adversary interface while
ignoring the signing oracle. \\
\bottomrule
\end{longtable}

The game modules have the following state interpretation.

\begin{longtable}{p{0.24\linewidth}p{0.66\linewidth}}
\toprule
\textbf{Game state or local} & \textbf{Meaning} \\
\midrule
\ec{Correctness.main} &
Initializes the random oracle, samples a key pair, asks the correctness
adversary for a message and context, signs them, verifies the signature, and
returns failure \(\neg ok\).  Thus a zero probability result is perfect
correctness for the formal model. \\
\ec{EUF\_CMA.sk} &
The secret key stored by the EUF-CMA game so the internal signing oracle can
answer adversarial signing queries. \\
\ec{EUF\_CMA.queries} &
The list of message-context pairs queried to the signing oracle.  It is reset
to the empty list before the adversary runs. \\
\ec{EUF\_CMA.O.sign} &
The signing oracle: it signs with \ec{sk}, logs \((m,ctx)\), and returns the
signature. \\
\ec{EUF\_CMA.main} &
The adaptive unforgeability experiment.  It accepts iff verification succeeds
and \ec{fresh\_msg} holds for the forgery message-context pair. \\
\ec{UF\_NMA.main} &
The no-message experiment.  It samples keys, gives only random-oracle access to
the adversary, and accepts iff verification succeeds. \\
\ec{SUF\_CMA.main} &
The strong unforgeability analogue of EUF-CMA.  It logs returned signatures and
requires freshness of the full signed query.  It is defined for completeness of
the generic interface, although the HAETAE top-level theorem uses EUF-CMA/NMA. \\
\bottomrule
\end{longtable}

The two freshness lemmas are proof-script sanity checks.  \ec{fresh\_msg\_nil}
and \ec{fresh\_sig\_nil} unfold the corresponding freshness predicate and solve
membership in the empty list by simplification.  No cryptographic assumption is
used.

The lemma \ec{nma\_as\_cma\_exact} is the main proof block in this file.  Its
mathematical statement is
\[
  \Pr[\ec{EUF\_CMA}(H,S,\ec{NMA\_As\_CMA}(A))]=
  \Pr[\ec{UF\_NMA}(H,S,A)].
\]
The proof uses \ec{byequiv} to compare two programs in lockstep.  The pre- and
postcondition \(\ec{=\{glob H, glob S, glob A\} ==> =\{res\}}\) says that if
the oracle, scheme, and adversary global states agree initially, then the two
experiments return the same Boolean result.  The \ec{inline} step expands the
adapter so the proof can see that it simply calls the NMA adversary.  Each
\ec{call ... by sim} aligns matching calls to the adversary, scheme, and random
oracle.  The final \ec{auto => /> } discharges the remaining first-order
obligations, including that the unused signing oracle made no logged query.

This closes the script-companion coverage for \file{Sig\_ROM.ec} at the level
needed for the HAETAE proof: every declaration, game state, and proof block in
the generic interface has a mathematical reading.  Future additions would only
be needed if later proofs begin using \ec{SUF\_CMA} in a security theorem.

\section{Script companion coverage for HAETAE Reductions}

\file{HAETAE\_Reductions.ec} is the arithmetic ledger for the proof.  It does
not define a cryptographic game.  Instead, it defines the real-valued loss and
hardness terms that appear in the final theorem, then proves the algebraic and
order facts needed to compose those terms.

\begin{longtable}{p{0.30\linewidth}p{0.60\linewidth}}
\toprule
\textbf{Operation family} & \textbf{Mathematical reading} \\
\midrule
\ec{mlwe\_hardness\_term} &
Abstract upper bound for the success probability of the MLWE reduction module.
It is not estimated in this file. \\
\ec{module\_sis\_hardness\_term} &
Abstract upper bound for the Module-SIS reduction module.  It is also not a
concrete bit-security estimate. \\
\ec{bimodal\_to\_msis\_reduction\_loss\_term} &
Loss slot for converting the bimodal self-target relation to Module-SIS.  The
current checked value is \(0\). \\
\ec{bimodal\_selftarget\_msis\_hardness\_term} &
The grouped bimodal hardness contribution:
\(\ec{module\_sis\_hardness\_term}+
\ec{bimodal\_to\_msis\_reduction\_loss\_term}\). \\
\ec{signature\_query\_count}, \ec{hash\_query\_count} &
The fixed counted-ROM wrapper constants \(q_S=q_H=2\).  These are not symbolic
adversarial query bounds. \\
\ec{challenge\_support\_bits\_lower\_bound},
\ec{challenge\_support\_cardinality\_lower\_bound} &
The challenge-support lower bound \(58\) bits and \(2^{58}\) elements, encoded
as the concrete real \(288230376151711744\). \\
\ec{counted\_rom\_collision\_term} &
Collision loss \((q_H+q_S+1)^2/N = 25/2^{58}\). \\
\ec{counted\_rom\_prequery\_term} &
Challenge prequery loss \(q_S(q_H+1)/N = 6/2^{58}\). \\
\ec{counted\_rom\_min\_entropy\_term} &
Counted min-entropy loss \(q_S(q_H+1)/N = 6/2^{58}\). \\
\ec{counted\_rom\_programming\_loss\_term} &
The counted-ROM sum \(25/2^{58}+6/2^{58}+6/2^{58}=37/2^{58}\). \\
\ec{fs\_with\_aborts\_*} and \ec{rejection\_sampling\_loss\_term} &
The older paper-style Fiat-Shamir-with-aborts/rejection ledger.  Its
min-entropy component rewrites to \(12\), so it is tracked as a vacuity warning
rather than a useful concrete loss. \\
\ec{haetae\_euf\_bound} &
Grouped legacy bound using the rejection-sampling ledger. \\
\ec{haetae\_euf\_counted\_rom\_bound} &
The final counted-ROM bound used by the top-level theorem. \\
\bottomrule
\end{longtable}

The proof families in this file should be read as follows.

\begin{longtable}{p{0.30\linewidth}p{0.60\linewidth}}
\toprule
\textbf{Lemma family} & \textbf{Proof-script reading} \\
\midrule
\ec{haetae\_euf\_boundE},
\ec{haetae\_euf\_bound\_groupedE},
\ec{haetae\_euf\_bound\_rejection\_groupedE},
\ec{haetae\_euf\_counted\_rom\_bound\_groupedE} &
These lemmas unfold definitions and regroup real sums.  The script pattern is
\ec{rewrite /definition ...; ring}; mathematically, it proves only associativity,
commutativity, and unfolding of named loss groups. \\
Nonnegativity lemmas &
Lemmas such as \ec{rom\_hash\_query\_budget\_nonnegative},
\ec{counted\_rom\_collision\_term\_nonnegative}, and
\ec{rejection\_sampling\_loss\_term\_nonnegative} prove that probabilities and
loss terms can be added or used in transitivity arguments without reversing
inequalities.  Tactics such as \ec{addr\_ge0}, \ec{mulr\_ge0}, and
\ec{divr\_ge0} correspond to elementary real-order closure. \\
Support positivity lemmas &
\ec{challenge\_support\_cardinality\_lower\_bound\_positive} and its
nonnegative corollary justify divisions by \(2^{58}\) and the use of
\ec{ltr\_pdivr\_mulr}. \\
Concrete counted-ROM rewrites &
\ec{counted\_rom\_collision\_termE},
\ec{counted\_rom\_prequery\_termE},
\ec{counted\_rom\_min\_entropy\_termE}, and
\ec{counted\_rom\_programming\_loss\_termE} compute the numerators \(25\), \(6\),
\(6\), and \(37\).  The scripts unfold query counts and use \ec{ring}. \\
Subunit and support-dominance lemmas &
\ec{counted\_rom\_programming\_loss\_term\_subunit},
\ec{counted\_rom\_prequery\_reprogramming\_term\_subunit}, and support
dominance lemmas show the counted-ROM losses are less than \(1\).  These facts
are meaningful concrete sanity checks for the counted ledger. \\
Cardinality-summary lemmas &
\ec{counted\_rom\_loss\_from\_query\_budgets\_and\_support} and related lemmas
package both the formula for the loss and the strict domination of the numerator
by the support size. \\
Vacuity lemmas &
\ec{fs\_with\_aborts\_min\_entropy\_termE},
\ec{fs\_with\_aborts\_min\_entropy\_term\_ge1}, and
\ec{rejection\_sampling\_loss\_term\_ge1} prove that the paper-style
rejection-sampling ledger is already at least \(1\), with a min-entropy
component exactly \(12\).  This is why final claims must preserve the vacuity
warning. \\
\ec{bimodal\_to\_msis\_reduction\_loss\_termE} &
This confirms the current adapter-loss slot is exactly \(0\).  It is a
definition-unfolding fact, not a new hardness theorem. \\
\bottomrule
\end{longtable}

This closes the script-companion coverage for \file{HAETAE\_Reductions.ec} as
an arithmetic ledger: every operation family and lemma family has a mathematical
role and a proof-script reading.  It does not close the cryptographic
assumption coverage for MLWE or Module-SIS; those remain attached to
\file{HAETAE\_Assumptions.ec} and the final theorem composition.

\section{Feedback 3 criticism for HopGames and Security}

Before this expansion, the answer for \file{HAETAE\_HopGames.ec} and
\file{HAETAE\_Security.ec} was still negative: a mathematician or cryptologist
could not reliably understand those EasyCrypt scripts from the report alone,
even with the scripts open.  The evidence is internal to the document.  The
coverage table still marked the declaration dictionary, state interpretation,
restriction explanations, lemma ledger, and proof-script commentary as partial.
The existing \file{HAETAE\_HopGames.ec} section only gave four implementation
layers and a few representative lemma names; it did not explain the long chain
of transcript games, counted-ROM disciplines, sampler bad events, erasure
equivalences, and exact-hyperball bridges.  The existing
\file{HAETAE\_Security.ec} section correctly said that the file assembles the
final theorem family, but it did not tell a reader how the successive
composition lemmas differ, which premises are imported from hop games, which
premises are hardness assumptions, and which steps are just checked real-order
algebra.

The main logical risk was not that the EasyCrypt scripts were unchecked.  The
risk was that the prose made the proof look flatter than it is.  In particular,
the scripts distinguish:
\begin{itemize}
\item exact equivalences, such as erasing transcript logging or replacing one
sampler wrapper by an extensionally equal wrapper;
\item Hoare invariants, such as preservation of transcript soundness, bad-event
flags, and query budgets;
\item probabilistic inequalities, such as splitting a success event into a
clean branch and a bad branch;
\item external or conditional premises, such as MLWE hardness, Module-SIS
hardness, and the structural-to-exact-hyperball sampling obligation;
\item explanatory or coarse fallback facts, especially bounds that are true
because \ec{rejection\_sampling\_loss\_term} is already at least \(1\).
\end{itemize}

Therefore, the revision below does not claim a new cryptographic theorem.  It
adds a companion map that explains what the machine checker proves in these two
files, what remains a premise supplied by another checked file or an external
assumption, and which statements are explanatory bridges or vacuity warnings.

\section{Script companion coverage for HAETAE HopGames}

\file{HAETAE\_HopGames.ec} is the largest script because it turns the generic
EUF-CMA game into transcript-aware and sampler-aware games.  It is not a single
game hop.  It is a stack of games, invariants, exact equivalences, and bad-event
bounds.  The table below is the reading guide for its main proof layers.

\begin{longtable}{p{0.28\linewidth}p{0.62\linewidth}}
\toprule
\textbf{Script layer} & \textbf{Mathematical and proof-script reading} \\
\midrule
\ec{SigningSimulator}, \ec{RealSigningSimulator},
\ec{EUF\_CMA\_SimulatedSign} &
Abstracts the signing oracle behind a simulator interface.  The simulated game
has the same EUF-CMA success predicate as \ec{SIG.EUF\_CMA}, but calls
\ec{Sim.sign} and records the queried message-context pairs.  The checker
verifies procedure calls and state updates; no hardness assumption is used. \\
\ec{transcript\_log\_consistent},
\ec{transcript\_log\_ready},
\ec{internal\_transcript\_state\_sound} &
Defines the invariant connecting three views of signing history:
\ec{queries}, \ec{transcripts}, and \ec{records}.  Read the lemmas ending in
\ec{\_nil}, \ec{\_cons}, \ec{\_valid}, and \ec{\_no\_min\_entropy} as list
invariant lemmas.  They are checked structural facts about logs, not
cryptographic reductions. \\
Budget predicates &
\ec{budgeted\_programmed\_query\_count\_discipline},
\ec{budgeted\_paper\_sim\_signing\_count\_discipline},
\ec{budgeted\_paper\_sim\_sampler\_freshness\_discipline}, and related
predicates are the explicit query-budget conditions.  They are the reason the
counted-ROM terms in \file{HAETAE\_Reductions.ec} can use the concrete
constants \(q_H=q_S=2\). \\
\ec{sampler\_rom\_covered} &
States that every sampler-expand query already present in the lazy-ROM map is
accounted for either as an adversary hash query or as a sampler query.  The
``fresh after clean seed'' lemmas convert absence from the two logs plus a
clean flag into absence from the ROM table. \\
Public-simulation operations &
\ec{public\_sim\_*} operations reconstruct the public pieces of a signing
transcript from public coins and the public key.  Their \ec{\_E} lemmas are
definition-expansion facts used in equivalence proofs. \\
Paper-simulation sample operations &
\ec{paper\_sim\_*}, \ec{paper\_sim\_sample\_from\_public\_coins},
\ec{paper\_sim\_sample\_from\_real\_signing\_coins}, and
\ec{paper\_sim\_sample\_from\_rejection\_attempt} normalize several sampler
representations to a common \ec{paper\_sim\_signature\_sample}.  The checker
proves field equalities and well-formedness; this is explanatory plumbing for
later game hops. \\
Sampler modules &
\ec{ROMPaperSimSigningSampler}, \ec{RealSigningPaperSimSampler},
\ec{ROSigningAttemptPaperSimSampler}, \ec{HAETAERejectionPaperSimSampler},
\ec{ExactHyperballHAETAERejectionPaperSimSampler},
\ec{ExactHyperballPaperSimSampler}, and \ec{ROExactHyperballPaperSimSampler}
are alternate implementations of the signing-sample source.  The proof uses
them as named endpoints in NMA games. \\
NMA wrappers &
\ec{ROMInternalTranscriptAsNMA},
\ec{ROMInternalTranscriptPublicSimAsNMA},
\ec{ROMInternalTranscriptPaperSimAsNMA}, and
\ec{ROMInternalTranscriptBudgetedPaperSimAsNMA} turn an adaptive adversary plus
a simulated signing oracle into an NMA adversary that internally runs the
transcript game.  This is the central structural trick of the proof. \\
Counted and programmed transcript games &
\ec{EUF\_CMA\_ROMInternalTranscriptCountedSign},
\ec{EUF\_CMA\_ROMProgrammedTranscriptCountedSign},
\ec{EUF\_CMA\_ROMProgrammedTranscriptBudgetedCountedSign}, and
\ec{EUF\_CMA\_ROMProgrammedTranscriptBudgetedSampledSign} expose the random
oracle bad flags, programmed sites, and bounded hash/sign counters. \\
Invariant sections &
\ec{InternalTranscriptSignSound},
\ec{CountedROMInternalTranscriptDiscipline},
\ec{ProgrammedROMTranscriptDiscipline},
\ec{BudgetedProgrammedROMTranscriptDiscipline}, and
\ec{BudgetedSampledROMTranscriptDiscipline} prove that the game states stay
sound after oracle calls and adversary calls.  These are Hoare proofs:
\ec{proc}, \ec{wp}, \ec{call}, \ec{rnd}, and \ec{auto} discharge preservation
of explicit predicates. \\
\ec{BudgetedProgrammedToSampledDirectBridge} &
Proves equivalence between a counted game that samples signing coins directly
and a sampled game using \ec{DirectSigningCoinSampler}.  The long relational
postconditions state equality of all shared state fields and bad flags. \\
Erasure and exactness sections &
\ec{HAETAEROMInternalTranscriptExact},
\ec{ROMInternalTranscriptStructuralNMA},
\ec{ROMInternalTranscriptPublicSimNMA},
\ec{ROMInternalTranscriptROMPaperSimNMA},
\ec{ROMInternalTranscriptRealSigningPaperSimNMA},
\ec{ROMInternalTranscriptROSigningAttemptPaperSimNMA},
\ec{TranscriptLoggingErasure}, and \ec{RealSigningSimulatorExact} are
\ec{byequiv} sections.  They prove equality or one-sided implication of success
events by matching calls and preserving corresponding state. \\
Exact-hyperball lifting layer &
\ec{structural\_to\_exact\_hyperball\_paper\_sample\_loss\_obligation} is a
conditional per-call sampling-loss premise.  Lemmas around
\ec{concrete\_ro\_signing\_attempt\_to\_exact\_hyperball\_*} transport this
premise through lazy-ROM freshness and clean-event hypotheses. \\
Concrete \ec{HAETAE\_RO.FRO} layer &
\ec{ConcreteROMInternalTranscriptROSigningAttemptExactHyperballPaperSimNMA}
specializes the abstract oracle lifting to the concrete lazy random oracle.
Several lemmas here still conclude with coarse probability bounds because the
rejection term is already at least \(1\). \\
Rejection and no-fallback exactness &
\ec{ROMInternalTranscriptHAETAERejectionPaperSimNMA} and
\ec{ROMInternalTranscriptExactHyperballPaperSimNMA} connect rejection-sampling
and exact-hyperball samplers.  The ``no fallback'' lemmas rely on support facts
about \ec{dexact\_hyperball\_signing\_attempt\_state}. \\
Special-soundness interface layer &
\ec{BimodalSpecialSoundness\_Game}, \ec{SpecialSoundness\_Game},
\ec{BimodalTranscriptExtractor\_As\_ModuleSIS}, and the
\ec{special\_soundness\_interfaces\_*} lemmas are interface checks for the
forking/special-soundness boundary.  They do not prove a new lattice theorem;
they package obligations imported from the assumption and transcript files. \\
\bottomrule
\end{longtable}

The main mutable state of \file{HAETAE\_HopGames.ec} has the following
mathematical interpretation.

\begin{longtable}{p{0.28\linewidth}p{0.62\linewidth}}
\toprule
\textbf{State component} & \textbf{Meaning} \\
\midrule
\ec{pk\_current}, \ec{sk\_current} &
The public and secret key currently installed in a transcript game or NMA
wrapper.  Invariants typically require \(\ec{pk\_current} =
\ec{public\_key\_of\_secret}(\ec{sk\_current})\). \\
\ec{queries} &
The generic EUF-CMA signing log, containing only message-context pairs. \\
\ec{transcripts} &
The transcript log used by ROM programming and special-soundness reasoning.
Validity and no-min-entropy-failure predicates are stated over this list. \\
\ec{records} &
The synchronized list of message, context, signature, and transcript records.
It is the bridge between the generic signing log and the transcript log. \\
\ec{C.hash\_queries} &
The counted lazy-ROM log of hash queries.  Challenge-query counts are derived
from this list. \\
\ec{C.programmed\_sites} and \ec{C.signature\_sites} &
The random-oracle sites that have been programmed, and the sites induced by
signing transcripts.  Conflict-freedom of these sites is the ROM programming
discipline. \\
\ec{C.bad\_prequery}, \ec{C.bad\_reprogram},
\ec{C.bad\_min\_entropy} &
Bad-event flags for prequerying the challenge site, attempting inconsistent
reprogramming, and hitting a transcript min-entropy failure.  The counted-ROM
proof bounds these flags rather than hiding them in prose. \\
\ec{adversary\_hash\_count}, \ec{signing\_count} &
Budget counters.  The scripts intentionally saturate them at the fixed budgets
instead of leaving the number of queries symbolic. \\
\ec{adversary\_hash\_queries},
\ec{sampler\_expand\_queries},
\ec{sampler\_bad\_prequery} &
The budgeted paper-simulation NMA wrapper's sampler-specific query log and bad
event.  These variables support the split between clean sampling and sampler
prequery failure. \\
Sampler \ec{sk\_current} fields &
Sampler modules store the secret key needed to produce signing samples.  The
equivalence proofs require equality of these fields across paired samplers. \\
\bottomrule
\end{longtable}

The proof-script patterns in this file should be read as follows.

\begin{longtable}{p{0.28\linewidth}p{0.62\linewidth}}
\toprule
\textbf{Pattern} & \textbf{Reader's interpretation} \\
\midrule
\ec{hoare[...]} with invariant preservation &
The checker proves that each oracle or adversary call preserves an explicit
state predicate.  These blocks are not probability bounds; they are deterministic
or almost-sure invariant arguments over program state. \\
\ec{byphoare} converting Hoare facts to probabilities &
Used for zero-probability bad-event lemmas such as transcript bad-forgery and
min-entropy bad-event claims.  The mathematical step is: if a Hoare proof shows
the event is impossible, its probability is \(0\). \\
\ec{byequiv} and long relational postconditions &
The checker compares two games in lockstep.  Equalities such as
\ec{=\{glob H, glob A\}} and field-by-field equalities are the formal statement
that the two executions expose the same interface and return the same result,
or that a left success implies a right success. \\
\ec{Pr[mu\_split]}, \ec{Pr[mu\_sub]}, \ec{Pr[mu\_or]} &
These are probability algebra steps.  They split a success event into clean and
bad branches, drop conjuncts to obtain upper bounds, or union-bound bad events. \\
\ec{ler\_trans}, \ec{ler\_add}, and \ec{smt(mu\_bounded)} &
These are real-order composition steps.  They do not add cryptographic content;
they combine already-proved inequalities and use that probabilities lie between
\(0\) and \(1\). \\
Coarse fallback bounds &
When a proof shows a desired inequality using
\ec{rejection\_sampling\_loss\_term\_ge1} and \ec{mu\_bounded}, the result is
machine-checked but vacuous as a
concrete-security estimate.  The report must keep this distinction visible. \\
\bottomrule
\end{longtable}

For \file{HAETAE\_HopGames.ec}, the machine-checked content is the typing,
Hoare judgments, equivalence judgments, probability inequalities, and real-order
side conditions appearing in the script.  The MLWE and Module-SIS hardness
terms are not proved here.  The structural-to-exact-hyperball paper-sample loss
is a conditional premise when it appears as
\ec{structural\_to\_exact\_hyperball\_paper\_sample\_loss\_obligation}.  The
prose explanations in this section are only a reading guide for those checked
statements.

\subsection{Lemma-ledger companion for HAETAE HopGames}

\paragraph{Lemma-ledger criticism for HAETAE HopGames.}
The current report plus \file{HAETAE\_HopGames.ec} is still not a
source-independent reconstruction manual.  The script has more than eight
thousand lines and many theorem surfaces that look similar but have different
logical status: exact \ec{equiv} results, Hoare invariant preservation,
probability upper bounds, counted-ROM bad-event bounds, rejection-sampling
premises, and coarse fallback inequalities.  Without a ledger, a reader could
mistake a checked erasure equivalence for a sampling-distance theorem, or could
miss that the structural-to-exact-hyperball bridge is conditional when it is
used as an antecedent.  The compact rows below identify the proof obligations
that matter for the final \(\EUFCMA\) path while keeping the EasyCrypt source
as the authoritative tactic-level artifact.

{\footnotesize
\begin{longtable}{p{0.20\linewidth}p{0.25\linewidth}p{0.16\linewidth}p{0.18\linewidth}p{0.15\linewidth}}
\toprule
\textbf{EasyCrypt item and lines} & \textbf{Mathematical statement} &
\textbf{Category} & \textbf{Premises and imports} & \textbf{Downstream role} \\
\midrule
\ec{require import} and \ec{import} block; 1--21 &
Imports the generic signature games, HAETAE scheme model, distributions,
assumption games, transcript predicates, loss terms, rejection facts, ROM
interface, and counted-ROM programming layer. &
Imported/library fact and external-assumption boundary. &
Uses \file{Sig\_ROM.ec}, \file{HAETAE\_Assumptions.ec},
\file{HAETAE\_Transcript.ec}, \file{HAETAE\_Rejection.ec}, and
\file{HAETAE\_ROM\_Programming.ec}. &
Provides all game, hardness, transcript, and bad-event symbols later composed
by \file{HAETAE\_Security.ec}. \\
\midrule
\ec{SigningSimulator}, \ec{RealSigningSimulator},
\ec{EUF\_CMA\_SimulatedSign}; 23--80 &
Introduces an abstract signing-simulator interface, realizes it by the real
scheme signer, and defines a simulator-backed EUF-CMA game with signing-query
logging. &
Machine-checked game/module boundary. &
Requires a \ec{SIG.Scheme}, \ec{SIG.POracle}, and adversary module. &
Starting point for the simulator-erasure and real-signing exactness lemmas used
in the security composition. \\
\midrule
\ec{transcript\_log\_consistent}, \ec{transcript\_log\_ready},
\ec{internal\_transcript\_state\_sound}; 81--441 &
Defines and proves list invariants connecting message-context queries,
transcripts, signing records, public-key consistency, transcript validity, and
min-entropy clear logs. &
Machine-checked fact and invariant family. &
Uses signing-record soundness and ROM-programming no-failure predicates. &
Supplies the state predicates preserved by internal transcript games and bad
forgery bounds. \\
\midrule
\ec{EUF\_CMA\_TranscriptSimulatedSign},
\ec{EUF\_CMA\_InternalTranscriptSign},
\ec{EUF\_CMA\_ROMInternalTranscriptSign}; 442--612 &
Adds transcript logs to simulated signing, replaces the scheme call by internal
signing, and then exposes sampler, message-hash, challenge-hash, and matrix ROM
calls in the signing game. &
Machine-checked game modules. &
Uses transcript constructors, \ec{drandom\_coins}, \ec{sign\_internal}, and
typed ROM query constructors. &
Forms the main source games for the ROM-transcript hop in
\file{HAETAE\_Security.ec}. \\
\midrule
\ec{ROMInternalTranscriptAsNMA}; 613--664 &
Turns the adaptive ROM-internal transcript game into an NMA adversary wrapper
that internally simulates signing transcripts. &
Reduction adapter. &
Uses the adversary through the same oracle/signing interface and reconstructs
the secret key from the public-key seed field. &
First structural bridge from EUF-CMA to UF-NMA in the final composition. \\
\midrule
\ec{public\_sim\_*} operations and \ec{ROMInternalTranscriptPublicSimAsNMA};
665--854 &
Replaces secret-key-dependent signing components by public deterministic
reconstructions and proves the \ec{\_E} equalities against real signing. &
Machine-checked fact and reduction adapter. &
Uses deterministic vector/poly generation, public-key challenge terms, and
signature-field definitions. &
Enables exact replacement of structural NMA by public-simulation NMA. \\
\midrule
\ec{paper\_sim\_signature\_sample},
\ec{paper\_sim\_*}, public/real/rejection sample families; 855--1351 &
Normalizes public coins, real signing coins, and rejection-attempt states into
one paper-simulation sample type; proves well-formedness and field equality
lemmas. &
Machine-checked fact with explanatory sampler-model boundary. &
Uses signature well-formedness, rejection attempt predicates, and current
no-abort facts. &
Provides the sample carrier used by all paper-simulation NMA wrappers. \\
\midrule
\ec{PaperSimSigningSampler} and sampler modules; 1352--1586 &
Defines the sampling interface and seven sampler implementations: ROM public
sample, real signing sample, RO signing attempt, HAETAE rejection attempt,
exact-hyperball rejection, exact-hyperball, and RO exact-hyperball. &
Machine-checked module/procedure boundary. &
Uses \ec{drandom\_coins}, ROM sampler queries, rejection distributions, and
\ec{dexact\_hyperball\_signing\_attempt\_state}. &
Names the endpoints of the paper-simulation and exact-hyperball hop chain. \\
\midrule
\ec{StructuralAttemptPaperSample}, \ec{ExactHyperballPaperSample},
\ec{ROMInternalTranscriptPaperSimAsNMA},
\ec{ROMInternalTranscriptBudgetedPaperSimAsNMA}; 1587--1765 &
Defines direct one-procedure sample sources and NMA wrappers, including a
budgeted wrapper with adversary-hash count, signing count, sampler-query log,
and sampler-prequery bad flag. &
Machine-checked game modules and reduction adapters. &
Uses \ec{PaperSimSigningSampler}, hash-query budgets, signature-query budgets,
and sampler prequery predicates. &
Provides the budgeted NMA games whose probabilities are charged to counted-ROM
and rejection-sampling loss terms. \\
\midrule
\ec{ROMInternalTranscriptBudgetedPaperSimNMAQueryBudget}; 1766--2622 &
Proves signing-count preservation, sampler-expand point/log/budget bounds,
lazy-ROM sampler freshness, and clean sampler one-call facts. &
Machine-checked Hoare/probability invariant section. &
Premises include query-budget disciplines, sampler-ROM coverage, and
losslessness/freshness of concrete lazy ROM calls. &
Supplies the one-call and NMA-level sampler-prequery estimates used before the
exact-hyperball bridge. \\
\midrule
\ec{EUF\_CMA\_ROMInternalTranscriptCountedSign},
\ec{EUF\_CMA\_ROMProgrammedTranscriptCountedSign},
\ec{EUF\_CMA\_ROMProgrammedTranscriptBudgetedCountedSign},
\ec{EUF\_CMA\_ROMProgrammedTranscriptBudgetedSampledSign}; 2624--2978 &
Introduces counted/programmed ROM games, then adds explicit hash/sign budgets,
programmed challenge sites, and a sampled signing source. &
Machine-checked game modules. &
Uses \ec{CountedLazyROM}, \ec{ProgramChallenge}, challenge-query counters, and
direct signing-coin sampler interfaces. &
These games expose the bad flags later bounded by counted-ROM programming
loss. \\
\midrule
\ec{InternalTranscriptSignSound}; 2979--3574 &
Proves internal and ROM-internal transcript-game state soundness, transcript
validity, log readiness, no-min-entropy, signature-site-clearness, and bad
forgery zero/count bounds. &
Machine-checked Hoare and probability fact family. &
Uses internal transcript invariants and \ec{byphoare} zero-event reasoning. &
Feeds the ROM-transcript clear-site reductions in the security file. \\
\midrule
\ec{CountedROMInternalTranscriptDiscipline}; 3576--3769 &
Maintains counted-ROM budget state and derives bad-clear, bad-false, counted
bad zero, budget-sound, and subunit bad-event lemmas. &
Machine-checked invariant and probability boundary. &
Uses \ec{CountedLazyROM} state facts and counted-ROM programming bound
obligation. &
Connects counted bad flags to \ec{counted\_rom\_programming\_loss\_term}. \\
\midrule
\ec{ProgrammedROMTranscriptDiscipline}; 3771--3904 &
Shows programmed counted games preserve min-entropy-clear state and bounds
min-entropy and prequery/reprogramming bad events from a budget expression. &
Machine-checked invariant and probability bound. &
Uses transcript no-failure lemmas and counted-ROM bad-flag wrappers. &
Supplies the programmed-ROM discipline consumed by budgeted variants. \\
\midrule
\ec{BudgetedProgrammedROMTranscriptDiscipline}; 3906--4518 &
Proves challenge-query discipline, reprogram discipline, query-budget
preservation, adaptive prequery lifting, and checked concrete bad-event bounds
for the budgeted programmed game. &
Machine-checked invariant and query-budget family. &
Premises include \ec{hash\_query\_budget\_count}, challenge-query counts, and
honest-site conflict-freedom. &
Produces the budget-certified branch bounds used in counted-ROM security
composition. \\
\midrule
\ec{BudgetedSampledROMTranscriptDiscipline}; 4520--5228 &
Repeats the budget and bad-event discipline for the sampled-direct game and
proves sampler-message-hash prequery and prequery/reprogramming bad bounds. &
Machine-checked invariant and probability bound. &
Uses \ec{DirectSigningCoinSampler} and the ROM-programming prequery lemmas. &
Prepares the sampled-direct endpoint for equivalence with the programmed game. \\
\midrule
\ec{BudgetedProgrammedToSampledDirectBridge}; 5230--5599 &
Gives equivalences for hash, sign, adversary, and main procedures between
budgeted programmed counted and sampled-direct games; transfers bad-event
probabilities. &
Checked \ec{equiv} family and reduction adapter. &
Requires module restrictions preventing access to \ec{CountedLazyROM} internals
and shared game state. &
Lets the final proof use direct-sampler prequery estimates for the programmed
game. \\
\midrule
\ec{HAETAEROMInternalTranscriptExact}; 5601--5704 &
Proves oracle/adversary erasure equivalences and
\ec{haetae\_rom\_internal\_transcript\_erasure\_exact}. &
Checked exact equivalence. &
Uses inlining of \ec{HAETAE} and ROM-internal transcript games under module
separation restrictions. &
First exact hop from generic EUF-CMA to the ROM-internal transcript game in
\file{HAETAE\_Security.ec}. \\
\midrule
\ec{ROMInternalTranscriptStructuralNMA}; 5706--5868 &
Relates the ROM-internal transcript game to the structural NMA wrapper and
proves clear-site and ordinary structural NMA bounds. &
Checked equivalence/probability inequality. &
Uses oracle/adversary equivalences and transcript-log clear-site events. &
Main structural EUF-CMA-to-NMA reduction used by counted and legacy security
lemmas. \\
\midrule
\ec{ROMInternalTranscriptPublicSimNMA},
\ec{ROMInternalTranscriptROMPaperSimNMA},
\ec{ROMInternalTranscriptRealSigningPaperSimNMA},
\ec{ROMInternalTranscriptROSigningAttemptPaperSimNMA}; 5874--6624 &
Proves exact replacement steps among public simulation, ROM paper simulation,
real-signing paper simulation, and RO signing-attempt paper simulation. &
Checked exact equivalence family. &
Uses public-sim equalities, paper-sample equalities, and sampler module
restrictions. &
Builds the NMA-side sampler chain later consumed by Security. \\
\midrule
\ec{structural\_to\_exact\_hyperball\_paper\_sample\_loss\_obligation}
and \ec{concrete\_ro\_*loss} family; 6615--6986 &
Names the structural-to-exact-hyperball sampling-loss obligation and transports
that per-call premise through seed freshness, clean logs, and concrete lazy-ROM
losslessness. &
Conditional premise plus machine-checked probability inequalities. &
Premises include sampler freshness, old-log cleanliness, and rejection-sampling
loss terms. &
This is the key statistical-loss bridge; it must not be reported as a proved
paper sampler theorem unless the antecedent is discharged. \\
\midrule
\ec{ROMInternalTranscriptROSigningAttemptExactHyperballPaperSimNMA};
7002--7291 &
Lifts the structural-to-exact-hyperball loss through NMA games, budgeted
wrappers, clean conditioning, checked clean lifting, and coarse bounded loss. &
Machine-checked probability inequality family with conditional branches. &
Uses the sampling-loss obligation, budgeted NMA wrappers, clean-event lemmas,
and sometimes coarse \ec{mu\_bounded} reasoning. &
Provides the RO-signing-attempt to RO-exact-hyperball NMA bounds used by final
Security lemmas. \\
\midrule
\ec{ConcreteROMInternalTranscriptROSigningAttemptExactHyperballPaperSimNMA};
7318--7817 &
Specializes the abstract RO-signing-attempt/exact-hyperball bridge to
\ec{HAETAE\_RO.FRO} and proves signature-event preservation, losslessness,
projection, and concrete clean-lifting lemmas. &
Machine-checked concrete ROM lifting. &
Uses concrete lazy-ROM facts, \ec{dexact\_hyperball} support facts, and
sampler-event preservation. &
Supplies concrete evidence for the checked RO sampling hop used by the final
counted-ROM theorem. \\
\midrule
\ec{ROMInternalTranscriptHAETAERejectionPaperSimNMA},
\ec{ROMInternalTranscriptExactHyperballPaperSimNMA}; 7819--8230 &
Proves exact swaps from real signing to HAETAE rejection sampling and from
exact-hyperball rejection sampling to exact-hyperball paper simulation with no
fallback. &
Checked exact equivalence and no-fallback theorem family. &
Uses rejection acceptance/no-abort lemmas from \file{HAETAE\_Rejection.ec}. &
Removes rejection-specific wrappers before the final exact-hyperball NMA
assumption is applied. \\
\midrule
\ec{ROMInternalTranscriptPaperSimHop}; 8232--8254 &
Turns a public-sim-to-paper-sim sampling hop into an NMA probability bound with
\ec{rejection\_sampling\_loss\_term}. &
Machine-checked probability inequality. &
Premise: a concrete sampling-hop inequality for the two NMA wrappers. &
Used by Security when composing paper-simulation NMA bounds. \\
\midrule
\ec{TranscriptLoggingErasure}, \ec{RealSigningSimulatorExact}; 8256--8422 &
Shows transcript logging is observationally erased and the real signing
simulator matches the scheme signing oracle. &
Checked exact equivalence. &
Uses module restrictions excluding game internals and simulator state. &
Connects the simulator-based games in Security back to the generic
\ec{SIG.EUF\_CMA} game. \\
\midrule
\ec{ForkingAdversary}, \ec{TranscriptExtractor},
\ec{BimodalTranscriptExtractor}, \ec{BimodalSpecialSoundness\_Game},
\ec{SpecialSoundness\_Game},
\ec{BimodalTranscriptExtractor\_As\_ModuleSIS}; 8424--8499 &
Defines the forking adversary and extractor interfaces and the games that test
bimodal self-target MSIS and Module-SIS extraction from transcript forks. &
Imported-interface/game boundary. &
Uses transcript extraction operators and Module-SIS adapters. &
Provides the special-soundness interface normalized before final hardness
composition. \\
\midrule
\ec{BimodalSpecialSoundnessLift},
\ec{special\_soundness\_interfaces\_*},
\ec{fs\_with\_aborts\_interfaces\_*}; 8500--8649 &
Proves the bimodal-to-Module-SIS special-soundness lift and packages
forking, public reconstruction, signer soundness, rejection-bound, and FS
interfaces as ready. &
Machine-checked interface theorem family. &
Uses \ec{forking\_extraction\_obligation\_holds},
\ec{bimodal\_to\_module\_sis\_lift\_obligation\_holds}, and reduction-term
nonnegativity. &
Supplies the interface premise consumed by several Security theorem variants. \\
\midrule
\ec{structural\_to\_exact\_hyperball\_*} and
\ec{dexact\_hyperball\_*boundary\_*}; 8651--8800 &
Relates structural-to-exact-hyperball obligations to attempt/coin-sample-pair
loss premises and proves exact-hyperball rejection/no-fallback distribution
boundary lemmas. &
Machine-checked implications from conditional premises and model-boundary facts. &
Uses \file{HAETAE\_Rejection.ec} exact-hyperball support and loss-obligation
lemmas. &
Closes local boundary checks but does not prove paper/implementation sampler
equivalence. \\
\midrule
Module restriction families; 1768--1771, 2981--2984, 3578--3585,
3773--3775, 3908--3911, 4522--4525, 5232--5236, 5603--5606,
5708--5711, 5876--5879, 6043--6046, 6214--6217, 6420--6423,
7004--7008, 7821--7824, 8032--8035, 8234--8241, 8258--8264,
8342--8346, 8502--8505 &
Declares which adversary, oracle, sampler, and game modules cannot access one
another's private state. &
Imported/module-separation fact checked by EasyCrypt. &
Depends on EasyCrypt's module system and the declared game interfaces. &
Prevents circular calls and hidden-state leaks in every equivalence and Hoare
proof used by the final security composition. \\
\bottomrule
\end{longtable}
}

\section{Script companion coverage for HAETAE Security}

\file{HAETAE\_Security.ec} is the composition file.  It does not implement the
sampler or ROM machinery itself.  Instead, it imports the previous scripts and
assembles theorem families by repeatedly applying already-proved hops, hardness
premises, and arithmetic rewrites.

\begin{longtable}{p{0.30\linewidth}p{0.60\linewidth}}
\toprule
\textbf{Section or lemma family} & \textbf{Mathematical and proof-script reading} \\
\midrule
\ec{correctness\_obligation} and
\ec{correctness\_obligation\_holds} &
Defines perfect correctness as verification of every internally generated
signature.  The proof unfolds \ec{keygen\_internal} and applies
\ec{verify\_internal\_sign\_internal}.  This is checked internal correctness,
not a security reduction. \\
\ec{correctness\_perfect} &
Shows the generic correctness game fails with probability \(0\).  The proof
inlines the game and HAETAE procedures, uses Hoare reasoning, samples keys and
coins, and discharges the final verification facts with
\ec{valid\_signature\_sign\_internal} and \ec{verify\_norm\_ok\_current}. \\
\ec{NoSigningSecurity} &
Uses \ec{SIG.nma\_as\_cma\_exact} to embed an NMA adversary as a CMA adversary
with no signing queries.  The lemma composes an NMA-to-hardness premise, the
bimodal-to-Module-SIS premise, and MLWE/Module-SIS hardness bounds into the
legacy \ec{haetae\_euf\_bound}. \\
Generic \ec{euf\_cma\_security} lemmas &
These are algebraic composition lemmas.  Their inputs are a Fiat-Shamir
hop bound, an NMA hardness bound, a bimodal-to-Module-SIS bound, and the two
hardness-term bounds.  The proof uses \ec{ler\_trans}, \ec{ler\_add}, and
bound rewrites. \\
\ec{FullReductionWithSimulatorHop} &
Replaces the generic EUF-CMA game by \ec{EUF\_CMA\_SimulatedSign} using
\ec{real\_signing\_simulator\_exact}.  The section is an exact-game
interface bridge followed by the same real-order composition as the generic
lemma. \\
\ec{FullReductionWithROMTranscriptHop} &
Moves from ordinary signing to ROM-internal transcript signing.  The
\ec{rom\_internal\_transcript\_hop\_from\_clear\_site\_*} lemmas split success
into a clear-site branch plus a bad-forgery branch and charge the bad branch to
either the legacy FS-with-aborts terms or the counted-ROM programming term. \\
\ec{MechanizedBimodalModuleSISLift} &
\ec{bimodal\_to\_module\_sis\_reduction\_bound} rewrites the checked exact
adapter \ec{bimodal\_msis\_as\_module\_sis\_exact} plus the zero adapter-loss
term.  It is not an independent proof of Module-SIS hardness. \\
\ec{MechanizedROMTranscriptBimodalLift} &
Combines the ROM transcript hop with the mechanized bimodal-to-Module-SIS
adapter.  The counted version targets \ec{haetae\_euf\_counted\_rom\_bound};
the legacy version targets \ec{haetae\_euf\_bound}. \\
\ec{MechanizedROMTranscriptConstructedNMALift} &
This is the final composition lane.  It progressively replaces structural NMA,
public-simulation NMA, paper-simulation NMA, real-signing paper simulation,
HAETAE-rejection simulation, exact-hyperball simulation, and random-oracle
exact-hyperball simulation, then applies MLWE and bimodal/Module-SIS bounds. \\
\ec{euf\_cma\_security\_from\_checked\_ro\_sampling\_hop\_and\_counted\_bimodal}
&
This is the representative checked counted-ROM theorem shape.  Its premise is
an NMA bound for \ec{ROExactHyperballPaperSimSampler} plus
\ec{rejection\_sampling\_loss\_term}.  The conclusion is the formal
EUF-CMA bound \ec{haetae\_euf\_counted\_rom\_bound}. \\
\ec{euf\_cma\_security\_from\_budgeted\_paper\_sample\_lifting\_and\_counted\_bimodal}
&
This variant exposes the budgeted lifting premises explicitly:
unbudgeted-to-budgeted for the signing-attempt sampler,
budgeted-to-unbudgeted for the exact-hyperball sampler, the structural
sampling-loss obligation, and a final NMA bound containing
\(\ec{rom\_signature\_query\_budget}\cdot
\ec{rejection\_sampling\_loss\_term}\). \\
\bottomrule
\end{longtable}

The final composition section uses module restrictions to enforce information
flow.  The exclusions should be read by category, not as arbitrary syntax.

\begin{longtable}{p{0.30\linewidth}p{0.60\linewidth}}
\toprule
\textbf{Restricted modules} & \textbf{Why the restriction matters} \\
\midrule
\ec{SIG.EUF\_CMA} and \ec{EUF\_CMA\_*} games &
Prevent the adversary or oracle from inspecting game-private secret keys,
query logs, transcript logs, bad flags, or return events. \\
\ec{ROMInternalTranscript*AsNMA} wrappers &
Prevent circular access to the NMA wrappers that are being used as simulated
adversaries in the reduction. \\
Paper-simulation sampler modules &
Restrictions on \ec{ROMPaperSimSigningSampler},
\ec{RealSigningPaperSimSampler}, \ec{ROSigningAttemptPaperSimSampler},
\ec{HAETAERejectionPaperSimSampler},
\ec{ExactHyperballHAETAERejectionPaperSimSampler},
\ec{ExactHyperballPaperSimSampler}, and
\ec{ROExactHyperballPaperSimSampler} prevent adversaries from using sampler
internals except through the declared signing-oracle interface. \\
\ec{HAETAE} &
When excluded from the adversary, prevents direct calls to the scheme module
outside the intended oracle/game interface. \\
\ec{H} and \ec{A} exclusions in \ec{Samp} &
Ensure a sampler cannot call back into the random oracle or adversary except
through the interfaces explicitly passed to it by the wrapper. \\
\ec{Bmlwe}, \ec{Bbimodal}, and \ec{Bsis} &
These are abstract reduction modules.  Their success probabilities are bounded
by premises or assumption terms; the security file composes those bounds but
does not implement the reductions' internals. \\
\bottomrule
\end{longtable}

The proof-script patterns in \file{HAETAE\_Security.ec} are deliberately
simple.  A \ec{rewrite} of an exactness lemma changes the game whose probability
is being bounded without loss.  A \ec{rewrite} of a grouped bound such as
\ec{haetae\_euf\_counted\_rom\_bound\_groupedE} unfolds the arithmetic ledger.
An \ec{apply} of a previous security lemma selects the next composition theorem.
Nested \ec{ler\_trans} and \ec{ler\_add} blocks compose inequalities and
preserve nonnegative loss terms.  The checker verifies all of these real-order
steps, but the cryptographic assumptions remain exactly the assumptions named
in the premises: MLWE hardness, Module-SIS hardness, NMA-to-hardness bounds,
and sampling-hop obligations.

Consequently, this file should be read as a machine-checked composition
argument, not as a source of new assumptions.  The statements that are
machine-checked are the conditional theorem statements and their real-order
proofs.  The hardness estimates are assumed through the real-valued terms and
module-game premises.  The prose in this report explains that structure; it is
not itself part of the EasyCrypt trusted base.

\subsection{Lemma-ledger companion for HAETAE Security}

\paragraph{Lemma-ledger criticism for HAETAE Security.}
The existing companion coverage correctly says that
\file{HAETAE\_Security.ec} is a composition file, but it still underexplains
the theorem boundaries.  The script has many lemmas whose conclusions look like
the same EUF-CMA bound, while the premises expose different hop surfaces:
legacy FS-with-aborts terms, counted-ROM programming loss, exact erasure,
mechanized bimodal-to-Module-SIS lifting, exact-hyperball sampling, checked RO
sampling, or budgeted paper-sample lifting.  A reader cannot reconstruct the
logical path unless each lemma family is classified as checked arithmetic
composition, imported hop theorem, external hardness premise, or conditional
sampling premise.  The ledger below is therefore a theorem-boundary map, not a
claim that the report alone replays the EasyCrypt proofs.

{\footnotesize
\begin{longtable}{p{0.20\linewidth}p{0.25\linewidth}p{0.16\linewidth}p{0.18\linewidth}p{0.15\linewidth}}
\toprule
\textbf{EasyCrypt item and lines} & \textbf{Mathematical statement} &
\textbf{Category} & \textbf{Premises and imports} & \textbf{Downstream role} \\
\midrule
\ec{require import} and \ec{import} block; 1--19 &
Imports generic signature games, HAETAE scheme/algebra/distributions,
reductions, assumptions, rejection facts, transcript predicates,
ROM-programming facts, HopGames, and real-order lemmas. &
Imported/library fact and external-assumption boundary. &
Uses \file{HAETAE\_HopGames.ec} as the main checked hop provider and
\file{HAETAE\_Assumptions.ec} for hardness games. &
Defines the entire theorem-composition environment. \\
\midrule
\ec{correctness\_obligation},
\ec{correctness\_obligation\_holds}; 21--34 &
Every internally generated signature verifies for the formal HAETAE model. &
Machine-checked fact. &
Uses \ec{keygen\_internal}, \ec{sign\_internal}, and
\ec{verify\_internal\_sign\_internal}. &
Local correctness fact used before security composition; not a hardness
assumption. \\
\midrule
\ec{correctness\_perfect}; 35--70 &
The generic correctness game fails with probability zero for the formal scheme. &
Machine-checked Hoare/probability fact. &
Requires module restrictions \ec{A <: SIG.CORR\_ADV \{-H, -HAETAE\}} and
algebra lemmas for valid signatures and current verification norms. &
Certifies perfect correctness of the EasyCrypt model. \\
\midrule
\ec{NoSigningSecurity}, \ec{euf\_cma\_security\_no\_signing}; 73--124 &
Composes an NMA bound, a bimodal-to-Module-SIS bound, MLWE hardness, and
Module-SIS hardness to bound an NMA-as-CMA adversary. &
Machine-checked composition with external cryptographic assumptions. &
Premises are NMA, MLWE, Module-SIS, and bimodal reduction probabilities; uses
\ec{SIG.nma\_as\_cma\_exact}. &
Baseline no-signing theorem; not the main adaptive-signing result. \\
\midrule
\ec{euf\_cma\_security},
\ec{euf\_cma\_security\_with\_hop\_interfaces}; 127--208 &
Generic EUF-CMA composition from an FS-with-aborts hop, NMA bound,
bimodal-to-Module-SIS bound, and MLWE/Module-SIS hardness assumptions. &
Machine-checked real-order composition; conditional premises. &
Requires \ec{fs\_with\_aborts\_interfaces\_ready} only as an explicit interface
premise in the second lemma. &
Template used by later sections after replacing the first hop with more
specific HopGames lemmas. \\
\midrule
\ec{FullReductionWithSimulatorHop}; 211--304 &
Uses \ec{real\_signing\_simulator\_exact} to replace generic EUF-CMA by the
simulated-sign game, then applies the generic composition theorem. &
Checked exact-hop composition. &
Imports the simulator exactness lemma from HopGames and external hardness
premises. &
Connects the generic theorem shape to the simulator game surface. \\
\midrule
\ec{FullReductionWithROMTranscriptHop}; 306--405 &
Uses \ec{haetae\_rom\_internal\_transcript\_erasure\_exact}; splits
ROM-internal transcript success into clear-site and bad-site branches, charging
the bad branch to legacy FS terms or counted-ROM programming loss. &
Machine-checked probability split and checked hop composition. &
Uses \ec{rom\_internal\_transcript\_bad\_forgery\_bound} and
\ec{rom\_internal\_transcript\_bad\_forgery\_counted\_bound}. &
Introduces the counted-ROM path used by the final theorem family. \\
\midrule
\ec{MechanizedBimodalModuleSISLift}; 407--459 &
Proves the bimodal self-target MSIS game is bounded by the Module-SIS adapter
game plus the zero adapter-loss term, then composes that into EUF-CMA bounds. &
Machine-checked adapter equality and real-order composition. &
Uses \ec{bimodal\_msis\_as\_module\_sis\_exact} and
\ec{bimodal\_to\_msis\_reduction\_loss\_termE}. &
Removes the need for an arbitrary \ec{Bsis} module in later counted-ROM
statements. \\
\midrule
\ec{MechanizedROMTranscriptBimodalLift}; 461--596 &
Combines the ROM-transcript hop with the mechanized bimodal-to-Module-SIS
adapter, both for legacy and counted-ROM bounds, and adds clear-site variants. &
Machine-checked theorem-boundary composition. &
Premises are ROM hop or clear-site reductions, an NMA bound, MLWE hardness, and
Module-SIS hardness for the adapter game. &
Main bridge from HopGames' ROM transcript lemmas to grouped final bounds. \\
\midrule
\ec{MechanizedROMTranscriptConstructedNMALift} declarations; 598--631 &
Declares the final composition modules and excludes all relevant game,
wrapper, sampler, and adversary internals from each other. &
Module restriction family. &
Depends on EasyCrypt's module separation checker and all prior game/module
interfaces. &
Protects every exact equivalence and reduction application in the final
composition lane. \\
\midrule
\ec{euf\_cma\_security\_from\_structural\_nma\_and\_counted\_bimodal}
through
\ec{euf\_cma\_security\_from\_real\_signing\_paper\_sim\_nma\_and\_counted\_bimodal};
633--744 &
Progressively rewrites the NMA bound through structural, public-sim,
paper-sim, ROM-paper-sim, and real-signing paper-sim wrappers. &
Checked exact-hop composition with external hardness premises. &
Uses HopGames exactness lemmas such as
\ec{rom\_internal\_nma\_public\_sim\_exact} and
\ec{rom\_internal\_nma\_real\_signing\_paper\_sim\_exact}. &
Moves the final theorem toward the sampler implementations used by the checked
RO sampling hop. \\
\midrule
\ec{euf\_cma\_security\_from\_haetae\_rejection\_*},
\ec{exact\_hyperball\_haetae\_rejection\_nma\_bound\_*},
\ec{haetae\_rejection\_nma\_bound\_*}; 746--842 &
Transfers NMA bounds between HAETAE rejection sampling and exact-hyperball
samplers, with or without an added rejection-sampling loss term. &
Checked theorem-boundary composition; conditional sampling premise when a hop
bound is an antecedent. &
Uses \ec{rom\_internal\_nma\_real\_signing\_haetae\_rejection\_paper\_sim\_exact}
and exact-hyperball no-fallback lemmas. &
Builds the rejection-to-exact-hyperball bridge used by final counted-ROM
theorems. \\
\midrule
\ec{real\_signing\_ro\_exact\_hyperball\_*},
\ec{haetae\_rejection\_ro\_exact\_hyperball\_*}; 843--972 &
Transfers real-signing and HAETAE-rejection NMA bounds through RO signing
attempt and RO exact-hyperball samplers, including budgeted lifting with a
signature-budget multiplier. &
Machine-checked probability composition with conditional sampling premises. &
Uses \ec{rom\_internal\_nma\_ro\_signing\_attempt\_ro\_exact\_hyperball\_*}
from HopGames and the
\ec{structural\_to\_exact\_hyperball\_paper\_sample\_loss\_obligation}
antecedent where required. &
Provides the sampling-hop premises used by the final checked and budgeted
theorem variants. \\
\midrule
\ec{euf\_cma\_security\_from\_exact\_hyperball\_*},
\ec{euf\_cma\_security\_from\_ro\_exact\_hyperball\_*},
\ec{euf\_cma\_security\_from\_ro\_signing\_attempt\_*}; 974--1099 &
Composes exact-hyperball, RO exact-hyperball, and RO signing-attempt NMA hops
with MLWE and Module-SIS hardness to conclude the counted-ROM EUF-CMA bound. &
Machine-checked final-bound composition. &
Premises include NMA hop inequalities plus MLWE and Module-SIS hardness
probability bounds. &
These are alternate final theorem surfaces for different supplied sampling-hop
evidence. \\
\midrule
\ec{real\_signing\_ro\_exact\_hyperball\_hop\_checked\_from\_loss\_budget},
\ec{real\_signing\_ro\_exact\_hyperball\_hop\_from\_paper\_sample\_lifting};
1102--1138 &
Derives the real-signing-to-RO-exact-hyperball hop either from the checked
loss-budget lemma or from an explicit paper-sample lifting premise. &
Machine-checked implication; conditional premise for the lifting variant. &
Uses HopGames' \ec{rom\_internal\_nma\_ro\_signing\_attempt\_ro\_exact\_hyperball\_*}
lemmas. &
Supplies the first premise for the final checked RO sampling theorem. \\
\midrule
\ec{euf\_cma\_security\_from\_paper\_sample\_lifting\_and\_counted\_bimodal},
\ec{euf\_cma\_security\_from\_budgeted\_paper\_sample\_lifting\_and\_counted\_bimodal};
1140--1214 &
Final counted-ROM bounds that expose either the paper-sample lifting premise or
the budgeted unbudgeted/budgeted conversion premises explicitly. &
Machine-checked theorem-boundary composition with conditional sampling
premises. &
Requires MLWE hardness, Module-SIS hardness, an RO exact-hyperball NMA bound,
and sampling-lifting premises. &
Documents which assumptions must be supplied to avoid relying solely on coarse
checked fallback bounds. \\
\midrule
\ec{euf\_cma\_security\_from\_checked\_ro\_sampling\_hop\_and\_counted\_bimodal};
1217--1237 &
Representative checked counted-ROM theorem: if the RO exact-hyperball NMA
probability plus rejection-sampling loss is bounded by MLWE plus bimodal
success, and MLWE/Module-SIS are bounded by their hardness terms, then the
formal HAETAE EUF-CMA game is bounded by
\ec{haetae\_euf\_counted\_rom\_bound}. &
Machine-checked final theorem boundary with external hardness premises. &
Uses \ec{real\_signing\_ro\_exact\_hyperball\_hop\_checked\_from\_loss\_budget}
and counted-ROM grouped arithmetic. &
Best current final theorem shape, but still conditional on the formal ROM
model and rejection/sampling loss ledger. \\
\midrule
\ec{euf\_cma\_security\_from\_real\_signing\_exact\_hyperball\_hop\_and\_counted\_bimodal};
1239--1269 &
Alternative final theorem starting from a real-signing-to-exact-hyperball hop
plus an exact-hyperball NMA bound. &
Machine-checked final theorem boundary with conditional premises. &
Requires MLWE/Module-SIS hardness and the supplied exact-hyperball NMA
inequality. &
Records a non-RO exact-hyperball route through the same counted-ROM final
bound. \\
\bottomrule
\end{longtable}
}

\subsection{Proof-obligation replay for the final Security theorem path}

\paragraph{Replay-level criticism.}
The lemma ledger above is enough to locate the final theorem surfaces, but it
is not yet enough for a mathematician or cryptologist to reconstruct the exact
EasyCrypt proof obligations used by the two terminal theorem shapes.  In
particular, the report previously did not spell out which antecedent is passed
to which \ec{apply}, where the random-oracle/rejection loss enters, which
HopGames lemma is imported, or whether a sampling hop is proved non-vacuously,
proved by a coarse boundedness fallback, or simply assumed as a theorem
premise.  The replay below closes that gap at the proof-obligation level.  It
is still not a tactic transcript: the EasyCrypt source remains necessary for
checking module restrictions, exact syntax, local names, and SMT side
conditions.

\paragraph{Common final-composition frame.}
Both terminal theorem paths live in
\ec{MechanizedROMTranscriptConstructedNMALift}.  The modules declared at
\file{HAETAE\_Security.ec} lines 598--631 are part of the theorem boundary:
\ec{H} is a signature oracle, \ec{A} is the EUF-CMA adversary, and
\ec{Bmlwe}/\ec{Bbimodal} are reduction modules for the MLWE and bimodal
self-target MSIS games.  EasyCrypt checks the module-separation restrictions,
but the numerical hardness estimates for
\ec{Pr[MLWE\_Game(Bmlwe).main]} and
\ec{Pr[ModuleSIS\_Game(BimodalMSIS\_As\_ModuleSIS(Bbimodal)).main]} enter as
premises.  The final arithmetic target is the imported bound expression
\ec{haetae\_euf\_counted\_rom\_bound}; expanding that expression is a
reduction-file exercise, not a new theorem in \file{HAETAE\_Security.ec}.

The common counted-ROM composition engine is
\ec{euf\_cma\_security\_from\_haetae\_rejection\_paper\_sim\_nma\_and\_counted\_bimodal}
at lines 746--766.  It rewrites the real-signing paper-simulation NMA game to
the HAETAE rejection paper-simulation NMA game using
\ec{rom\_internal\_nma\_real\_signing\_haetae\_rejection\_paper\_sim\_exact},
then passes the NMA, MLWE, and Module-SIS premises into the earlier
real-signing-paper-simulation counted-ROM theorem.  Thus the terminal
theorems do not redo the EUF-CMA game-hop proof; they supply a suitable NMA
premise for that engine and forward the two hardness premises.

\paragraph{Replay for the checked RO-sampling final theorem.}
The theorem
\ec{euf\_cma\_security\_from\_checked\_ro\_sampling\_hop\_and\_counted\_bimodal}
appears at \file{HAETAE\_Security.ec} lines 1217--1237.  Its statement says:
if the RO exact-hyperball paper-simulation NMA success probability plus
\ec{rejection\_sampling\_loss\_term} is bounded by MLWE success plus bimodal
self-target MSIS success, and if those MLWE/Module-SIS games are bounded by
their hardness terms, then the formal HAETAE EUF-CMA game is bounded by
\ec{haetae\_euf\_counted\_rom\_bound}.

{\scriptsize
\begin{longtable}{p{0.08\linewidth}p{0.25\linewidth}p{0.34\linewidth}p{0.25\linewidth}}
\toprule
\textbf{Step} & \textbf{EasyCrypt evidence} & \textbf{Proof obligation replay} &
\textbf{Status boundary} \\
\midrule
1 &
Lines 1230--1232:
\ec{move=> ro\_exact\_nma\_hop mlwe\_bound msis\_bound} and
\ec{apply euf\_cma\_security\_from\_ro\_exact\_hyperball\_sampling\_hop\_and\_counted\_bimodal}. &
The final theorem is reduced to the RO exact-hyperball counted-ROM theorem at
lines 1036--1070.  That theorem requires four inputs: a real-signing to RO
exact-hyperball NMA hop, the RO exact-hyperball NMA-to-MLWE/bimodal premise,
the MLWE hardness premise, and the Module-SIS hardness premise. &
Machine-checked proof replay step; not an external assumption. \\
\midrule
2 &
Lines 1233 and 1102--1114:
\ec{real\_signing\_ro\_exact\_hyperball\_hop\_checked\_from\_loss\_budget}. &
The first required input is discharged by proving
\[
\Prb[\ec{RealSigningPaperSimSampler}]
 \le
\Prb[\ec{ROExactHyperballPaperSimSampler}]
 + \ec{rejection\_sampling\_loss\_term}.
\]
The local proof applies
\ec{real\_signing\_ro\_exact\_hyperball\_hop\_from\_ro\_signing\_attempt\_hop}
and supplies the HopGames lemma
\ec{rom\_internal\_nma\_ro\_signing\_attempt\_ro\_exact\_hyperball\_loss\_bound}. &
Machine-checked implication, but its imported HopGames input is a coarse
fallback bound. \\
\midrule
3 &
\file{HAETAE\_Security.ec} lines 865--884 and
\file{HAETAE\_HopGames.ec} lines 7291--7314. &
The real-signing sampler is first rewritten to the RO-signing-attempt sampler
by \ec{rom\_internal\_nma\_real\_signing\_ro\_attempt\_paper\_sim\_exact}.
The remaining RO-signing-attempt to RO exact-hyperball inequality is bounded in
HopGames by \ec{mu\_bounded}, nonnegativity of the right probability, and
\ec{rejection\_sampling\_loss\_term\_ge1}. &
Machine-checked, but vacuous as a concrete sampling-distance estimate.  The
script comment explicitly marks this as a coarse fallback. \\
\midrule
4 &
Lines 1234 and 1218--1223:
\ec{ro\_exact\_nma\_hop}. &
The theorem consumes, rather than proves, the premise
\[
\Prb[\ec{ROExactHyperballPaperSimSampler}] +
\ec{rejection\_sampling\_loss\_term}
\le
\Prb[\ec{MLWE\_Game}] + \Prb[\ec{BimodalSelfTargetMSIS\_Game}].
\]
This premise is the cryptographic/sampler reduction surface for the RO
exact-hyperball NMA endpoint. &
Conditional premise.  The report must not describe it as derived in this final
lemma. \\
\midrule
5 &
Lines 1235--1236 and 1224--1226. &
\ec{mlwe\_bound} and \ec{msis\_bound} are forwarded unchanged.  The MSIS
premise is for the adapter game
\ec{ModuleSIS\_Game(BimodalMSIS\_As\_ModuleSIS(Bbimodal))}. &
External cryptographic assumption boundary expressed as theorem premises. \\
\midrule
6 &
Lines 1057--1069 and 746--766. &
The RO exact-hyperball theorem internally applies
\ec{euf\_cma\_security\_from\_haetae\_rejection\_paper\_sim\_nma\_and\_counted\_bimodal}.
The real-order step is \ec{ler\_trans}: first convert HAETAE rejection NMA to
RO exact-hyperball NMA with one rejection-loss term, then use
\ec{ro\_exact\_nma\_hop}. &
Machine-checked probability algebra and exact-hop composition. \\
\bottomrule
\end{longtable}
}

The phrase ``checked RO sampling hop'' should therefore be read narrowly.  The
EasyCrypt theorem is checked, and it imports a checked HopGames bound, but that
HopGames bound currently uses the fallback
\ec{rejection\_sampling\_loss\_term\_ge1}.  It is a valid formal implication,
not a non-vacuous concrete sampler-distance proof.

\paragraph{Replay for the real-signing exact-hyperball final theorem.}
The theorem
\ec{euf\_cma\_security\_from\_real\_signing\_exact\_hyperball\_hop\_and\_counted\_bimodal}
appears at \file{HAETAE\_Security.ec} lines 1239--1269.  Its statement says:
if real-signing paper simulation is bounded by exact-hyperball paper
simulation plus one rejection-sampling loss, if exact-hyperball paper
simulation plus that same loss is bounded by MLWE plus bimodal success, and if
MLWE/Module-SIS are bounded by their hardness terms, then the formal HAETAE
EUF-CMA game is bounded by \ec{haetae\_euf\_counted\_rom\_bound}.

{\scriptsize
\begin{longtable}{p{0.08\linewidth}p{0.25\linewidth}p{0.34\linewidth}p{0.25\linewidth}}
\toprule
\textbf{Step} & \textbf{EasyCrypt evidence} & \textbf{Proof obligation replay} &
\textbf{Status boundary} \\
\midrule
1 &
Lines 1259--1261:
\ec{move=> real\_exact\_hop exact\_nma\_hop mlwe\_bound msis\_bound} and
\ec{apply euf\_cma\_security\_from\_exact\_hyperball\_sampling\_hop\_and\_counted\_bimodal}. &
The theorem reduces to the exact-hyperball counted-ROM theorem at lines
1007--1034.  That theorem requires a HAETAE-rejection to exact-hyperball
sampling hop, an exact-hyperball NMA-to-MLWE/bimodal premise, and the two
hardness premises. &
Machine-checked proof replay step. \\
\midrule
2 &
Lines 1262--1264 and 819--841:
\ec{haetae\_rejection\_exact\_hyperball\_sampling\_hop\_from\_real\_signing\_exact\_hyperball\_hop}. &
The required HAETAE-rejection sampling hop is derived from the supplied
\ec{real\_exact\_hop}.  The proof rewrites real-signing paper simulation to
HAETAE-rejection paper simulation and exact-hyperball rejection paper
simulation to exact-hyperball paper simulation. &
Machine-checked exact-rewrite composition, but only conditional on the
\ec{real\_exact\_hop} premise. \\
\midrule
3 &
Lines 1240--1246:
\ec{real\_exact\_hop}. &
This premise is the non-RO sampler bridge:
\[
\Prb[\ec{RealSigningPaperSimSampler}]
\le
\Prb[\ec{ExactHyperballPaperSimSampler}]
+\ec{rejection\_sampling\_loss\_term}.
\]
The final theorem does not prove this sampling distance. &
Conditional premise.  It may be discharged by another lemma or external
argument, but not by lines 1239--1269 alone. \\
\midrule
4 &
Lines 1247--1252:
\ec{exact\_nma\_hop}. &
This premise bounds the exact-hyperball NMA endpoint plus one
\ec{rejection\_sampling\_loss\_term} by MLWE plus bimodal success.  It is the
reduction/hardness interface for the exact-hyperball endpoint. &
Conditional premise connecting sampler endpoint to hardness games. \\
\midrule
5 &
Lines 1253--1255. &
\ec{mlwe\_bound} and \ec{msis\_bound} are forwarded unchanged, with Module-SIS
stated for \ec{BimodalMSIS\_As\_ModuleSIS(Bbimodal)}. &
External MLWE and Module-SIS hardness assumption boundary. \\
\midrule
6 &
Lines 1028--1033 and 787--817. &
The exact-hyperball counted-ROM theorem applies
\ec{euf\_cma\_security\_from\_haetae\_rejection\_paper\_sim\_nma\_and\_counted\_bimodal}
after deriving a HAETAE-rejection NMA bound by \ec{ler\_trans}.  The
intermediate rewrite uses
\ec{rom\_internal\_nma\_exact\_hyperball\_rejection\_paper\_sim\_no\_fallback\_exact}. &
Machine-checked real-order arithmetic plus exact no-fallback rewrite. \\
\bottomrule
\end{longtable}
}

The second final theorem is therefore stronger as a replay surface for
non-vacuous sampling analysis: it exposes the real-signing-to-exact-hyperball
hop as an explicit premise instead of hiding it behind the current coarse RO
fallback.  It is also less self-contained: the user must supply or prove that
sampling-hop premise separately before the theorem gives a concrete end-to-end
security estimate.

\paragraph{Replay checklist for reimplementers.}
To replay either final theorem in EasyCrypt, a reimplementation should first
instantiate the module restrictions from lines 598--631, then prove or assume
the relevant NMA endpoint premise, then provide MLWE and Module-SIS hardness
bounds for the declared reductions, and only then invoke the final composition
lemma.  The only arithmetic performed in the final theorem bodies is forwarding
premises through earlier lemmas, \ec{ler\_trans}, exact \ec{rewrite} steps, and
SMT-dispatched probability facts such as \ec{mu\_bounded}.  None of these
steps proves MLWE hardness, Module-SIS hardness, SHAKE-to-ROM soundness,
implementation equivalence, or a non-vacuous rejection-sampler distance unless
the corresponding premise or imported lemma already supplies it.

\subsection{Non-vacuous replacement roadmap for the RO exact-hyperball fallback}

\paragraph{Fallback criticism.}
The current EasyCrypt scripts and the preceding report text do identify the
coarse fallback, but they still do not give enough information to replace it.
In \file{HAETAE\_HopGames.ec}, the comment at lines 7287--7290 says that the
intended replacement is a push-forward of the structural-to-exact-hyperball
attempt loss through \ec{paper\_sim\_sample\_from\_rejection\_attempt}, lifted
across the NMA wrapper once the lazy-ROM sampler law is exposed.  The checked
lemma immediately below, however, proves
\ec{rom\_internal\_nma\_ro\_signing\_attempt\_ro\_exact\_hyperball\_loss\_bound}
by a boundedness argument: the left probability is at most \(1\), the right
probability is nonnegative, and
\ec{rejection\_sampling\_loss\_term\_ge1} makes the additive loss at least
\(1\).  That is machine-checked, but it proves no meaningful closeness between
the RO-signing-attempt sampler and the RO exact-hyperball sampler.

Thus the report can truthfully say that the final RO-sampling theorem is
checked, but it cannot say that the present proof contains a non-vacuous
random-oracle sampler-distance theorem.  A reader trying to reimplement the
proof needs the following roadmap, not just the theorem name.

\paragraph{Desired replacement statement.}
The strongest drop-in replacement would prove the same inequality as the
current fallback, but without using \ec{rejection\_sampling\_loss\_term\_ge1}:
\[
\begin{aligned}
&\Prb[\ec{UF\_NMA}(\ec{ROMInternalTranscriptPaperSimAsNMA}
  (A,\ec{ROSigningAttemptPaperSimSampler}))] \\
&\quad \le
\Prb[\ec{UF\_NMA}(\ec{ROMInternalTranscriptPaperSimAsNMA}
  (A,\ec{ROExactHyperballPaperSimSampler}))] +
\ec{rejection\_sampling\_loss\_term}.
\end{aligned}
\]
If \ec{structural\_to\_exact\_hyperball\_paper\_sample\_loss\_obligation}
is not globally discharged, the honest replacement should expose it as an
antecedent, as the existing boundary lemmas already do.  If the proof must go
through budgeted wrappers, the resulting theorem may instead have
\ec{rom\_signature\_query\_budget * rejection\_sampling\_loss\_term} plus a
sampler-prequery term; in that case the final Security bound must be updated
rather than silently reusing the one-loss theorem.

\paragraph{Current fallback proof structure.}
{\scriptsize
\begin{longtable}{p{0.18\linewidth}p{0.29\linewidth}p{0.28\linewidth}p{0.17\linewidth}}
\toprule
\textbf{Script item} & \textbf{What the current script proves} &
\textbf{Why it is insufficient} & \textbf{Status} \\
\midrule
\file{HAETAE\_HopGames.ec} 7016--7034 &
States that a per-mode structural paper-sample loss obligation plus an
already-lifted NMA inequality gives the same NMA inequality. &
This is only an identity wrapper over the supplied lifted premise; it does not
lift the per-call sampler loss through adversarial NMA execution. &
Machine-checked implication with conditional premise. \\
\midrule
Lines 7057--7095 &
Creates a local one-step fact from
\ec{structural\_attempt\_exact\_hyperball\_paper\_sample\_one\_step\_loss},
but the final proof still closes by \ec{mu\_bounded} and a loss-at-least-one
argument. &
The one-step fact is not connected to the NMA result, so the proof remains
vacuous as a distribution comparison. &
Machine-checked coarse fallback. \\
\midrule
Lines 7097--7152 &
Splits sampled NMA success into a clean branch and the
\ec{sampler\_bad\_prequery} bad event, then charges the bad event to the
finite-event-logic bound
\ec{budgeted\_paper\_sim\_sampler\_bad\_prequery\_nma\_fel\_bound}. &
This is the right shape for a non-vacuous proof, but it requires a genuine
clean-lifting inequality as input. &
Machine-checked conditional bridge. \\
\midrule
Lines 7291--7314 &
Proves the unbudgeted RO-signing-attempt to RO exact-hyperball NMA bound by
\ec{left\_le1}, \ec{right\_ge0}, and
\ec{rejection\_sampling\_loss\_term\_ge1}. &
The result is true for any right-hand game with nonnegative probability; it
does not use sampler structure, ROM freshness, or the structural loss
obligation. &
Machine-checked but vacuous. \\
\midrule
Lines 7733--7815 &
Specializes the concrete \ec{HAETAE\_RO.FRO} budgeted NMA lifting path, but
the clean-conditioned lemma at lines 7733--7765 again uses boundedness and
loss-at-least-one. &
Concrete ROM specialization exists, yet the central clean sampler comparison
is still not non-vacuous. &
Machine-checked coarse fallback. \\
\bottomrule
\end{longtable}
}

\paragraph{Replacement proof obligations.}
{\scriptsize
\begin{longtable}{p{0.20\linewidth}p{0.36\linewidth}p{0.20\linewidth}p{0.16\linewidth}}
\toprule
\textbf{Obligation} & \textbf{Required content} & \textbf{Existing evidence} &
\textbf{Boundary} \\
\midrule
Per-call sampler push-forward &
Show that sampling with \ec{ROSigningAttemptPaperSimSampler.sample\_with\_seed}
at a fresh sampler site is bounded by
\ec{ROExactHyperballPaperSimSampler.sample\_with\_seed} plus
\ec{rejection\_sampling\_loss\_term}, for every postcondition on the resulting
\ec{paper\_sim\_signature\_sample}. &
The intended local facts are lines 6625--6676:
\ec{concrete\_ro\_signing\_attempt\_to\_exact\_hyperball\_sample\_with\_seed\_*}.
They use the structural paper-sample obligation and freshness. &
Machine-checked once all premises are available. \\
\midrule
Distributional push-forward &
Connect \ec{dsigning\_attempt\_state} and
\ec{dexact\_hyperball\_signing\_attempt\_state} through
\ec{paper\_sim\_sample\_from\_rejection\_attempt}; do not compare raw internal
states when the games observe paper samples or signatures. &
The obligation is named at lines 6615--6623.  Lines 8651--8722 show how attempt
or coin-sample-pair loss obligations imply the paper-sample obligation. &
Conditional unless the lower-level attempt or coin-sample-pair loss is proved. \\
\midrule
Lazy-ROM sampler law &
Expose that a fresh \ec{sampler\_expand\_query} returns signing coins with law
\ec{drandom\_coins}, and that projections through \ec{ro\_signing\_coins} are
the only component used by the sampler. &
\file{HAETAE\_ROM.ec} lines 99--109 prove
\ec{sampler\_expand\_query\_dro\_output\_ro\_signing\_coins}.  HopGames uses it
in sample-with-seed structural lemmas around lines 2285--2480. &
Machine-checked formal ROM fact; not SHAKE/FIPS evidence. \\
\midrule
Freshness and coverage invariant &
Maintain \ec{sampler\_rom\_covered} and show that a clean seed has not already
been queried by the adversary or sampler before the sampler call. &
Lines 168--199 give \ec{sampler\_rom\_covered\_fresh\_after\_clean\_seed} and
\ec{sampler\_rom\_covered\_old\_log\_clean\_boundary}; lines 6678--6954 package
old-log and self-logged clean surfaces. &
Machine-checked invariant family. \\
\midrule
Clean-event one-call lifting &
Lift the per-call sample comparison to one signing-oracle call under the clean
condition
\ec{! sampler\_bad\_prequery} and an active signing budget. &
Lines 7413--7731 provide projection facts between \ec{O.sign} and
\ec{sample\_with\_seed}, including active signing and exact-hyperball
projections. &
Partly machine-checked, but the NMA-level clean lifting still needs a
non-vacuous composition proof. \\
\midrule
Sampler-prequery bad event &
Split total NMA success into clean success plus \ec{sampler\_bad\_prequery},
then bound the bad event by query budgets and challenge support. &
Lines 2062--2223 prove
\ec{budgeted\_paper\_sim\_sampler\_bad\_prequery\_nma\_fel\_bound}; lines
7097--7152 already use the correct union-bound skeleton. &
Machine-checked bad-event bound. \\
\midrule
Budgeted-to-unbudgeted adapters &
Move from unbudgeted NMA wrappers to budgeted wrappers and back without
changing the observable success event, or explicitly charge any extra loss. &
The final Security budgeted theorem already exposes
\ec{attempt\_unbudgeted\_to\_budgeted} and
\ec{exact\_budgeted\_to\_unbudgeted} premises at lines 1170--1214. &
Conditional unless exact adapters are proved for the chosen theorem surface. \\
\midrule
Final real-order composition &
Use \ec{ler\_trans}, \ec{ler\_add}, \ec{mu\_sub}, and \ec{mu\_or} to combine
clean lifting and bad-event bounds, then feed the result into
\file{HAETAE\_Security.ec}. &
Existing scripts show the arithmetic pattern at HopGames lines 7097--7152 and
Security lines 1036--1070. &
Machine-checked arithmetic once the non-vacuous premises are supplied. \\
\bottomrule
\end{longtable}
}

\paragraph{Where the structural obligation enters.}
The structural-to-exact-hyperball premise must enter before any NMA-level
probability comparison: it is a one-call distributional statement over
\ec{dsigning\_attempt\_state} and
\ec{dexact\_hyperball\_signing\_attempt\_state}.  In the current script this is
visible at lines 6625--6653 and 6965--6984, where the per-call
\ec{sample\_with\_seed} and \ec{StructuralAttemptPaperSample.sample} bounds
invoke \ec{structural\_to\_exact\_hyperball\_paper\_sample\_loss\_obligation}.
A non-vacuous NMA theorem must propagate exactly that premise through signing
oracle calls, adversary calls, clean-event preservation, and sampler-prequery
splitting.  Introducing it only at the final \(\EUFCMA\) composition level is
too late unless an imported lemma has already performed this lifting.

\paragraph{Proof status after this roadmap.}
The proof-engineering spike found no direct non-vacuous replacement for the
unconditional theorem as written.  No new EasyCrypt proof was implemented,
because existing lemmas stop at either a one-call sampler comparison or an
identity wrapper over an already-supplied NMA-level premise.  The missing
machine-checked statement is one of the following theorem surfaces:
\[
\begin{aligned}
&\Prb[\ec{UF\_NMA}(\ec{ROMInternalTranscriptPaperSimAsNMA}
  (A,\ec{ROSigningAttemptPaperSimSampler}))] \\
&\quad \le
\Prb[\ec{UF\_NMA}(\ec{ROMInternalTranscriptPaperSimAsNMA}
  (A,\ec{ROExactHyperballPaperSimSampler}))] +
\ec{rejection\_sampling\_loss\_term},
\end{aligned}
\]
derived from
\ec{structural\_to\_exact\_hyperball\_paper\_sample\_loss\_obligation}
without \ec{rejection\_sampling\_loss\_term\_ge1}; or the budgeted clean-event
variant
\[
\begin{aligned}
&\Prb[\ec{UF\_NMA}(\ec{ROMInternalTranscriptBudgetedPaperSimAsNMA}
  (A,\ec{ROSigningAttemptPaperSimSampler})) :
  \ec{res} \land \neg\ec{sampler\_bad\_prequery}]\\
&\quad \le
\Prb[\ec{UF\_NMA}(\ec{ROMInternalTranscriptBudgetedPaperSimAsNMA}
  (A,\ec{ROExactHyperballPaperSimSampler}))] +
\ec{rom\_signature\_query\_budget}\cdot
\ec{rejection\_sampling\_loss\_term},
\end{aligned}
\]
plus exact or monotone budgeted-to-unbudgeted adapters.  In the second route,
the checked \ec{budgeted\_paper\_sim\_sampler\_bad\_prequery\_nma\_fel\_bound}
then supplies the bad-event charge, but it does not itself compare the samplers.

The exact local blockage is therefore not the marginal sampler law:
\ec{concrete\_ro\_signing\_attempt\_to\_exact\_hyperball\_sample\_with\_seed\_*}
and \ec{sampler\_expand\_query\_dro\_output\_ro\_signing\_coins} already give
the per-call ingredients.  The missing proof is the adaptive lifting through
\ec{ROMInternalTranscriptPaperSimAsNMA} or its budgeted variant, including
signing-oracle sequencing, clean-event preservation, freshness/coverage, and
loss accumulation over at most the signing-query budget.  Until that theorem is
added, the roadmap above is explanatory documentation.  The machine-checked
facts are the existing EasyCrypt lemmas cited by line range.  The structural
sampler-distance premise, SHAKE-to-ROM interpretation, and any implementation
equivalence remain outside the checked fallback lemma unless future EasyCrypt
code proves them or exposes them explicitly as assumptions.

\paragraph{Focused budgeted clean-event proof spike.}
The follow-up proof spike inspected the two existing \ec{fel} patterns in
\file{HAETAE\_HopGames.ec}: the sampler-prequery bound around lines
2069--2125 and the sampled-direct ROM-programming prequery bound around lines
4935--5000.  Both are machine-checked bad-event accumulators for one
execution.  They sum the probability of a flag becoming true over a counter,
but they do not compare two oracle implementations and do not propagate a
continuation predicate through \ec{A.forge}.  Consequently, the current
\ec{fel} surface cannot directly prove the adaptive lossy-oracle statement
needed for
\ec{concrete\_rom\_internal\_budgeted\_nma\_ro\_signing\_attempt\_ro\_exact\_hyperball\_clean\_conditioned\_lifting}.

The checked one-call ingredients are present.  In particular,
\ec{concrete\_budgeted\_o\_sign\_clean\_one\_call\_loss\_from\_self\_logged\_surface}
connects an active clean signing call using
\ec{ROSigningAttemptPaperSimSampler} to the corresponding
\ec{ROExactHyperballPaperSimSampler} call with one
\ec{rejection\_sampling\_loss\_term} charge, assuming the structural sampler
obligation, \ec{sampler\_rom\_covered}, a fresh self-logged seed, active
\ec{signing\_count}, public-key equality, and equality of the sampler secret
keys.  The source-side projection is
\ec{concrete\_budgeted\_ro\_signing\_attempt\_o\_sign\_to\_self\_logged\_sample\_active\_projection};
the target-side projection is
\ec{concrete\_budgeted\_ro\_exact\_hyperball\_sample\_with\_seed\_to\_o\_sign\_active\_projection}.

The proof-infrastructure spike sharpened the blocker further.  The current
one-call loss lemma quantifies over predicates of type
\ec{signature -> bool}.  The continuation needed inside an adaptive adversary
is not merely a predicate on the returned signature: after \ec{O.sign} returns,
the rest of \ec{A.forge} can depend on adversary private state, the lazy-ROM
table, public query and transcript logs, signing records, counters, and future
oracle calls.  Thus the missing induction principle must be stateful.  Its
forge-level shape is: for every continuation predicate \(C\) over the final
forge result and the public final state later used by \ec{UF\_NMA.main}, and
for related attempt/exact memories with
\[
  k =
  \ec{signature\_query\_budget\_count}
  - \ec{ROMInternalTranscriptBudgetedPaperSimAsNMA.signing\_count},
\]
prove that the attempt game on the clean branch is bounded by the exact game
plus \(k\cdot\ec{rejection\_sampling\_loss\_term}\).  The active
\ec{O.sign} premise must itself be stateful:
\[
\Pr[\ec{O\_attempt.sign} : R(\ec{res},\ec{post\_state}) \land
      \neg\ec{sampler\_bad\_prequery}]
\le
\Pr[\ec{O\_exact.sign} : R(\ec{res},\ec{post\_state})]
  + \ec{rejection\_sampling\_loss\_term},
\]
where \(R\) includes the post-call relation required by the recursive
continuation.  The existing signature-only one-call theorem is a component of
that active-call premise, not the adversary-level lifting principle.

The missing theorem is now isolated as a conditional EasyCrypt proof
obligation, not as an external cryptographic assumption.  A non-vacuous proof
would require an adaptive lifting rule with the following invariant before and
after every active signing-oracle call: both games expose the same
adversary-visible public key, hash log, transcript log, signing-record log,
adversary hash counter, and signing counter; the two sampler secret keys are
equal; the attempt side satisfies \ec{sampler\_rom\_covered}; the attempt side
is on the clean branch
\(\neg\ec{ROMInternalTranscriptBudgetedPaperSimAsNMA.sampler\_bad\_prequery}\);
and the next sampler seed has not already appeared in the adversary or sampler
logs.  At each active call, the one-call loss lemma should be instantiated with
the continuation success predicate induced by the rest of \ec{A.forge}; the
constant charge should then be summed over at most
\ec{signature\_query\_budget\_count} active calls to obtain
\ec{rom\_signature\_query\_budget} times
\ec{rejection\_sampling\_loss\_term}.

No new proof was added in this spike because the repository currently exposes
exact \ec{equiv}/\ec{byequiv} adversary-call patterns, local \ec{phoare}
inequalities, and bad-event \ec{fel} machinery, but not the required stateful
relational/probability-loss adversary-call rule.  The status is therefore: the
local one-call facts remain machine-checked; the clean-event adaptive lifting
theorem remains an unproved EasyCrypt proof obligation; the structural
sampler-distance premise remains a conditional premise; and SHAKE-to-ROM and
implementation-equivalence claims remain external to this checked theorem path.

\section{Feedback 3 criticism for assumptions, transcripts, and ROM programming}

Before this expansion, the report still did not give enough evidence for a
reader to understand \file{HAETAE\_Assumptions.ec},
\file{HAETAE\_Transcript.ec}, and \file{HAETAE\_ROM\_Programming.ec} at the
level required by the new \file{HAETAE\_HopGames.ec} and
\file{HAETAE\_Security.ec} companion sections.  The old roadmap named the
files and listed representative lemma names, but it did not explain which
objects were external assumption interfaces, which were checked adapter
equalities, which transcript predicates fed special soundness, and which
random-oracle flags were later charged to counted-ROM loss terms.

That gap matters because the downstream proof scripts rely on these files
without restating their mathematics.  The hop-game script assumes the reader
already knows the meaning of \ec{valid\_forking\_pair},
\ec{transcript\_valid}, \ec{actual\_signature\_programming\_site},
\ec{bad\_prequery}, \ec{bad\_reprogram}, and
\ec{bad\_min\_entropy}.  The security script assumes the reader knows that
\ec{MLWE\_Game}, \ec{ModuleSIS\_Game}, and
\ec{BimodalSelfTargetMSIS\_Game} are assumption surfaces rather than proved
cryptanalytic estimates.  Without the companion material below, the report
could still let a reader follow the names of the final theorem while missing
the logical boundary between checked reductions, external hardness premises,
conditional sampling premises, and explanatory prose.

\section{Script companion coverage for HAETAE Assumptions}

\file{HAETAE\_Assumptions.ec} defines the formal hardness surfaces and the
checked adapters between the bimodal self-target relation and Module-SIS.  It
does not estimate concrete MLWE or Module-SIS security.  Its role is to give
the rest of the proof typed games and exact interface adapters.

\begin{longtable}{p{0.30\linewidth}p{0.60\linewidth}}
\toprule
\textbf{Script item} & \textbf{Mathematical and checker role} \\
\midrule
\ec{mlwe\_instance},
\ec{dmlwe\_real},
\ec{dmlwe\_random} &
Abstract MLWE instance space and two lossless distributions.  The script does
not define their concrete algebraic distribution; it exposes a typed
distinguishing game surface.  Losslessness is checked from the declared
\ec{[lossless]} annotations. \\
\ec{module\_sis\_instance},
\ec{module\_sis\_solution},
\ec{module\_sis\_valid} &
Module-SIS instances are pairs \((A,t)\), solutions are \ec{polyvecl} vectors,
and validity is the equation \(A\cdot s=t\) plus well-formedness.  The lemmas
around \ec{module\_sis\_eval} and \ec{module\_sis\_zero\_valid} are checked
algebraic sanity facts. \\
\ec{bimodal\_selftarget\_msis\_instance} and solution aliases &
The bimodal self-target relation is represented by the same carrier types as
Module-SIS in this model.  This is a formal interface decision, not a
cryptanalytic theorem. \\
\ec{bimodal\_to\_module\_sis\_instance},
\ec{bimodal\_to\_module\_sis\_solution},
\ec{bimodal\_to\_module\_sis\_valid} &
Identity adapter from bimodal self-target MSIS to Module-SIS.  EasyCrypt checks
that validity is preserved by unfolding definitions. \\
\ec{MLWE\_Distinguisher},
\ec{MLWE\_Real\_Game},
\ec{MLWE\_Random\_Game} &
Classical distinguishing surface for MLWE.  These games sample from the
abstract distributions and call an adversary.  No hardness bound is proved
here; bounds enter later as premises or as \ec{mlwe\_hardness\_term}. \\
\ec{ModuleSIS\_Solver} and
\ec{ModuleSIS\_Relation\_Game} &
Relation-game surface for a solver that outputs an optional Module-SIS
solution.  Success means \ec{sol <> None} and \ec{module\_sis\_valid}. \\
\ec{BimodalSelfTargetMSIS\_Solver} and relation game &
Analogous relation-game surface for the bimodal self-target relation.  The
success event is later related exactly to Module-SIS through the identity
adapter. \\
\ec{bimodal\_solver\_lift\_exact} and
\ec{bimodal\_mode\_solver\_lifted\_distribution\_exact} &
\ec{byequiv} proofs showing that solving the bimodal relation and then viewing
the result as Module-SIS has the same success probability as the lifted
Module-SIS game.  These are checked equivalences, not hardness estimates. \\
\ec{MLWE\_Reduction},
\ec{BimodalSelfTargetMSIS\_Reduction},
\ec{ModuleSIS\_Reduction} &
One-procedure reduction interfaces used by the final theorem family.  They are
opaque modules whose success probabilities are bounded outside this file. \\
\ec{MLWE\_Game},
\ec{BimodalSelfTargetMSIS\_Game},
\ec{ModuleSIS\_Game} &
Thin wrappers around the corresponding reduction modules.  They exist so that
security theorems can state probabilities uniformly. \\
\ec{BimodalMSIS\_As\_ModuleSIS} and
\ec{bimodal\_msis\_as\_module\_sis\_exact} &
Final exact adapter used in \file{HAETAE\_Security.ec}.  The proof is a
\ec{byequiv} that calls the same underlying module on both sides. \\
\bottomrule
\end{longtable}

The assumption boundary is therefore exact: checked material in this file
includes type-correct games, adapter definitions, algebraic validity lemmas,
losslessness of declared distributions, and equivalence lemmas for adapters.
External material includes the actual hardness of MLWE, Module-SIS, and the
bimodal self-target relation.  The final security file may assume or compose
bounds for \ec{Pr[MLWE\_Game(...).main():res]} and
\ec{Pr[ModuleSIS\_Game(...).main():res]}, but those bounds are not derived in
\file{HAETAE\_Assumptions.ec}.

\section{Script companion coverage for HAETAE Transcript}

\file{HAETAE\_Transcript.ec} defines the transcript vocabulary used by
random-oracle programming and special soundness.  It is the bridge between a
signature object and the public proof transcript extracted from that signature.

\begin{longtable}{p{0.30\linewidth}p{0.60\linewidth}}
\toprule
\textbf{Transcript family} & \textbf{Mathematical reading} \\
\midrule
\ec{transcript} &
An eight-component tuple consisting of public key, message, context,
commitment high bits, commitment low bits, message hash, challenge, and
signature. \\
Projection operations &
\ec{transcript\_pk}, \ec{transcript\_message},
\ec{transcript\_context}, \ec{transcript\_commitment\_highbits},
\ec{transcript\_commitment\_lowbits}, \ec{transcript\_message\_hash},
\ec{transcript\_challenge}, and \ec{transcript\_signature} are tuple
projections.  Their role is readability and stable proof targets. \\
\ec{transcript\_of\_signature} &
Builds a transcript from a signature by extracting the commitment, message-hash,
and challenge fields from the signature.  This is used whenever signing games
append a transcript to their log. \\
\ec{transcript\_challenge\_query} and
\ec{transcript\_signature\_challenge\_query} &
Two views of the challenge random-oracle query: one from transcript fields and
one from signature fields.  Valid transcripts make them equal. \\
\ec{transcript\_verifies},
\ec{transcript\_fields\_match\_signature},
\ec{transcript\_valid} &
Validity means the signature verifies, the signature is well-formed, and the
stored transcript fields match the signature fields.  This is the predicate
used by ROM programming and special-soundness lemmas. \\
\ec{transcript\_challenge\_matches} and
\ec{transcript\_signature\_challenge\_matches} &
Relate a random-oracle output to the transcript or signature challenge.  Valid
transcripts let the proof replace one predicate by the other. \\
\ec{same\_fork\_target} and
\ec{distinct\_fork\_challenges} &
Forking-pair structure: two transcripts must share public key, message,
context, commitment high bits, commitment low bits, and message hash, while
having different challenges. \\
\ec{valid\_forking\_pair} &
The exact precondition for extraction: same target, distinct challenges, and
both transcripts valid.  This avoids the logical error of claiming extraction
from one accepting transcript. \\
\ec{forking\_difference\_solution},
\ec{fork\_module\_matrix},
\ec{fork\_module\_target} &
Build the bimodal self-target MSIS instance and candidate solution from a
valid pair of transcripts.  The solution is the response difference. \\
\ec{extract\_bimodal\_selftarget\_msis\_solution} and
\ec{extract\_module\_sis\_solution} &
Partial extractors that return \ec{None} unless the pair is valid.  The
Module-SIS extractor applies the adapter from \file{HAETAE\_Assumptions.ec}. \\
Honest-transcript operations &
\ec{transcript\_from\_honest\_signing},
\ec{honest\_transcript\_verifies}, and
\ec{honest\_transcript\_challenge} connect deterministic internal signing to
the transcript vocabulary consumed by ROM programming. \\
\bottomrule
\end{longtable}

The lemma families in this file are proof-script glue for downstream scripts.

\begin{longtable}{p{0.30\linewidth}p{0.60\linewidth}}
\toprule
\textbf{Lemma family} & \textbf{What EasyCrypt checks} \\
\midrule
Signature-field consistency &
\ec{transcript\_fields\_match\_signature\_challenge\_query},
\ec{transcript\_fields\_match\_signature\_challenge\_matches},
\ec{transcript\_valid\_signature\_challenge\_query},
\ec{transcript\_valid\_signature\_challenge\_matches}, and
\ec{transcript\_valid\_signature\_challengeE} unfold definitions to show that
valid transcript fields and signature fields agree. \\
Public-equation and reconstruction lemmas &
\ec{transcript\_valid\_public\_equation},
\ec{transcript\_valid\_commitment\_reconstruction},
\ec{transcript\_reconstructionE}, and
\ec{transcript\_active\_reconstructionE} justify using verification equations
as public reconstruction equations in special-soundness reasoning. \\
Forking-pair query and challenge lemmas &
\ec{valid\_forking\_pair\_challenge\_query},
\ec{valid\_forking\_pair\_signature\_challenge\_query}, and
\ec{valid\_forking\_pair\_distinct\_signature\_challenges} convert the abstract
forking-pair predicate into equal challenge-query sites and distinct challenge
values. \\
Extraction obligation lemmas &
\ec{extract\_bimodal\_forking\_some},
\ec{extract\_bimodal\_forking\_nonempty},
\ec{bimodal\_forking\_solution\_valid}, and
\ec{forking\_extraction\_obligation\_holds} prove that a valid forking pair
produces a valid bimodal solution. \\
Module-SIS lift obligations &
\ec{bimodal\_to\_module\_sis\_lift\_obligation\_holds},
\ec{module\_sis\_extraction\_from\_bimodal}, and
\ec{special\_soundness\_from\_bimodal\_lift} package the adapter from bimodal
extraction to Module-SIS extraction. \\
Public reconstruction obligations &
\ec{fork\_public\_reconstruction\_obligation\_holds} and
\ec{special\_soundness\_with\_public\_reconstruction\_from\_bimodal\_lift}
record that public reconstruction accompanies extraction when later proofs need
both facts. \\
\bottomrule
\end{longtable}

The checked content here is definitional and algebraic: projections, tuple
equality, validity propagation, extraction from a valid pair, and adapter
validity.  The explanatory boundary is important.  This file does not prove a
forking lemma that produces two transcripts from an adversary; it only states
and checks what follows once a valid forking pair is available.

\subsection{Lemma-ledger companion for HAETAE Transcript}

\paragraph{Lemma-ledger criticism for HAETAE Transcript.}
The preceding companion section explains the purpose of
\file{HAETAE\_Transcript.ec}, but by itself it is still not sufficient for a
mathematician or cryptologist to reconstruct the proof obligations in that
script.  The evidence is in the script structure: lines 50--112 define several
different challenge-query, matching, validity, and forking predicates, while
lines 470--604 turn those predicates into extraction and special-soundness
obligations.  Without a ledger, a reader can easily conflate a checked
definition such as \ec{valid\_forking\_pair}, a theorem-boundary predicate such
as \ec{special\_soundness\_obligation}, and the external cryptographic meaning
of the Module-SIS hardness game imported through \file{HAETAE\_Assumptions.ec}.
The report plus script can support reconstruction only if the reader is told
which operators are pure transcript syntax, which lemmas are checked algebraic
consequences, and which adversarial production of forks is left to downstream
HopGames reductions.

\begin{longtable}{p{0.20\linewidth}p{0.25\linewidth}p{0.16\linewidth}p{0.18\linewidth}p{0.15\linewidth}}
\toprule
\textbf{EasyCrypt item and lines} & \textbf{Mathematical statement} &
\textbf{Category} & \textbf{Premises and imports} & \textbf{Downstream role} \\
\midrule
\ec{AllCore}, \ec{List}, \ec{HAETAE\_Params}, \ec{HAETAE\_Algebra},
\ec{HAETAE\_ROM}, \ec{HAETAE\_Assumptions}; 1--9 &
Imports the tuple/list environment, HAETAE algebra, ROM query vocabulary, and
SIS adapter vocabulary used by all transcript predicates. &
Imported/library fact; external-assumption boundary for the hardness-side
interfaces. &
No local proof; relies on imported EasyCrypt and project theories. &
Makes transcript extraction compatible with HopGames and the final
\file{HAETAE\_Security.ec} reduction chain. \\
\midrule
\ec{transcript}; 11--12 &
Models a transcript as public key, message, context, commitment high bits,
commitment low bits, message hash, challenge, and signature. &
Machine-checked fact. &
Imports HAETAE carrier types from algebra and parameters. &
Carrier for transcript logs, forking games, and special-soundness extraction
in \file{HAETAE\_HopGames.ec}. \\
\midrule
\ec{signature\_commitment\_highbits},
\ec{signature\_commitment\_lowbits}, \ec{signature\_message\_hash},
\ec{signature\_challenge}; 14--29 &
Defines the transcript-visible fields extracted from a signature, ignoring
mode, key, message, and context except for typing. &
Machine-checked fact. &
Uses \ec{sig\_*} projections from the algebraic signature model. &
Lets later proofs compare stored transcript fields with signature fields
without unpacking signatures repeatedly. \\
\midrule
\ec{transcript\_pk} through \ec{transcript\_signature},
\ec{transcript\_of\_signature}; 31--48 &
Gives tuple projections and the canonical transcript obtained from a public
signature. &
Machine-checked fact. &
Definitional tuple projections plus the signature-field operators above. &
Used by signing-record logs and ROM programming sites in HopGames. \\
\midrule
\ec{transcript\_challenge\_query},
\ec{transcript\_signature\_challenge\_query}; 50--63 &
Builds the challenge random-oracle query from either stored transcript fields
or the embedded signature fields. &
Machine-checked fact. &
Uses \ec{challenge\_hash\_query} from \file{HAETAE\_ROM.ec}. &
Supplies the query equality needed by ROM programming and valid-fork query
alignment. \\
\midrule
\ec{transcript\_verifies}, \ec{transcript\_challenge\_matches},
\ec{transcript\_signature\_challenge\_matches},
\ec{transcript\_fields\_match\_signature}, \ec{transcript\_valid},
\ec{transcript\_public\_equation}; 65--95 &
Defines verification, output-challenge matching, field consistency, full
transcript validity, and the public verification equation. &
Machine-checked fact. &
Uses \ec{verify\_internal}, \ec{valid\_signature},
\ec{ro\_challenge\_hash}, and \ec{verify\_public\_equation}. &
Defines the invariant consumed by valid-forking-pair and transcript-log soundness
lemmas in HopGames. \\
\midrule
\ec{same\_fork\_target}, \ec{distinct\_fork\_challenges},
\ec{valid\_forking\_pair}; 97--112 &
A valid fork is two valid transcripts for the same public target with distinct
challenges. &
Conditional premise and machine-checked predicate. &
Requires both transcripts to satisfy \ec{transcript\_valid}. &
Primary precondition for extraction games; no lemma here proves that an
adversary produces such a pair. \\
\midrule
\ec{transcript\_response}, \ec{transcript\_response\_aux},
\ec{transcript\_public\_key\_challenge\_term},
\ec{active\_public\_reconstruction},
\ec{inactive\_public\_reconstruction}; 114--129 &
Extracts response components and the public-key challenge contribution used in
verification reconstruction. &
Machine-checked fact. &
Imports signature response projections, \ec{public\_key\_challenge\_term},
and \ec{polyveck\_zero}. &
Feeds public-reconstruction lemmas and special-soundness obligations. \\
\midrule
\ec{forking\_difference\_solution},
\ec{fork\_response\_aux\_difference},
\ec{fork\_public\_challenge\_relation},
\ec{fork\_response\_aux\_relation},
\ec{fork\_left\_active\_public\_challenge\_relation},
\ec{fork\_right\_active\_public\_challenge\_relation},
\ec{fork\_public\_reconstruction\_cases}; 131--166 &
Defines the response difference and public reconstruction relations obtained
from two forked transcripts. &
Machine-checked fact. &
Uses vector subtraction/addition and public-key challenge terms from algebra. &
Local algebraic vocabulary for later extraction and reconstruction obligations. \\
\midrule
\ec{fork\_matrix\_source}, \ec{fork\_module\_matrix},
\ec{fork\_module\_target},
\ec{bimodal\_selftarget\_msis\_instance\_of\_fork},
\ec{extract\_bimodal\_selftarget\_msis\_solution},
\ec{module\_sis\_instance\_of\_fork},
\ec{extract\_module\_sis\_solution}; 168--215 &
Constructs the bimodal self-target MSIS instance, candidate solution, and
Module-SIS adapter output from a valid fork. &
Machine-checked fact with external-assumption boundary for the hardness game. &
Uses deterministic matrix generation and \ec{bimodal\_to\_module\_sis\_*}
from \file{HAETAE\_Assumptions.ec}. &
Provides the objects returned by HopGames special-soundness and Module-SIS
reduction modules. \\
\midrule
\ec{bimodal\_extraction\_success}; 216--220 &
States that bimodal extraction returns a nonempty solution and that the
solution validates for the fork-derived instance. &
Theorem-boundary predicate. &
Depends on the extractor and \ec{bimodal\_selftarget\_msis\_valid}. &
Used to package extraction success before lifting to Module-SIS. \\
\midrule
\ec{extract\_bimodal\_forking\_some},
\ec{extract\_bimodal\_forking\_nonempty}; 222--238 &
A valid forking pair makes the bimodal extractor return exactly the response
difference, hence not \ec{None}. &
Machine-checked fact. &
Premise: \ec{valid\_forking\_pair md tr1 tr2}. &
Internal support for \ec{forking\_extraction\_obligation\_holds}. \\
\midrule
\ec{transcript\_fields\_match\_signature\_challenge\_*},
\ec{transcript\_valid\_signature\_challenge\_*}; 240--288 &
Field consistency and transcript validity make the stored challenge query,
signature challenge query, and challenge-match predicates agree. &
Machine-checked fact. &
Premises: \ec{transcript\_fields\_match\_signature} or
\ec{transcript\_valid}. &
Required by ROM programming to switch between transcript and signature views
of the programmed challenge site. \\
\midrule
\ec{transcript\_valid\_public\_equation},
\ec{transcript\_valid\_commitment\_reconstruction}; 299--331 &
A valid transcript satisfies the public verification equation and reconstructs
the stored commitment high bits. &
Machine-checked fact. &
Unfolds \ec{verify\_internal}, \ec{verify\_public\_equation}, and transcript
validity. &
Connects transcript validity to public reconstruction used in special
soundness. \\
\midrule
\ec{valid\_forking\_pair\_*}; 333--461 &
Forking pairs have public equations, common challenge-query sites, distinct
signature challenges, equal reconstructed high bits, and public challenge
relations. &
Machine-checked fact family. &
Premise: \ec{valid\_forking\_pair}; relies on the field-query and public
reconstruction lemmas above. &
Supplies the proof obligations consumed by HopGames' special-soundness
interface readiness lemmas. \\
\midrule
\ec{extraction\_success}, \ec{forking\_extraction\_obligation},
\ec{fork\_public\_reconstruction\_obligation}; 470--484 &
Defines the Module-SIS extraction-success predicate and the two universal
obligations over all valid forking pairs. &
Conditional premise and theorem boundary. &
Quantifies over \ec{valid\_forking\_pair}; uses Module-SIS validity. &
Named assumptions/targets for downstream interface readiness; they are not
adversarial fork-generation lemmas. \\
\midrule
\ec{forking\_difference\_solution\_wf},
\ec{fork\_module\_matrix\_rows\_wf},
\ec{fork\_module\_matrix\_size},
\ec{bimodal\_forking\_solution\_valid},
\ec{forking\_extraction\_obligation\_holds},
\ec{fork\_public\_reconstruction\_obligation\_holds}; 486--527 &
Proves well-formedness, matrix size, bimodal solution validity, and the two
local fork obligations. &
Machine-checked fact. &
Uses algebraic well-formedness, \ec{mode\_k}, and previous fork lemmas. &
Discharges the transcript-side premises used by HopGames special-soundness
wrappers. \\
\midrule
\ec{bimodal\_to\_module\_sis\_lift\_obligation},
\ec{bimodal\_to\_module\_sis\_lift\_obligation\_holds}; 535--549 &
States and proves that every valid bimodal self-target solution lifts to a
valid Module-SIS solution under the imported adapter. &
Machine-checked fact, imported-interface boundary. &
Uses \ec{bimodal\_to\_module\_sis\_valid} from assumptions. &
Bridge from bimodal extraction to the Module-SIS game used in final security. \\
\midrule
\ec{special\_soundness\_obligation},
\ec{special\_soundness\_with\_public\_reconstruction\_obligation},
\ec{module\_sis\_extraction\_from\_bimodal},
\ec{special\_soundness\_*}; 551--604 &
Packages bimodal extraction plus the lift obligation into Module-SIS
special-soundness, optionally preserving public reconstruction. &
Machine-checked fact; theorem boundary. &
Premises are the fork, lift, and public-reconstruction obligations explicitly
listed in the lemmas. &
Used by HopGames to certify the special-soundness interface before composing
with NMA/CMA reductions. \\
\midrule
\ec{transcript\_from\_honest\_signing},
\ec{honest\_transcript\_verifies},
\ec{transcript\_of\_signature\_*},
\ec{honest\_transcript\_challenge}; 607--660 &
Builds the transcript produced by honest signing and proves it verifies, has
fields matching the signature, and carries the expected challenge hash. &
Machine-checked fact. &
Uses \ec{sign\_internal}, \ec{verify\_internal\_sign\_internal}, and signature
field projections. &
Supports signing-record soundness, min-entropy no-failure facts, and transcript
logging in ROM programming and HopGames. \\
\bottomrule
\end{longtable}

\section{Script companion coverage for HAETAE ROM Programming}

\file{HAETAE\_ROM\_Programming.ec} defines the counted random-oracle
programming surface consumed by \file{HAETAE\_HopGames.ec}.  It is where
transcript challenge sites become programmable ROM sites and where bad-event
flags are named.

\begin{longtable}{p{0.30\linewidth}p{0.60\linewidth}}
\toprule
\textbf{ROM-programming object} & \textbf{Mathematical reading} \\
\midrule
\ec{programming\_site} &
A pair \((q,y)\) consisting of a random-oracle query and the output to program
at that query. \\
\ec{programming\_site\_of\_transcript},
\ec{signature\_programming\_site\_of\_transcript},
\ec{actual\_signature\_programming\_site} &
Ways to identify the challenge site induced by a transcript or by the
signature fields inside that transcript.  Validity lemmas show the transcript
and signature views agree. \\
\ec{honest\_signing\_programming\_site} &
The challenge site generated by honest signing coins.  This is the random
site whose freshness is needed for safe ROM programming. \\
\ec{programmed\_site\_query\_min\_entropy\_bound} &
Point-probability bound for honest programming-site queries over signing
coins.  This is the min-entropy input used for prequery bounds. \\
\ec{programming\_site\_matches\_transcript} and
\ec{reprogramming\_failure} &
A programmed site is correct for a transcript when it has the transcript's
challenge query and output matching the transcript challenge.  Failure is the
negation of that match. \\
\ec{challenge\_entropy\_support\_ok},
\ec{min\_entropy\_failure},
\ec{transcript\_log\_min\_entropy\_clear} &
Challenge-support predicates.  They are used to prove that honest transcripts
do not trigger the min-entropy bad flag. \\
\ec{fs\_with\_aborts\_failure} and
\ec{transcript\_log\_signature\_programming\_sites\_clear} &
Legacy Fiat-Shamir-with-aborts failure combines reprogramming failure and
min-entropy failure.  The clear-log predicate says no transcript in the signing
log fails at its actual signature programming site. \\
\ec{counted\_rom\_event} and \ec{counted\_trace\_*} &
Event vocabulary for hash queries, signature sites, programmed sites, and
min-entropy checks.  These are counters and traces, not probability bounds by
themselves. \\
\ec{ro\_query\_is\_challenge} and
\ec{challenge\_query\_log\_count} &
Classifies query constructors and counts distinct challenge queries in a log.
Message, matrix-expansion, and sampler-expansion queries are not challenge
queries. \\
\ec{programming\_site\_prequeried},
\ec{programming\_site\_reprograms},
\ec{programming\_site\_fresh} &
The three core bad-event predicates: a query was already asked, a site conflicts
with an earlier programmed output for the same query, or neither bad event
holds. \\
\ec{rom\_counted\_programming\_bad} and
\ec{rom\_counted\_state\_bad} &
Boolean summary of the counted bad state: min-entropy failure, prequery, or
reprogramming conflict. \\
\bottomrule
\end{longtable}

The stateful ROM module is central to the downstream proof.

\begin{longtable}{p{0.30\linewidth}p{0.60\linewidth}}
\toprule
\textbf{\ec{CountedLazyROM} state or method} & \textbf{Meaning} \\
\midrule
\ec{hash\_queries} &
Log of queries passed to \ec{get}.  Challenge-query budgets are derived from
this log. \\
\ec{programmed\_sites} &
Sites passed to \ec{program}.  Reprogramming is detected by comparing a new
site against this list. \\
\ec{signature\_sites} &
Sites observed from signing transcripts.  Hop games use this to connect
transcript validity and ROM programming. \\
\ec{bad\_prequery} &
Set when \ec{program(site)} is called after the site's query already appears in
\ec{hash\_queries}. \\
\ec{bad\_reprogram} &
Set when \ec{program(site)} conflicts with a previously programmed output for
the same query. \\
\ec{bad\_min\_entropy} &
Set when observing a transcript whose challenge is outside the expected sparse
support. \\
\ec{init}, \ec{get}, \ec{program},
\ec{observe\_signature\_site},
\ec{observe\_transcript}, \ec{bad} &
\ec{init} clears logs and flags; \ec{get} logs a query; \ec{program} sets
prequery/reprogram flags and calls the underlying oracle's \ec{set};
\ec{observe\_transcript} records the induced signature site and updates the
min-entropy flag; \ec{bad} returns the disjunction of the three flags. \\
\bottomrule
\end{longtable}

The lemma families needed by \file{HAETAE\_HopGames.ec} and
\file{HAETAE\_Security.ec} are:

\begin{longtable}{p{0.30\linewidth}p{0.60\linewidth}}
\toprule
\textbf{Lemma family} & \textbf{Proof-script reading} \\
\midrule
\ec{CountedLazyROMStateFacts} &
Hoare facts showing that \ec{init} clears flags, \ec{get} preserves clear
flags or min-entropy clear state, \ec{program} preserves bad-clear state when a
site is fresh, and \ec{bad} is false when all flags are clear.  These are state
invariants, not loss bounds. \\
Programming-site self-match lemmas &
\ec{programming\_site\_self\_matches},
\ec{programming\_site\_self\_no\_reprogramming\_failure}, and
\ec{fs\_with\_aborts\_failureE} unfold the definitions connecting a transcript,
an oracle output, and a programming site. \\
Challenge-query counting lemmas &
\ec{challenge\_query\_log\_count\_cons\_le} and the nonchallenge-cons lemmas
show how query logs affect the counted challenge budget.  These facts are used
inside budgeted HopGames invariants. \\
Honest-site entropy and spread lemmas &
\ec{honest\_signing\_programming\_site\_query\_entropy\_token},
\ec{signing\_distribution\_programming\_site\_query\_spread}, and
\ec{honest\_signing\_programming\_site\_query\_spread\_bound} reduce
programming-site min-entropy to the signing-coin token-spread assumption from
the distribution layer. \\
Prequery probability bounds &
\ec{programmed\_site\_prequery\_one\_step\_*},
\ec{direct\_signing\_coin\_sampler\_prequery\_bound}, and
\ec{direct\_signing\_coin\_sampler\_message\_hash\_prequery\_bound} bound the
probability that a freshly sampled signing site was already queried.  These
are the local ingredients for counted-ROM prequery terms. \\
Reprogramming conflict lemmas &
\ec{honest\_signing\_programming\_sites\_no\_conflict},
\ec{programming\_site\_log\_conflict\_free\_for\_honest\_*},
\ec{programmed\_site\_reprogram\_one\_step\_conflict\_free\_zero}, and
\ec{programmed\_site\_prequery\_reprogram\_one\_step\_concrete\_bound} show
that honest sites do not conflict and combine reprogramming with prequery
bounds. \\
Counted bad-flag budget lemmas &
\ec{counted\_lazy\_rom\_bad\_flags\_budget\_sound},
\ec{counted\_lazy\_rom\_bad\_flags\_query\_budget\_bound}, and
\ec{counted\_lazy\_rom\_bad\_flags\_query\_budget\_bound\_subunit} express the
condition under which the probability of \ec{CountedLazyROMBadGame} is charged
to \ec{counted\_rom\_programming\_loss\_term}. \\
Transcript no-failure lemmas &
\ec{honest\_transcript\_no\_min\_entropy\_failure},
\ec{signing\_transcript\_relation\_no\_min\_entropy\_failure},
\ec{signing\_record\_log\_sound\_transcripts\_no\_min\_entropy}, and
\ec{fs\_with\_aborts\_no\_failure\_*} explain why honest signing records do not
trigger the legacy failure predicate. \\
Valid-transcript programming-site lemmas &
\ec{transcript\_valid\_signature\_programming\_site\_matches},
\ec{transcript\_valid\_actual\_signature\_programming\_site\_matches}, and
\ec{valid\_forking\_pair\_programming\_site\_query} connect transcript validity
and forking pairs to identical programming-query sites. \\
Bound-obligation lemmas &
\ec{counted\_rom\_programming\_bound\_obligation\_holds} and
\ec{fs\_with\_aborts\_bound\_obligation\_holds} prove only nonnegativity of
the named loss terms.  They do not prove that the loss is small; the reductions
file records separately that the old rejection ledger is vacuous. \\
\bottomrule
\end{longtable}

The checked content in \file{HAETAE\_ROM\_Programming.ec} is therefore a set of
definitions, Hoare invariants, and probability bounds for explicit bad events.
The file does not justify replacing SHAKE by a random oracle, does not give a
concrete MLWE/SIS estimate, and does not by itself complete the final
EUF-CMA theorem.  It supplies the exact bad-event vocabulary and counted-loss
interfaces that HopGames and Security compose.

\subsection{Lemma-ledger companion for HAETAE ROM Programming}

\paragraph{Lemma-ledger criticism for HAETAE ROM Programming.}
The existing ROM-programming companion describes the main object families, but
it still leaves a reconstruction gap.  The script is not merely a dictionary of
bad-event names: it defines a stateful counted ROM module at lines 250--301,
Hoare invariants at lines 319--500, min-entropy and prequery bounds at
lines 617--1386, bad-flag budget wrappers at lines 1455--1478, and
valid-transcript programming-site closure lemmas at lines 1639--1723.  A reader
who sees only the family names cannot determine which statements are checked
probability bounds, which are conditional on query-budget hypotheses, which are
formal ROM facts, and which are external SHAKE/FIPS modeling assumptions.  The
ledger below therefore separates formal ROM programming from implementation
hash assumptions and from the final cryptographic hardness composition.

\begin{longtable}{p{0.20\linewidth}p{0.25\linewidth}p{0.16\linewidth}p{0.18\linewidth}p{0.15\linewidth}}
\toprule
\textbf{EasyCrypt item and lines} & \textbf{Mathematical statement} &
\textbf{Category} & \textbf{Premises and imports} & \textbf{Downstream role} \\
\midrule
\ec{AllCore}, \ec{Real}, \ec{Distr}, \ec{List}, \ec{FSet},
\ec{PROM}, \ec{StdOrder}, \ec{Mu\_mem}, and local imports; 1--16 &
Imports probability, finite-set, random-oracle, transcript, rejection, and
loss-term infrastructure. &
Imported/library fact; external ROM-assumption boundary. &
No local proof; relies on EasyCrypt libraries and project theories. &
Defines the formal ROM environment used by HopGames and indirectly by
\file{HAETAE\_Security.ec}. \\
\midrule
\ec{programming\_site}, \ec{programming\_site\_query},
\ec{programming\_site\_output}; 18--22 &
A programmable site is a pair of a ROM query and the output to install at that
query. &
Machine-checked fact. &
Uses \ec{ro\_query} and \ec{ro\_output} from \file{HAETAE\_ROM.ec}. &
Base type for every prequery, reprogramming, and transcript-site predicate. \\
\midrule
\ec{programming\_site\_of\_transcript},
\ec{signature\_programming\_site\_of\_transcript},
\ec{transcript\_signature\_challenge\_output},
\ec{actual\_signature\_programming\_site}; 23--37 &
Builds the challenge programming site from transcript fields, signature fields,
or the signature's own challenge output. &
Machine-checked fact. &
Depends on transcript challenge queries and \ec{ro\_output\_of\_challenge}. &
Used when HopGames switches between transcript and signature views of a
programmed challenge. \\
\midrule
\ec{honest\_signing\_programming\_site},
\ec{honest\_signing\_programming\_site\_query}; 39--50 &
Defines the programming site induced by an honest signing transcript generated
from secret key, message, context, and signing coins. &
Machine-checked fact. &
Uses \ec{transcript\_from\_honest\_signing}. &
Source of the randomly sampled site in prequery and reprogramming-loss lemmas. \\
\midrule
\ec{programmed\_site\_query\_min\_entropy\_bound\_for},
\ec{programmed\_site\_query\_min\_entropy\_bound}; 51--60 &
States that every query in a log has at most a given point probability as an
honest signing programming-site query. &
Conditional premise. &
Uses \ec{mu} over a signing-coin distribution; default instance is
\ec{drandom\_coins}. &
Input condition for one-step prequery probability bounds in HopGames. \\
\midrule
\ec{programming\_site\_matches\_transcript},
\ec{reprogramming\_failure},
\ec{challenge\_entropy\_support\_ok},
\ec{min\_entropy\_failure}; 65--77 &
Defines correct transcript programming, reprogramming failure, sparse challenge
support, and the min-entropy bad event. &
Machine-checked predicate; external modeling boundary for entropy support. &
Uses transcript challenge matching and \ec{challenge\_sparse}. &
Bad-event vocabulary for programmed transcript games. \\
\midrule
\ec{transcript\_log\_min\_entropy\_clear},
\ec{signing\_record\_log\_min\_entropy\_clear},
\ec{fs\_with\_aborts\_failure},
\ec{transcript\_signature\_programming\_site\_clear},
\ec{transcript\_log\_signature\_programming\_sites\_clear},
\ec{fs\_with\_aborts\_bound\_obligation}; 80--99 &
Lifts bad-event predicates to transcript/signing logs and names the legacy
Fiat-Shamir-with-aborts loss-obligation nonnegativity check. &
Machine-checked predicate; theorem-boundary obligation. &
Depends on reprogramming and min-entropy predicates plus reduction loss terms. &
Used by HopGames transcript-log invariants; one clear-log predicate is also
visible in \file{HAETAE\_Security.ec}. \\
\midrule
\ec{counted\_rom\_event} and \ec{counted\_event\_is\_*}; 103--131 &
Defines counted events for hash queries, signature sites, programmed sites, and
min-entropy checks, with boolean classifiers. &
Machine-checked fact. &
Pure local sum-type analysis. &
Audit vocabulary for counted ROM traces and bad-flag explanations. \\
\midrule
\ec{ro\_query\_is\_challenge},
\ec{challenge\_query\_log\_count},
\ec{challenge\_query\_log\_count\_*}; 133--192 &
Classifies challenge-hash queries and proves how consing challenge or
nonchallenge queries changes the distinct challenge-query count. &
Machine-checked fact; query-budget invariant. &
Uses finite-set cardinality over filtered query logs. &
Discharges HopGames budget obligations for message-hash, matrix, sampler, and
challenge query steps. \\
\midrule
\ec{counted\_trace\_*}; 193--202 &
Counts the four counted event families by filtering a trace. &
Machine-checked fact. &
Uses list filtering and event classifiers. &
Provides trace-accounting vocabulary, though most downstream use goes through
higher-level bad-flag lemmas. \\
\midrule
\ec{programming\_sites\_same\_query},
\ec{programming\_sites\_conflict},
\ec{programming\_site\_prequeried},
\ec{programming\_site\_reprograms},
\ec{programming\_site\_fresh},
\ec{programming\_transcript\_fresh}; 205--234 &
Defines prequery, reprogramming conflict, freshness, and transcript freshness. &
Machine-checked predicate. &
Depends only on query equality, output inequality, query logs, programmed-site
logs, and min-entropy failure. &
Core bad-event guards for ROM programming games. \\
\midrule
\ec{rom\_counted\_programming\_bad},
\ec{rom\_counted\_state\_bad},
\ec{counted\_rom\_programming\_bound\_obligation}; 235--247 &
Summarizes min-entropy, prequery, and reprogramming bad states and names the
nonnegativity obligation for the counted loss term. &
Machine-checked predicate; theorem-boundary obligation. &
Uses \ec{counted\_rom\_programming\_loss\_term}. &
Connects local bad flags to the counted-ROM loss term later composed by
HopGames and Security. \\
\midrule
\ec{CountedLazyROM}; 250--301 &
Stateful wrapper over an oracle with logs for hash queries, programmed sites,
signature sites, and flags for prequery, reprogramming, and min-entropy bad
events. &
Formal ROM modeling fact; module/procedure boundary. &
Requires an \ec{Oracle}; uses \ec{O.get} and \ec{O.set}. &
Instantiated inside HopGames to expose counted bad events.  It is not a SHAKE
implementation theorem. \\
\midrule
\ec{SigningCoinSampler}, \ec{DirectSigningCoinSampler}; 302--313 &
Abstracts coin sampling and gives the direct sampler from \ec{drandom\_coins}. &
Machine-checked module interface. &
Uses \ec{drandom\_coins} and its imported distribution facts. &
Target of \ec{phoare} prequery bounds used in sampled signing games. \\
\midrule
\ec{CountedLazyROMStateFacts}; 315--500 &
Hoare invariants: initialization clears state, \ec{get} preserves clear flags
and nonchallenge budgets, \ec{observe\_transcript} preserves clear state when
min-entropy is clear, \ec{program} preserves clear state when fresh, and
\ec{bad} is false under clear flags. &
Machine-checked fact. &
Premises include freshness, nonchallenge-query classification, and
min-entropy-clear conditions. &
Supplies HopGames state-invariant steps for counted ROM hybrids. \\
\midrule
\ec{CountedROMInterface}, \ec{CountedROMClient},
\ec{CountedLazyROMBadGame}, \ec{ProgramChallenge},
\ec{TranscriptProgrammer}, \ec{TranscriptProgrammerFromSite}; 501--549 &
Defines the client interface, bad-game wrapper, direct challenge-programming
procedure, and transcript-to-site programmer. &
Formal ROM modeling fact; module/procedure boundary. &
Depends on \ec{CountedLazyROM} and the generic oracle interface. &
Gives the game API whose bad probability is bounded by later lemmas. \\
\midrule
\ec{programming\_site\_self\_*},
\ec{fs\_with\_aborts\_failureE},
\ec{programming\_site\_prequeried\_*},
\ec{honest\_signing\_programming\_site\_queryE}; 550--610 &
Unfolds self-matching, failure decomposition, prequery preservation, and the
honest signing challenge query. &
Machine-checked fact family. &
Uses definitions from transcript, ROM, and honest signing. &
Local rewriting library for later no-failure, query-budget, and prequery
bounds. \\
\midrule
\ec{programmed\_site\_entropy\_token\_from\_query},
\ec{honest\_signing\_programming\_site\_query\_entropy\_token},
\ec{honest\_signing\_programming\_site\_query\_token\_injective},
\ec{signing\_distribution\_programming\_site\_query\_spread}; 611--671 &
Extracts an entropy token from a challenge query and bounds the probability
that a signing distribution hits any fixed programming-site query. &
Machine-checked fact; conditional premise for arbitrary distributions. &
Needs \ec{signing\_randomness\_token\_spread} for arbitrary distributions and
the signing-coin token lemmas from distributions. &
Min-entropy engine for prequery loss. \\
\midrule
\ec{honest\_signing\_programming\_site\_query\_spread},
\ec{honest\_signing\_programming\_site\_query\_spread\_bound},
\ec{honest\_signing\_programming\_site\_query\_spread\_bound\_for},
\ec{honest\_signing\_programming\_site\_query\_is\_challenge}; 673--713 &
Instantiates the query-spread bound to \ec{drandom\_coins} and proves honest
programming queries are challenge queries. &
Machine-checked fact. &
Uses \ec{signing\_coin\_distribution\_token\_spread}. &
Turns entropy assumptions into concrete challenge-query budget bounds. \\
\midrule
\ec{programming\_site\_prequeried\_honest\_*},
\ec{honest\_signing\_programming\_site\_outputE},
\ec{honest\_signing\_programming\_site\_functional\_output},
\ec{honest\_signing\_programming\_sites\_no\_conflict},
\ec{actual\_signature\_programming\_site\_*}; 713--883 &
Shows nonchallenge message-hash consing does not create honest-site prequeries,
honest sites have functional outputs, honest sites do not conflict, and actual
signature sites equal honest signing sites. &
Machine-checked fact. &
Unfolds \ec{sign\_internal}, challenge query construction, and signature-field
definitions. &
Supports message-hash-aware prequery bounds and no-reprogramming arguments in
HopGames. \\
\midrule
\ec{programming\_site\_conflict\_free},
\ec{programming\_site\_log\_conflict\_free\_for\_honest},
\ec{programming\_site\_conflict\_free\_no\_reprograms},
\ec{actual\_signature\_programming\_site\_sign\_internal\_conflict\_step},
\ec{programmed\_site\_reprogram\_one\_step\_conflict\_free\_zero}; 884--968 &
Defines conflict-free programmed-site logs and proves honest sites remain
conflict-free, making one-step reprogramming probability zero under the
conflict-free premise. &
Machine-checked fact; conditional premise for the zero-probability lemma. &
Premise: all honest sites are conflict-free against the current log. &
Separates reprogramming loss from prequery loss in the counted-ROM argument. \\
\midrule
\ec{programmed\_site\_prequery\_one\_step\_*}; 969--1230 &
Bounds the probability that a sampled honest site is already in a query log,
using either finite-set cardinality, total budget, challenge-query budget, or
concrete counter hypotheses. &
Machine-checked probability bound; query-budget invariant. &
Premises include min-entropy bounds, cardinality bounds, and
\ec{hc <= hash\_query\_budget\_count}. &
Main prequery-loss family used by HopGames sampled-signing hybrids. \\
\midrule
\ec{programmed\_site\_prequery\_one\_step\_message\_hash\_counter\_bound},
\ec{direct\_signing\_coin\_sampler\_*}; 1231--1351 &
Lifts the prequery bound through a leading message-hash query and into
\ec{phoare} statements for the direct signing-coin sampler. &
Machine-checked probability bound. &
Premises: nonnegative and bounded challenge-query counter plus query-count
hypothesis. &
Used when HopGames samples signing coins after message-hash interactions. \\
\midrule
\ec{programmed\_site\_prequery\_one\_step\_concrete\_bound},
\ec{programmed\_site\_prequery\_reprogram\_one\_step\_concrete\_bound}; 1353--1386 &
Combines concrete prequery and conflict-free reprogramming reasoning into a
single one-step bad-event bound. &
Machine-checked probability bound; conditional premise. &
Requires a cardinality budget and conflict-free honest-site premise. &
Local one-step bound feeding counted-ROM bad-event accounting. \\
\midrule
\ec{programmed\_site\_query\_min\_entropy\_bound\_constant\_impossible},
\ec{programming\_site\_reprograms\_conflict\_cons},
\ec{rom\_counted\_programming\_bad\_*},
\ec{counted\_rom\_programming\_bound\_obligation\_holds}; 1404--1448 &
Records sanity checks: constant low-entropy query generation contradicts a
subunit bound; conflicting cons creates reprogramming; min-entropy/prequery
conditions imply the counted bad predicate; the counted loss term is
nonnegative. &
Machine-checked fact. &
Uses signing-coin losslessness and nonnegativity of reduction loss terms. &
Guards against vacuous min-entropy and bad-state accounting mistakes. \\
\midrule
\ec{CountedLazyROMBadFlags},
\ec{counted\_lazy\_rom\_bad\_flags\_budget\_sound},
\ec{counted\_lazy\_rom\_bad\_flags\_*}; 1455--1478 &
Packages the probability of \ec{CountedLazyROMBadGame} and states that, under
a budget-soundness premise, it is bounded by
\ec{counted\_rom\_programming\_loss\_term} and is subunit. &
Conditional premise and machine-checked probability wrapper. &
Premise: \ec{counted\_lazy\_rom\_bad\_flags\_budget\_sound} for the exact
game probability. &
Theorem boundary used by HopGames when charging bad events to the counted-ROM
loss term. \\
\midrule
\ec{fs\_with\_aborts\_no\_programming\_failure},
\ec{transcript\_signature\_challenge\_output\_matches},
\ec{honest\_transcript\_challenge\_entropy\_support},
\ec{honest\_transcript\_no\_min\_entropy\_failure},
\ec{signing\_transcript\_relation\_*},
\ec{signing\_record\_log\_sound\_*},
\ec{fs\_with\_aborts\_no\_failure\_*}; 1491--1624 &
Proves honest transcripts and sound signing records do not trigger
min-entropy or legacy Fiat-Shamir-with-aborts failure when programming matches. &
Machine-checked fact. &
Uses signing transcript relations, \ec{challenge\_from\_seed\_sparse} through
the challenge-sparsity path, and transcript validity. &
Used by HopGames to maintain clear transcript logs and no-failure side
conditions. \\
\midrule
\ec{transcript\_valid\_programming\_site\_*},
\ec{transcript\_valid\_actual\_signature\_programming\_site\_*},
\ec{valid\_forking\_pair\_signature\_programming\_site\_query},
\ec{valid\_forking\_pair\_programming\_site\_query},
\ec{valid\_forking\_pair\_self\_programming\_sites\_no\_reprogramming};
1639--1723 &
Valid transcripts make transcript and signature programming-site queries agree;
valid forks share programming-site queries; self-programmed valid sites do not
count as reprogramming failures. &
Machine-checked fact. &
Premises: \ec{transcript\_valid}, \ec{valid\_forking\_pair}, and challenge
matches. &
Bridge from transcript special-soundness facts to ROM programming sites in
HopGames. \\
\midrule
\ec{fs\_bound\_parts\_nonnegative},
\ec{fs\_with\_aborts\_bound\_obligation\_holds}; 1739--1745 &
Proves nonnegativity of the legacy Fiat-Shamir-with-aborts reprogramming and
min-entropy loss terms. &
Machine-checked fact; theorem-boundary obligation. &
Uses nonnegativity lemmas from \file{HAETAE\_Reductions.ec}. &
Allows downstream bounds to add these terms safely; it does not prove the old
paper-style rejection ledger is tight or nonvacuous. \\
\bottomrule
\end{longtable}

\section{Feedback 3 criticism for the foundational dependency layer}

The current report is still not enough for a mathematician or cryptologist to
understand the EasyCrypt proof scripts from the report plus scripts alone at
the layer below \file{HAETAE\_HopGames.ec} and \file{HAETAE\_Security.ec}.
The missing material is not cosmetic.  The hop and final-security files consume
definitions from the formal HAETAE model as if they were ordinary mathematical
objects: parameters, polynomial carriers, signature fields, random-oracle query
constructors, signing-attempt distributions, rejection predicates, and concrete
hash wrappers.  Without a companion explanation for those files, a reader can
see that a theorem compiles but cannot tell which part of the argument is an
algebraic fact, which part is a model definition, which part is a conditional
sampling premise, and which part is only explanatory idealization.

The main evidence from the scripts is as follows.

\begin{itemize}
\item \file{HAETAE\_Security.ec} proves correctness through
\ec{verify\_internal\_sign\_internal}, \ec{valid\_signature\_sign\_internal},
and \ec{verify\_norm\_ok\_current}.  Those facts are not security
assumptions; they are checked algebra/model facts from \file{HAETAE\_Algebra.ec}.
\item \file{HAETAE\_HopGames.ec} repeatedly samples from
\ec{drandom\_coins}, \ec{dsigning\_attempt\_state}, and
\ec{dexact\_hyperball\_signing\_attempt\_state}.  A reader must know that the
current ``exact hyperball'' distribution is a checked bounded-product carrier,
with comments in \file{HAETAE\_Distributions.ec} marking the paper-faithful
hyperball/sphere sampler as remaining correspondence work.
\item The rejection bridge uses conditional premises such as
\ec{structural\_to\_exact\_hyperball\_paper\_sample\_loss\_obligation} in
\file{HAETAE\_HopGames.ec} and obligation families in
\file{HAETAE\_Rejection.ec}.  Those are not proved by the top-level theorem;
they are explicitly consumed as premises or discharged only after accepting
the current checked carrier/loss boundary.
\item The scheme module in \file{HAETAE\_Scheme.ec} makes exactly three
random-oracle query kinds in signing: sampler expansion, message hash, and
challenge hash.  The ROM programming proof relies on this query discipline.
Without the \file{HAETAE\_ROM.ec} query/output dictionary, the bad-event and
query-budget lemmas are opaque.
\item \file{HAETAE\_Algebra.ec} imports \file{HAETAE\_FIPS202.ec}, which
imports \file{HAETAE\_Keccak1600.ec}.  These files provide checked byte-list
and size facts for the formal model, not a standard-model proof that SHAKE is
a random oracle.
\end{itemize}

This section closes that gap at the level needed to read the scripts.  It does
not claim that the report alone can replace line-by-line EasyCrypt source.
The checked source remains authoritative.

\begin{longtable}{p{0.24\linewidth}p{0.34\linewidth}p{0.32\linewidth}}
\toprule
\textbf{Category} & \textbf{What belongs here in this layer} & \textbf{What must not be inferred} \\
\midrule
Machine-checked facts &
Definitions, typing, losslessness, support, size, well-formedness, equality,
Hoare, equivalence, and probability inequalities proved inside the manifest
files. &
They do not automatically prove equality with the HAETAE paper specification,
with C/Jasmin code, or with a concrete security estimate. \\
\midrule
External cryptographic assumptions &
MLWE, Module-SIS, and bimodal self-target MSIS hardness surfaces from
\file{HAETAE\_Assumptions.ec} and real-valued hardness/loss terms from
\file{HAETAE\_Reductions.ec}. &
The foundational files below do not estimate these assumptions; they only
define the formal carriers and adapters used by the reductions. \\
\midrule
Conditional premises &
Sampling and hop premises such as structural-to-exact hyperball lifting and
random-oracle reprogramming bounds when they appear as antecedents of a
security theorem. &
They must not be presented as already discharged unless a checked lemma in the
same manifest proves the exact antecedent. \\
\midrule
Imported/library facts &
EasyCrypt distribution, list, finite-map, real-order, and random-oracle
library facts, plus the local Keccak/FIPS wrapper files. &
Imported facts are part of the verification environment; a standalone proof
archive must include the version and import contract. \\
\midrule
Explanatory prose &
Descriptions such as ``formal HAETAE model'', ``bounded carrier'', and
``paper-shaped sampler interface''. &
Such prose is not a theorem.  It guides reading and marks boundaries that the
checked scripts either prove, assume, or leave conditional. \\
\bottomrule
\end{longtable}

\section{Script companion coverage for HAETAE Params}

\file{HAETAE\_Params.ec} is the numeric parameter base of the formal model.
It has no cryptographic assumption and no probability theorem.  Its
machine-checked content is definitional arithmetic for modes and positivity
facts consumed by algebra, distributions, rejection sampling, and final
correctness.

\begin{longtable}{p{0.26\linewidth}p{0.36\linewidth}p{0.28\linewidth}}
\toprule
\textbf{Script object} & \textbf{Mathematical reading} & \textbf{Downstream role} \\
\midrule
\ec{seedbytes}, \ec{crhbytes}, \ec{n}, \ec{q}, \ec{dq} &
Global byte, polynomial-degree, and modulus constants.  The checked fact
\ec{dqE} rewrites \(dq\) to \(129026\). &
Used throughout the algebraic carriers, seed slicing, SHAKE wrappers, public
key packing, challenge size, and distribution lengths. \\
\midrule
\ec{mode = [ Mode2 | Mode3 | Mode5 ]} &
The formal model has three parameter modes.  There is no abstract
well-formedness condition on modes beyond membership in this finite type. &
\ec{HAETAE\_Scheme.ec} declares one abstract \ec{haetae\_mode}; all later
security statements are mode-parametric through this value. \\
\midrule
\ec{mode\_k}, \ec{mode\_l}, \ec{mode\_tau} &
Mode-dependent vector dimensions and sparse-challenge weight prefix. &
Used by matrix/vector dimensions, challenge sparsity, transcript extraction,
and Module-SIS/bimodal relation shapes. \\
\midrule
\ec{mode\_b0sq}, \ec{mode\_b1sq}, \ec{mode\_b2sq} &
Mode-dependent squared-norm bounds for secret, response, and verification
checks. &
Used by \ec{secret\_norm\_ok}, \ec{response\_norm\_ok}, and
\ec{verify\_norm\_ok} in \file{HAETAE\_Algebra.ec}, then by rejection-failure
events and correctness. \\
\midrule
\ec{mode\_ln}, \ec{mode\_alpha\_hint}, \ec{mode\_m} &
Formal signing-sample bound, hint-decomposition base, and \(m=l-1\). &
Used by bounded/hyperball carriers, decomposition lemmas, public-key
verification columns, and rejection-sampling thresholds. \\
\midrule
\ec{mode\_crypto\_bytes}, \ec{mode\_polyq\_packedbytes},
\ec{mode\_publickeybytes} &
Formal byte-length constants for signature and public-key encodings. &
Used to prove packing well-formedness and that the pack branch in rejection
sampling is never the abort source in the current model. \\
\midrule
Positivity and size lemmas &
\ec{mode\_k\_gt0}, \ec{mode\_l\_gt0}, \ec{mode\_tau\_le\_n},
\ec{mode\_crypto\_bytes\_gt0}, and related lemmas are checked by mode case
analysis. &
They discharge side conditions for \ec{dlist}, \ec{mkseq}, support predicates,
norm arithmetic, and real-probability nonnegativity proofs. \\
\bottomrule
\end{longtable}

Boundary: the parameter file checks internal arithmetic consistency of the
formal constants.  It does not prove that these constants are the claimed
HAETAE standard parameter sets or that they yield a concrete bit-security
level.

\subsection{Lemma-ledger companion for HAETAE Params}

\paragraph{Lemma-ledger criticism for HAETAE Params.}
The previous companion coverage identified the parameter families, but it did
not give enough detail for a reader to reconstruct which later EasyCrypt
obligations are discharged merely by unfolding numeric constants and which
claims are outside \file{HAETAE\_Params.ec}.  The gap matters because these
constants are used everywhere: vector dimensions, challenge weight, norm
bounds, byte lengths, and sampler supports all inherit their side conditions
from this file.  The script proves only finite-mode arithmetic facts about the
formal constants.  It does not prove correspondence with any external HAETAE
standard document, does not justify a concrete security level, and does not
introduce MLWE, Module-SIS, ROM, or sampler assumptions.  The ledgers below
make that boundary explicit.

\paragraph{Imported-interface and assumption ledger for HAETAE Params.}
{\footnotesize
\begin{longtable}{p{0.23\linewidth}p{0.16\linewidth}p{0.28\linewidth}p{0.23\linewidth}}
\toprule
\textbf{EasyCrypt object} & \textbf{Lines / category} &
\textbf{Mathematical statement paraphrase} &
\textbf{Local premises, imports, downstream role} \\
\midrule
\ec{AllCore} &
1. Imported/library fact. &
Provides the base logic and arithmetic environment used by the parameter
definitions and case-split proofs. &
No local premise.  All later imports of \file{HAETAE\_Params.ec} inherit this
base proof environment. \\
\midrule
Formal parameter table &
5--62. Machine-checked fact. &
The file fixes one formal table of values for \ec{Mode2}, \ec{Mode3}, and
\ec{Mode5}. &
The table is consumed by \file{HAETAE\_Algebra.ec},
\file{HAETAE\_Distributions.ec}, \file{HAETAE\_Rejection.ec}, and the hop and
security scripts through those imports. \\
\midrule
External HAETAE-standard correspondence &
5--62. External cryptographic assumption. &
The file does not prove that the numeric table matches a particular HAETAE
submission, standard, or concrete security claim. &
Any such correspondence must be justified outside this EasyCrypt file.  The
checked security theorem is relative to these formal constants. \\
\midrule
Cryptographic assumptions introduced locally &
1--113. External cryptographic assumption. &
None are introduced in \file{HAETAE\_Params.ec}. &
MLWE, Module-SIS, Fiat-Shamir, ROM, rejection-sampling, and implementation
equivalence assumptions enter only in later files. \\
\bottomrule
\end{longtable}
}

\paragraph{Definition ledger for HAETAE Params.}
{\footnotesize
\begin{longtable}{p{0.25\linewidth}p{0.14\linewidth}p{0.28\linewidth}p{0.23\linewidth}}
\toprule
\textbf{EasyCrypt object(s)} & \textbf{Lines / category} &
\textbf{Mathematical statement paraphrase} &
\textbf{Local premises, imports, downstream role} \\
\midrule
\ec{seedbytes}, \ec{crhbytes}, \ec{n}, \ec{q}, \ec{dq} &
5--9. Machine-checked fact. &
Fixes seed length, CRH byte length, polynomial degree, modulus, and doubled
modulus for the formal model. &
No local premise.  Used by algebraic carriers, SHAKE seed slices,
coefficient-reduction facts, byte packing, and seed/CRH distributions. \\
\midrule
\ec{mode} &
11. Machine-checked fact. &
Declares the three formal modes \ec{Mode2}, \ec{Mode3}, and \ec{Mode5}. &
No separate validity predicate exists.  \file{HAETAE\_Scheme.ec} later exposes
a single abstract \ec{haetae\_mode}; downstream theorems are relative to that
formal mode. \\
\midrule
\ec{mode\_k}, \ec{mode\_l}, \ec{mode\_tau} &
13--26. Machine-checked fact. &
Defines per-mode public-vector dimension, secret-vector dimension, and sparse
challenge prefix weight. &
Used by \file{HAETAE\_Algebra.ec} for vector/matrix and challenge shape,
\file{HAETAE\_Distributions.ec} for product-distribution lengths, and
\file{HAETAE\_Rejection.ec} for attempt-state dimensions. \\
\midrule
\ec{mode\_b0sq}, \ec{mode\_b1sq}, \ec{mode\_b2sq} &
28--41. Machine-checked fact. &
Defines per-mode squared-norm thresholds for secret, response, and verifier
norm checks. &
Used by \ec{secret\_norm\_ok}, \ec{response\_norm\_ok}, and
\ec{verify\_norm\_ok}; rejection and security scripts inherit these as the
formal abort/correctness thresholds. \\
\midrule
\ec{mode\_ln}, \ec{mode\_alpha\_hint}, \ec{mode\_m} &
43--50. Machine-checked fact. &
Defines the signing-sample bound, hint decomposition base, and \(m=l-1\). &
Used by bounded and hyperball sample carriers, decomposition lemmas, public
verification columns, and rejection-sampling loss interfaces. \\
\midrule
\ec{mode\_crypto\_bytes}, \ec{mode\_polyq\_packedbytes},
\ec{mode\_publickeybytes} &
52--62. Machine-checked fact. &
Defines formal signature/public-key byte lengths and per-polynomial
\(q\)-packing lengths. &
Used by public-key packing roundtrips in \file{HAETAE\_Algebra.ec} and by
rejection packing-abort facts in \file{HAETAE\_Rejection.ec}. \\
\bottomrule
\end{longtable}
}

\paragraph{Lemma ledger for HAETAE Params.}
{\footnotesize
\begin{longtable}{p{0.25\linewidth}p{0.14\linewidth}p{0.28\linewidth}p{0.23\linewidth}}
\toprule
\textbf{EasyCrypt lemma(s)} & \textbf{Lines / category} &
\textbf{Mathematical statement paraphrase} &
\textbf{Local premises, imports, downstream role} \\
\midrule
\ec{seedbytes\_gt0}, \ec{crhbytes\_gt0}, \ec{n\_gt0},
\ec{q\_gt0}, \ec{dqE} &
65--69. Machine-checked fact. &
Proves positivity for global sizes/modulus and rewrites \ec{dq} to
\(129026\). &
No local hypotheses.  Discharges nonempty list, distribution, modulus, and
division side conditions in \file{HAETAE\_Algebra.ec} and
\file{HAETAE\_Distributions.ec}. \\
\midrule
\ec{mode\_k\_gt0}, \ec{mode\_l\_gt0}, \ec{mode\_tau\_gt0},
\ec{mode\_tau\_le\_n} &
71--80. Machine-checked fact. &
Every formal mode has positive vector dimensions and challenge weight bounded
by \(n\). &
Premise is only \ec{md : mode}; proofs are by case analysis.  Used by shape
lemmas, sparse challenge support, distribution support, and HopGames
dimension side conditions. \\
\midrule
\ec{mode\_b0sq\_gt0}, \ec{mode\_b1sq\_gt0},
\ec{mode\_b2sq\_gt0}, \ec{mode\_ln\_gt0},
\ec{mode\_alpha\_hint\_gt0}, \ec{mode\_m\_gt0} &
83--98. Machine-checked fact. &
Proves positivity of formal norm thresholds, signing-sample bound,
hint-decomposition base, and \(m\). &
Premise is \ec{md : mode}.  Used by norm proofs, bounded-sampler support,
decomposition arithmetic, keygen seed counters, and rejection threshold
lemmas. \\
\midrule
\ec{mode\_crypto\_bytes\_gt0},
\ec{mode\_polyq\_packedbytes\_gt0},
\ec{n\_le\_mode\_polyq\_packedbytes},
\ec{mode\_publickeybytes\_gt0} &
101--110. Machine-checked fact. &
Proves positivity of byte-size interfaces and that a formal \(q\)-packed
polynomial has at least \(n\) coefficient positions. &
Premise is \ec{md : mode}.  Used by public-key packing/unpacking,
verification-matrix reconstruction, and rejection packing-abort proofs. \\
\bottomrule
\end{longtable}
}

\section{Script companion coverage for HAETAE Algebra}

\file{HAETAE\_Algebra.ec} is the formal HAETAE model core.  It defines the
carrier types and deterministic operations that the security games manipulate.
The file is machine-checked, but its definitions are part of the model rather
than a proof that the model is identical to the paper algorithm or an
implementation.

\begin{longtable}{p{0.25\linewidth}p{0.39\linewidth}p{0.26\linewidth}}
\toprule
\textbf{Script object family} & \textbf{Mathematical reading} & \textbf{Downstream role} \\
\midrule
Carrier types &
\ec{byte}, \ec{seed}, \ec{crh}, \ec{random\_coins}, \ec{message},
\ec{context}, \ec{coeff}, \ec{poly}, \ec{challenge}, \ec{polyveck},
\ec{polyvecl}, \ec{matrix}, \ec{pkey}, \ec{skey}, and \ec{signature} are
list- and tuple-based carriers. &
These are the concrete type instantiations used when
\file{HAETAE\_Scheme.ec} clones \file{Sig\_ROM.ec}.  All games and reductions
share these carriers. \\
\midrule
Zero values and well-formedness &
\ec{poly\_zero}, vector zeroes, \ec{matrix\_zero},
\ec{poly\_wf}, \ec{polyveck\_wf}, \ec{polyvecl\_wf},
\ec{crh\_wf}, \ec{challenge\_wf}, and unit/sparse predicates define local
shape constraints. &
Used by signature validity, transcript validity, distribution support, and
sampling support lemmas. \\
\midrule
Polynomial and vector operations &
\ec{coeff\_mod}, \ec{poly\_add}, \ec{poly\_mul}, \ec{poly\_dot},
\ec{matrix\_vec\_mul}, \ec{polyveck\_add}, \ec{polyveck\_sub}, and
\ec{polyvecl\_sub} define the algebraic operations over the formal list model. &
The checked \ec{\_wf} and coefficient-bound lemmas make later Module-SIS and
verification equations well typed and well shaped. \\
\midrule
Decomposition and public-key packing families &
\ec{coeff\_decompose\_z1\_*}, \ec{poly\_highbits},
\ec{poly\_lowbits}, \ec{polyveck\_vk\_highbits},
\ec{polyveck\_highbits\_hint}, packing/unpacking operations, and roundtrip
lemmas describe the formal byte/list encoding layer. &
Used by public-key reconstruction, verification-matrix construction, and
transcript public-equation lemmas.  These are checked formal encodings, not a
complete C/Jasmin parsing theorem. \\
\midrule
FIPS/SHAKE seed flow &
\ec{haetae\_shake256\_bytes}, \ec{haetae\_xof256\_*},
\ec{haetae\_keygen\_rhoprime}, \ec{haetae\_keygen\_sigma}, and
\ec{haetae\_keygen\_key} wrap \file{HAETAE\_FIPS202.ec}. &
Used to define formal key-generation seeds and public matrices.  This is
checked byte-list plumbing, not a random-oracle theorem about SHAKE. \\
\midrule
Key-generation relation &
\ec{public\_matrix\_seed}, \ec{public\_secret\_vector\_seed},
\ec{public\_error\_vector\_seed}, \ec{public\_key\_vector\_seed},
\ec{public\_key\_relation}, \ec{public\_key\_equation\_holds}, and
\ec{public\_key\_of\_secret} define the formal public key relation. &
Used by \file{HAETAE\_Transcript.ec} extraction, by Module-SIS instances, and
by \file{HAETAE\_Security.ec} correctness. \\
\midrule
Signature fields &
\ec{valid\_signature}, \ec{sig\_commitment\_highbits},
\ec{sig\_commitment\_lowbits}, \ec{sig\_message\_hash},
\ec{sig\_challenge}, \ec{sig\_response}, \ec{sig\_response\_aux}, and
\ec{sig\_hint} define the tuple layout of a signature. &
Used by transcript construction, challenge-query predicates, rejection
attempt states, and verification. \\
\midrule
Deterministic signing model &
\ec{commitment\_raw}, \ec{commitment\_lowbits}, \ec{message\_hash},
\ec{challenge\_hash}, \ec{commitment\_challenge},
\ec{commitment\_highbits}, \ec{response\_vector}, \ec{hint\_vector}, and
\ec{signature\_token} define signing fields from secret key, message, context,
and coins. &
Used by the real, public-simulation, paper-simulation, and rejection-simulation
signing games in \file{HAETAE\_HopGames.ec}. \\
\midrule
Internal algorithms &
\ec{keygen\_internal}, \ec{sign\_internal},
\ec{verify\_challenge\_consistent}, \ec{verify\_public\_equation},
\ec{verify\_norm\_ok}, and \ec{verify\_internal} define the scheme core. &
\ec{HAETAE\_Security.ec} correctness reduces directly to
\ec{verify\_internal\_sign\_internal}.  \file{HAETAE\_Scheme.ec} exposes these
operations through the generic signature interface. \\
\bottomrule
\end{longtable}

The main checked lemma families are:

\begin{longtable}{p{0.30\linewidth}p{0.42\linewidth}p{0.18\linewidth}}
\toprule
\textbf{Lemma family} & \textbf{Checked fact} & \textbf{Boundary} \\
\midrule
Well-formedness and coefficient bounds &
\ec{poly\_add\_wf}, \ec{poly\_mul\_wf},
\ec{matrix\_vec\_mul\_wf}, vector well-formedness lemmas, coefficient-bound
lemmas such as \ec{polyveck\_add\_coeffs\_q\_bound}, decomposition
\ec{\_wf} lemmas, and poly-\(q\) bounds show that formal operations preserve
shape and ranges. &
Machine-checked. \\
\midrule
Seed/XOF and packing facts &
\ec{haetae\_reference\_keygen\_xof\_*},
\ec{haetae\_keygen\_*\_size}, \ec{haetae\_pack\_vk\_ref\_roundtrip\_*},
\ec{haetae\_public\_key\_roundtrip}, and public-key byte well-formedness
lemmas prove list-size and roundtrip facts for the model. &
Machine-checked formal-model facts; implementation equivalence remains outside
this file. \\
\midrule
Public-equation facts &
\ec{public\_key\_of\_secret\_relation},
\ec{public\_key\_of\_secret\_equation\_holds},
\ec{public\_key\_vector\_seed\_mode23\_keygen\_equation}, and
\ec{public\_key\_vector\_seed\_mode5\_keygen\_equation} prove the formal
verification equation shape. &
Machine-checked; used by transcripts and extraction. \\
\midrule
Signature validity and correctness &
\ec{valid\_signature\_sign\_internal},
\ec{verify\_challenge\_consistent\_sign\_internal},
\ec{response\_norm\_ok\_current}, \ec{verify\_norm\_ok\_current}, and
\ec{verify\_internal\_sign\_internal}. &
Machine-checked formal correctness, used by \file{HAETAE\_Security.ec}. \\
\bottomrule
\end{longtable}

Boundary: this file contains no MLWE/SIS hardness theorem and no ROM theorem.
It is the formal model that later security theorems prove statements about.

\subsection{Line-by-line proof-script companion for HAETAE Algebra}

Focused criticism for this pass: the preceding dependency-level table names
the major object families, but it is not enough for a mathematician or
cryptologist to reconstruct the proof obligations in
\file{HAETAE\_Algebra.ec}.  In particular, it did not say where the script
switches from model declarations to checked preservation lemmas, which
preconditions are theorem hypotheses rather than global facts, which proof
steps rely on EasyCrypt library lemmas such as \ec{size\_mkseq},
\ec{nth\_mkseq}, \ec{eq\_from\_nth}, and integer SMT, or why the final
\ec{verify\_internal\_sign\_internal} theorem is only a formal-model
correctness theorem.  The following line-range companion closes that gap at
proof-obligation granularity.  It is still not a replacement for reading the
source: tactic sequencing and all generated subgoals remain in the
\file{.ec} file and in EasyCrypt's checker.

\begin{longtable}{p{0.13\linewidth}p{0.18\linewidth}p{0.21\linewidth}p{0.38\linewidth}}
\toprule
\textbf{Lines} & \textbf{Script content} & \textbf{Classification} &
\textbf{Companion reading and downstream role} \\
\midrule
1--7 &
Imports and theory wrapper. &
Imported/library facts. &
\ec{AllCore}, \ec{IntDiv}, and \ec{List} provide boolean logic, integer
division/modulus, list constructors, \ec{size}, \ec{nth}, \ec{nseq},
\ec{mkseq}, \ec{range}, \ec{zip}, and \ec{foldl}.  \file{HAETAE\_Params.ec}
supplies all mode constants; \file{HAETAE\_FIPS202.ec} supplies the formal
SHAKE byte-list API.  Any proof later using \ec{smt}, \ec{size\_mkseq},
or SHAKE size lemmas imports those facts rather than reproving them here. \\
\midrule
9--29 &
Carrier declarations for bytes, seeds, polynomials, vectors, matrices, keys,
and signatures. &
Machine-checked definitions. &
These lines fix the formal universe: all algebra is over lists and tuples.
They are checked type declarations, not a theorem that these carriers equal
the HAETAE paper objects or Jasmin/C memory layouts.  \file{HAETAE\_Scheme.ec}
and the generic \file{Sig\_ROM.ec} clone use these exact carriers. \\
\midrule
31--106 &
Zero objects, well-formedness predicates, challenge sparsity predicates, and
their first size/well-formedness lemmas. &
Machine-checked facts with imported list facts. &
The proof pattern is definitional unfolding followed by \ec{size\_nseq},
\ec{all\_nseq}, \ec{allP}, and \ec{mem\_nth}.  These lemmas are the source of
later side conditions saying that a formal polynomial has length \(n\), a
\(k\)-vector has \ec{mode\_k md} rows, and a challenge has coefficients in
\(\{-1,0,1\}\). \\
\midrule
108--169 &
Coefficient arithmetic, polynomial operations, vector operations, matrix-vector
multiplication, and the list model of negacyclic multiplication. &
Machine-checked definitions; explanatory model boundary. &
The operations are definitions in the EasyCrypt model.  The file checks shape
and range facts about them later, but this block is not a mathematical proof
that the list operations are isomorphic to a quotient-ring implementation.
That correspondence would require a separate algebra or implementation
equivalence theorem. \\
\midrule
170--289 &
\ec{\_wf} preservation lemmas for arithmetic and matrix/vector operations. &
Machine-checked facts relying on imported list facts. &
The script repeatedly unfolds an operation, splits zero-special cases, proves
the length of \ec{mkseq} outputs, and uses induction for \ec{foldl} in
\ec{poly\_dot\_wf}.  These proofs discharge the routine obligations later
needed by \ec{valid\_signature}, transcript predicates, and verification
matrix construction. \\
\midrule
291--425 &
\ec{coeff\_q\_bound} and coefficient-range preservation lemmas. &
Machine-checked facts; imported integer arithmetic. &
The key obligation is that every arithmetic constructor returns coefficients
modulo \(q\).  EasyCrypt's integer division/modulus library and SMT prove
\ec{coeff\_mod\_q\_bound}; fold induction lifts that fact to dot products.
These range facts feed public-key packing and mode-dependent coefficient
packing bounds. \\
\midrule
427--562 &
High/low-bit decomposition, hint high bits, verification-key high bits, and
well-formedness of decomposition maps. &
Machine-checked definitions and facts. &
This block formalizes the bit-decomposition vocabulary used by public-key
packing and verification reconstruction.  The proofs check only list shape;
numeric range refinements appear in the next packing block.  Do not read this
as a standalone theorem about the paper decomposition unless the definition is
separately matched to the paper. \\
\midrule
564--672 &
Seed-derived public-key polynomials, SHAKE/XOF wrappers, seed-buffer layout,
and keygen seed slices. &
Machine-checked definitions plus imported hash facts. &
The script defines formal byte-list plumbing for the reference key-generation
seed flow.  Later lemmas use \file{HAETAE\_FIPS202.ec} size/domain facts.
This is not a ROM or SHAKE security assumption; it is deterministic byte-list
modeling. \\
\midrule
674--755 &
Formal keygen relation: matrix seed, secret/error vectors, rounding vector,
mode-specific public-key vector, and \ec{public\_key\_relation}. &
Machine-checked definitions. &
These are the definitions that extraction and verification use as the public
equation.  The theorem boundary is definitional: the relation is checked
inside the model, while cryptographic hardness of finding alternate witnesses
is handled outside this file by later reduction assumptions. \\
\midrule
757--841 &
Reference keygen aliases and seed-flow predicate. &
Machine-checked definitions and explanatory prose boundary. &
The reference names make the script readable against HAETAE naming conventions.
The file proves these aliases line up with the formal relation, but does not
certify a C or Jasmin reference implementation. \\
\midrule
841--902 &
Verification column/matrix construction, challenge embedding, public-key
challenge term, and reconstructed highbits. &
Machine-checked definitions. &
These lines define what \ec{verify\_public\_equation} later checks.  They are
used downstream by transcript public-equation predicates and rejection-attempt
field consistency. \\
\midrule
904--1265 &
Public-key byte encodings, simple pack/unpack definitions, reference
15-byte/16-bit packers, concrete packed-key aliases, and roundtrip predicates. &
Machine-checked definitions; explanatory implementation boundary. &
There are two encoding surfaces: a simple model packer and a reference-style
bit packer.  They are checked against their own definitions later.  The
presence of reference-style names should not be overread as a checked parser
equivalence to every external implementation. \\
\midrule
1267--1466 &
Seed/XOF and seed-buffer size/layout lemmas. &
Machine-checked facts with imported hash/list facts and conditional premises. &
Representative hypotheses include \ec{0 <= outlen}, \ec{0 <= len}, and
\ec{size sd = seedbytes}.  The proofs use unfolding, SHAKE size lemmas, and
offset/slice arithmetic.  The theorem boundary is deterministic size/layout
correctness, not pseudorandomness. \\
\midrule
1468--1767 &
Public matrix/vector well-formedness and poly-\(q\) range lemmas for formal and
reference keygen. &
Machine-checked facts with mode-case proof obligations. &
The script proves that generated rows and vectors have the right lengths and
coefficient bounds.  Mode-specific cases \ec{Mode2}, \ec{Mode3}, and
\ec{Mode5} are split explicitly where packing ranges differ. \\
\midrule
1769--1936 &
Verification-matrix well-formedness, seed pack/unpack, and byte-normalization
lemmas. &
Machine-checked facts relying on imported list/integer facts. &
The important proof idioms are \ec{eq\_from\_nth} for list equality,
\ec{nth\_cat} for concatenated encodings, and integer SMT for byte
normalization.  These lemmas make the public-key roundtrip theorems possible. \\
\midrule
1938--2753 &
Reference pack/unpack arithmetic for Mode2/3 15-byte blocks and Mode5 16-bit
blocks; vector-body and public-key reference roundtrips. &
Machine-checked facts with conditional premises. &
The conditions \ec{haetae\_polyq\_coeffs\_wf},
\ec{haetae\_polyveck\_polyq\_coeffs\_wf}, \ec{md = Mode2 \/ md = Mode3},
and \ec{size sd = seedbytes} are local theorem hypotheses, not global truths.
The proof decomposes byte offsets, proves each coefficient offset \(0,\ldots,7\),
then lifts coefficient equality to polynomial and vector equality using
\ec{eq\_from\_nth}. \\
\midrule
2771--3019 &
Simple model packing roundtrip, encoding well-formedness, and reference-key
well-formedness. &
Machine-checked facts with explicit hypotheses. &
The simple packer proves a model roundtrip for \ec{polyveck\_wf} inputs; the
reference packer proves a stronger poly-\(q\)-bounded roundtrip.  These facts
support keygen public-key recovery and verification-matrix construction in
later games. \\
\midrule
3022--3037 &
\ec{public\_key\_of\_secret\_relation} and
\ec{public\_key\_of\_secret\_equation\_holds}. &
Machine-checked facts. &
These are definitional correctness lemmas for honest keys.  They feed
transcript validity and extraction predicates, but they do not by themselves
prove a hardness statement about producing another relation witness. \\
\midrule
3039--3168 &
\ec{valid\_mode}, \ec{valid\_keypair}, \ec{valid\_signature}, signature
projection functions, encoders, deterministic signing-field generators, and
the signing entropy token extractor. &
Machine-checked definitions; explanatory model boundary. &
The formal signer is deterministic once coins are supplied.  The lowbits field
programs the first coefficient with the entropy token.  This modeling choice
is essential for ROM prequery and rejection-sampling arguments; it is not a
claim about concrete entropy extraction unless separately connected to the
implementation. \\
\midrule
3171--3371 &
Well-formedness/unit/sparsity lemmas for deterministic signing fields and
challenge hashes. &
Machine-checked facts. &
The proof pattern is constructor unfolding plus \ec{mkseqP}, \ec{allP}, and
case analysis on parity and \ec{i < mode\_tau md}.  These lemmas feed
\ec{valid\_signature\_sign\_internal}, transcript validity, and rejection
attempt acceptance. \\
\midrule
3373--3404 &
\ec{keygen\_internal}, \ec{sign\_internal},
\ec{verify\_challenge\_consistent}, \ec{verify\_reconstructed\_highbits},
and \ec{verify\_public\_equation}. &
Machine-checked definitions. &
This is the formal scheme core used by \file{HAETAE\_Scheme.ec}.  No
probability, ROM, or adversary reasoning occurs here; those enter in
\file{Sig\_ROM.ec}, \file{HAETAE\_HopGames.ec}, and
\file{HAETAE\_Security.ec}. \\
\midrule
3407--3767 &
Keygen validity, public-key equations, reference keygen equations, packed-key
roundtrips, verification-matrix correctness, and
\ec{valid\_signature\_sign\_internal}. &
Machine-checked facts with conditional premises. &
Mode-specific equations split \ec{Mode2}/\ec{Mode3} from \ec{Mode5};
roundtrip lemmas require \ec{size sd = seedbytes} or encoding
well-formedness.  The downstream role is to justify that honest formal keys
and signatures satisfy the predicates consumed by transcript and security
proofs. \\
\midrule
3777--3883 &
Norm definitions, current-signing norm lemmas, and
\ec{verify\_internal\_sign\_internal}. &
Machine-checked facts; theorem boundary. &
The final theorem proves: for this formal model, verifying an honestly signed
message under \ec{public\_key\_of\_secret md sk} succeeds.  The norm proofs
use unit-vector lemmas and mode case analysis.  This is a correctness theorem
for \ec{sign\_internal}/\ec{verify\_internal}, not HAETAE EUF-CMA security and
not a statement about an external implementation. \\
\bottomrule
\end{longtable}

The practical proof-script reading rule for this file is therefore: treat every
\ec{op} and \ec{type} as a checked model declaration, every \ec{lemma} as a
checked theorem only under its displayed hypotheses, every use of
\ec{smt}/list rewriting as imported EasyCrypt reasoning, and every connection
to the paper algorithm, SHAKE security, or implementation behavior as
explanatory unless another source file proves it.

\paragraph{Lemma-ledger criticism for HAETAE Algebra.}
The line-range companion above makes the file navigable, but it is still not
enough, by itself, for a mathematician or cryptologist to reconstruct the
proof obligations that later scripts cite.  It gives ranges, but not a stable
ledger of theorem surfaces: which names are definitions, which lemmas are
checked consequences of those definitions, which side conditions are local
premises, and which facts enter only because \file{AllCore}, \file{IntDiv},
\file{List}, \file{HAETAE\_Params.ec}, or \file{HAETAE\_FIPS202.ec} were
imported.  It also risks hiding the strongest boundary in the algebra file:
\file{HAETAE\_Algebra.ec} proves correctness of the formal
\ec{sign\_internal}/\ec{verify\_internal} model, not equivalence to the
HAETAE paper sampler, SHAKE-as-ROM security, or a Jasmin/C implementation.
The ledgers below close that documentation gap at the level of names later
used by \file{HAETAE\_Rejection.ec}, \file{HAETAE\_HopGames.ec}, and
\file{HAETAE\_Security.ec}.  Repeated coefficient-offset lemmas and sibling
\ec{\_wf} lemmas are grouped when they have the same proof obligation and the
same downstream role; the EasyCrypt source remains authoritative for exact
tactic order and generated subgoals.

\paragraph{Imported-interface ledger for HAETAE Algebra.}
{\footnotesize
\begin{longtable}{p{0.23\linewidth}p{0.16\linewidth}p{0.28\linewidth}p{0.23\linewidth}}
\toprule
\textbf{EasyCrypt interface} & \textbf{Lines / category} &
\textbf{Statement paraphrase} & \textbf{Premises, imports, downstream role} \\
\midrule
\ec{AllCore}, \ec{IntDiv}, \ec{List} &
1. Imported/library fact. &
Boolean reasoning, integer division/modulus, lists, \ec{size}, \ec{nth},
\ec{nseq}, \ec{mkseq}, \ec{range}, \ec{zip}, \ec{foldl}, and equality from
indexed lists. &
No local cryptographic premise.  Used in almost every algebra proof; later
Rejection and HopGames inherit these list and arithmetic facts through this
model. \\
\midrule
\file{HAETAE\_Params.ec} and \ec{HAETAE\_Params} &
2, 6. Imported/library fact. &
Provides the formal mode constants, dimensions, modulus, norm bounds, packing
lengths, positivity lemmas, and mode-case equations. &
Used by vector dimensions, sparse challenges, norm bounds, packing sizes, and
the rejection thresholds consumed in \file{HAETAE\_Rejection.ec}.  This is an
imported formal-parameter interface, not a proof that the constants match a
standard document. \\
\midrule
\file{HAETAE\_FIPS202.ec} and \ec{HAETAE\_FIPS202} &
2, 7. Imported/library fact. &
Provides byte-list SHAKE/XOF functions and size/domain lemmas such as
\ec{shake256\_size}, \ec{shake256\_squeeze\_size}, rate bytes, and domain
separator facts. &
Feeds keygen seed-flow lemmas.  It supplies deterministic byte-list facts, not
a random-oracle theorem and not SHAKE pseudorandomness. \\
\midrule
Cryptographic assumptions in this file &
1--3885. External cryptographic assumption. &
None are introduced by \file{HAETAE\_Algebra.ec}. &
MLWE, Module-SIS, ROM, Fiat-Shamir, rejection-sampling, paper-sampler, and
implementation-equivalence assumptions are outside this file and enter only in
later security/reduction layers. \\
\bottomrule
\end{longtable}
}

\paragraph{Definition ledger for HAETAE Algebra.}
{\footnotesize
\begin{longtable}{p{0.24\linewidth}p{0.15\linewidth}p{0.28\linewidth}p{0.23\linewidth}}
\toprule
\textbf{EasyCrypt object(s)} & \textbf{Lines / category} &
\textbf{Mathematical statement paraphrase} &
\textbf{Local premises, imports, downstream role} \\
\midrule
\ec{byte}, \ec{seed}, \ec{crh}, \ec{random\_coins}, \ec{xof256\_state},
\ec{message}, \ec{context}, \ec{coeff}, \ec{poly}, \ec{challenge},
\ec{polyveck}, \ec{polyvecl}, \ec{polyvecm}, \ec{matrix}, \ec{hint\_t},
\ec{sig\_token}, \ec{pkey}, \ec{encoded\_pkey}, \ec{skey}, \ec{signature} &
9--29. Machine-checked fact. &
Fixes the formal carrier universe as integers, lists, and tuples. &
No local premises.  Consumed by \file{HAETAE\_Scheme.ec} to instantiate
\file{Sig\_ROM.ec}; all Rejection, HopGames, and Security statements quantify
over these carriers. \\
\midrule
\ec{secret\_key\_of\_seed}, zero vectors/matrices, \ec{poly\_wf},
\ec{polyveck\_wf}, \ec{polyvecl\_wf}, \ec{crh\_wf},
\ec{challenge\_wf}, \ec{challenge\_sparse}, \ec{poly\_unit},
\ec{polyveck\_unit}, \ec{polyvecl\_unit} &
31--87. Machine-checked fact. &
Defines zero/default objects and shape, sparse-challenge, and unit-vector
predicates. &
Imports mode dimensions from \file{HAETAE\_Params.ec}.  These predicates are
the side conditions used by signature validity, distributions, rejection
samples, and transcript validity. \\
\midrule
\ec{coeff\_mod}, \ec{coeff\_add}, \ec{coeff\_neg}, \ec{coeff\_sub},
\ec{coeff\_mul}, \ec{coeff\_double}, \ec{poly\_coeff},
\ec{poly\_normalize}, \ec{poly\_add}, \ec{poly\_neg}, \ec{poly\_sub},
\ec{poly\_double}, \ec{poly\_mul}, \ec{poly\_dot},
\ec{matrix\_vec\_mul}, \ec{polyveck\_add}, \ec{polyveck\_sub},
\ec{polyveck\_double}, \ec{polyvecl\_sub} &
108--169. Machine-checked fact. &
Defines the list-based formal polynomial, vector, and matrix arithmetic used
by the HAETAE model. &
Imports \(q\), \(n\), and mode dimensions.  Used by public-key equations,
verification reconstruction, rejection reject2 balance, and extraction
relations.  This is not a separate quotient-ring implementation theorem. \\
\midrule
\ec{coeff\_q\_bound}, \ec{poly\_coeffs\_q\_bound},
\ec{polyveck\_coeffs\_q\_bound} &
291--297. Machine-checked fact. &
Names coefficient-range predicates for formal \(q\)-bounded values. &
No cryptographic premise.  Feeds packing poly-\(q\) bounds and public-key
vector well-formedness. \\
\midrule
\ec{coeff\_decompose\_z1\_low}, \ec{coeff\_decompose\_z1\_high},
\ec{coeff\_lsb}, \ec{coeff\_decompose\_vk\_low},
\ec{coeff\_decompose\_vk\_high}, \ec{coeff\_decompose\_hint\_high},
\ec{poly\_highbits}, \ec{poly\_lowbits}, \ec{poly\_vk\_lowbits},
\ec{poly\_vk\_highbits}, \ec{poly\_hint\_highbits}, vector lifts &
427--481. Machine-checked fact. &
Defines formal high/low-bit, verification-key, and hint decomposition
functions. &
Imports mode-specific hint bases.  Used by public-key packing, public-key
verification columns, signer commitments, and Rejection's public-equation
predicates.  Paper correspondence is explanatory unless separately proved. \\
\midrule
\ec{public\_key\_seed\_poly}, \ec{haetae\_shake256\_bytes},
\ec{haetae\_xof256\_*}, \ec{haetae\_seed\_slice},
\ec{haetae\_reference\_keygen\_*}, \ec{haetae\_keygen\_*},
\ec{haetae\_keygen\_sampler\_seed} &
564--678. Machine-checked fact. &
Defines deterministic seed expansion, XOF byte-list wrappers, seed-buffer
slices, and keygen seed components. &
Imports FIPS202 byte-list functions and parameter lengths.  Used by keygen and
packed-public-key lemmas.  This is deterministic plumbing, not SHAKE security
or ROM modeling. \\
\midrule
\ec{public\_matrix\_seed}, \ec{public\_secret\_vector\_seed},
\ec{public\_error\_vector\_seed}, \ec{public\_rounding\_vector\_seed},
\ec{haetae\_public\_key\_vector\_from\_relation},
\ec{public\_key\_relation}, \ec{public\_key\_equation\_holds},
reference-keygen aliases &
680--840. Machine-checked fact. &
Defines the formal public-key relation and the reference-named keygen views. &
Imports formal arithmetic and mode-specific branches.  Used by transcript
extraction, \file{HAETAE\_Security.ec} correctness, and Module-SIS relation
setup.  Hardness is not proved here. \\
\midrule
\ec{haetae\_mode23\_verification\_column},
\ec{haetae\_mode5\_verification\_column},
\ec{public\_key\_verification\_column},
\ec{haetae\_verification\_matrix\_from\_unpacked},
\ec{public\_verification\_matrix}, \ec{challenge\_embedding},
\ec{public\_key\_challenge\_term}, \ec{reconstructed\_highbits},
\ec{public\_key\_of\_secret} &
841--903. Machine-checked fact. &
Defines how public keys induce verification matrices, challenge terms, and
reconstructed highbits. &
Used by \ec{verify\_public\_equation}, Rejection field consistency,
HopGames paper/public simulation rewrites, and final correctness. \\
\midrule
\ec{haetae\_public\_key\_*}, \ec{haetae\_pack\_*},
\ec{haetae\_unpack\_*}, \ec{haetae\_byte\_norm},
\ec{haetae\_polyq\_*}, \ec{concrete\_haetae\_*} &
904--1265. Machine-checked fact. &
Defines the simple and reference-style public-key byte encoders/decoders and
their well-formedness predicates. &
Imports packing lengths and integer division.  Used by verification-matrix
correctness and key recovery lemmas.  Names containing ``concrete'' are still
formal EasyCrypt objects, not external implementation equivalence. \\
\midrule
\ec{valid\_mode}, \ec{valid\_keypair}, \ec{valid\_signature},
\ec{sig\_*} projections, \ec{encode\_poly}, \ec{encode\_polyveck},
\ec{encode\_pkey} &
3039--3062. Machine-checked fact. &
Defines formal validity and the exact tuple fields of signatures and encoded
public keys. &
Consumed by \file{HAETAE\_Transcript.ec}, Rejection attempt relations, HopGames
transcript games, and \file{HAETAE\_Security.ec}. \\
\midrule
\ec{deterministic\_*}, \ec{commitment\_source}, \ec{commitment\_raw},
\ec{response\_source}, \ec{response\_aux\_vector},
\ec{signing\_sample\_y1}, \ec{signing\_sample\_y2},
\ec{signing\_entropy\_token\_of\_coins}, \ec{commitment\_lowbits},
\ec{message\_hash}, \ec{challenge\_source}, \ec{challenge\_from\_seed},
\ec{challenge\_hash}, \ec{commitment\_challenge},
\ec{commitment\_highbits}, \ec{response\_vector}, \ec{hint\_vector},
\ec{signature\_token} &
3064--3169. Machine-checked fact. &
Defines the deterministic formal signing-field generator from a key, message,
context, and coins.  The first lowbits coefficient stores a coin-derived token. &
Used by Rejection attempt constructors, HopGames public/paper-simulation
rewrites, counted-ROM prequery arguments, and transcript soundness.  It is a
formal model of the signer, not a sampler correspondence proof. \\
\midrule
\ec{keygen\_internal}, \ec{sign\_internal},
\ec{verify\_challenge\_consistent},
\ec{verify\_reconstructed\_highbits}, \ec{verify\_public\_equation} &
3373--3405. Machine-checked fact. &
Defines the formal keygen, signer, challenge-consistency check, and public
equation check. &
Used directly by \file{HAETAE\_Scheme.ec}, Rejection relation lemmas, and
\file{HAETAE\_Security.ec} correctness. \\
\midrule
\ec{coeff\_norm\_sq}, \ec{poly\_norm\_sq}, \ec{polyvecl\_norm\_sq},
\ec{polyveck\_norm\_sq}, \ec{response\_norm\_ok},
\ec{verify\_norm\_ok}, \ec{verify\_internal} &
3777--3817. Machine-checked fact. &
Defines the formal norm predicates and final verifier conjunction. &
Imports mode norm bounds.  Rejection uses these as acceptance/abort channels;
Security uses \ec{verify\_internal\_sign\_internal} for correctness. \\
\bottomrule
\end{longtable}
}

\paragraph{Lemma ledger for HAETAE Algebra.}
{\footnotesize
\begin{longtable}{p{0.25\linewidth}p{0.14\linewidth}p{0.28\linewidth}p{0.23\linewidth}}
\toprule
\textbf{EasyCrypt lemma(s)} & \textbf{Lines / category} &
\textbf{Mathematical statement paraphrase} &
\textbf{Local premises, imports, downstream role} \\
\midrule
\ec{poly\_zero\_size}, \ec{polyveck\_zero\_size},
\ec{polyvecl\_zero\_size}, \ec{polyvecm\_zero\_size},
\ec{polyveck\_wf\_nth}, \ec{poly\_zero\_wf},
\ec{polyveck\_zero\_wf}, \ec{polyvecl\_zero\_wf} &
39--106. Machine-checked fact. &
Zero objects have the expected lengths; well-formed vectors expose
well-formed rows. &
Imports list facts \ec{size\_nseq}, \ec{allP}, \ec{mem\_nth}.  Used whenever
later scripts take a default row from a public key, commitment, or sample. \\
\midrule
\ec{polyvecl\_sub\_size}, arithmetic \ec{\_wf} lemmas through
\ec{polyvecl\_sub\_wf} &
170--289. Machine-checked fact. &
Formal arithmetic preserves polynomial and vector shape. &
Imports \ec{size\_mkseq}, \ec{mkseqP}, \ec{foldl} induction, and mode
positivity.  Downstream: signature validity, verification matrices, Rejection
state well-formedness. \\
\midrule
\ec{coeff\_mod\_q\_bound}, \ec{poly\_*coeffs\_q\_bound},
\ec{matrix\_vec\_mul\_coeffs\_q\_bound},
\ec{polyveck\_*coeffs\_q\_bound} &
298--425. Machine-checked fact. &
Formal arithmetic constructors produce \(q\)-bounded coefficients. &
Imports integer modulus facts and \ec{smt}.  Used by poly-\(q\) packing bounds
and public-key vector range lemmas. \\
\midrule
Decomposition \ec{\_wf} lemmas and poly-\(q\) bound lemmas,
including \ec{polyveck\_vk\_highbits\_mode23\_polyq\_wf} and
\ec{polyveck\_q\_bound\_mode5\_polyq\_wf} &
483--1057. Machine-checked fact. &
High/low-bit and hint decomposition maps preserve shape and satisfy the
formal poly-\(q\) bounds used for public-key packing. &
Local premises include \(q\)-bounded inputs and, for Mode2/3 VK highbits,
\ec{md = Mode2 \/ md = Mode3}.  Used by keygen public-vector and reference
packing roundtrip lemmas. \\
\midrule
\ec{public\_key\_seed\_poly\_wf},
\ec{public\_key\_seed\_polyvecl\_wf},
\ec{public\_key\_seed\_polyveck\_wf}, seed/XOF size and layout lemmas through
\ec{haetae\_keygen\_counter\_next\_gt} &
1267--1579. Machine-checked fact. &
Seed-derived objects and keygen seed buffers have the expected formal lengths,
and the XOF wrapper matches the SHAKE byte-list presentation. &
Local premises include \ec{0 <= outlen}, \ec{0 <= len}, and
\ec{size sd = seedbytes}.  Imports FIPS202 size/domain facts.  Used by keygen
well-formedness and reference seed-flow documentation. \\
\midrule
\ec{public\_matrix\_seed\_*}, \ec{public\_secret\_vector\_seed\_wf},
\ec{public\_error\_vector\_seed\_wf}, \ec{public\_rounding\_vector\_seed\_wf},
reference/keygen vector \ec{\_wf} and \ec{polyq\_wf} lemmas through
\ec{public\_key\_vector\_seed\_wf} &
1581--1839. Machine-checked fact. &
Formal public matrices, secret/error vectors, rounding vectors, and public-key
vectors are shape-correct and poly-\(q\)-bounded. &
Mode cases split Mode2/3 from Mode5.  Used by public-key packing,
verification-matrix construction, and \file{HAETAE\_Security.ec} correctness. \\
\midrule
\ec{public\_key\_mode23\_column\_wf},
\ec{haetae\_verification\_column\_from\_unpacked\_wf},
\ec{public\_key\_verification\_column\_wf},
\ec{public\_verification\_matrix\_size},
\ec{public\_verification\_matrix\_rows\_wf},
\ec{challenge\_embedding\_wf}, \ec{public\_key\_of\_secret\_wf},
seed pack/unpack lemmas &
1846--1953. Machine-checked fact. &
Verification matrices and challenge embeddings are well-shaped; seed packing
roundtrips under seed-size hypotheses. &
Imports list equality and mode dimensions.  Used by verification public
equations, transcripts, and packed-key recovery. \\
\midrule
\ec{haetae\_polyq\_coeff\_bound\_gt0},
\ec{haetae\_byte\_norm\_*}, \ec{haetae\_unpack15\_coeff0} through
\ec{haetae\_unpack15\_coeff7}, reference packer size/wf lemmas &
1955--2098. Machine-checked fact. &
Byte normalization is bounded/idempotent, and each 15-byte Mode2/3 coefficient
lane decodes to the original formal coefficient under range hypotheses. &
Local premises are explicit numeric ranges.  Imports integer division/modulus
SMT.  Used only for formal reference-style public-key packing. \\
\midrule
\ec{nth\_cat\_size\_plus}, \ec{haetae\_pack\_polyveck\_body\_ref\_nth},
Mode2/3 coefficient-offset lemmas, \ec{haetae\_unpack\_poly\_q\_ref\_pack\_mode23\_*},
\ec{haetae\_pack\_vk\_ref\_roundtrip\_mode23} &
2100--2583. Machine-checked fact. &
For Mode2/3, reference-style public-key packing followed by unpacking returns
the same formal polynomial/vector/key. &
Premises: \ec{md = Mode2 \/ md = Mode3}, seed size, and
\ec{haetae\_polyq\_*wf}.  Downstream: reference packed-key roundtrip and
verification-matrix correctness; not external parser equivalence. \\
\midrule
\ec{haetae\_unpack\_poly\_q\_ref\_pack\_mode5\_*},
\ec{haetae\_unpack\_polyveck\_body\_ref\_pack\_mode5},
\ec{haetae\_pack\_vk\_ref\_roundtrip\_mode5},
\ec{haetae\_public\_key\_ref\_roundtrip},
\ec{haetae\_keygen\_public\_key\_ref\_roundtrip} &
2585--2769. Machine-checked fact. &
For Mode5 and mode-generic reference keys, formal reference packing roundtrips
under the displayed well-formedness premises. &
Premises: seed size and poly-\(q\) well-formedness.  Used by concrete/ref
verification-matrix correctness lemmas. \\
\midrule
\ec{haetae\_pack\_polyveck\_body\_size},
\ec{haetae\_unpack\_poly\_q\_pack\_at},
\ec{haetae\_unpack\_polyveck\_body\_pack\_at},
\ec{haetae\_public\_key\_roundtrip}, public-key encoding/unpacking wf lemmas &
2771--2964. Machine-checked fact. &
The simpler model public-key packer is size-correct and roundtrips for
well-formed inputs. &
Premises: \ec{poly\_wf}, \ec{polyveck\_wf}, or unpacked-key well-formedness.
Used by keygen public-key roundtrip and matrix-correctness lemmas. \\
\midrule
\ec{packed\_haetae\_reference\_public\_key\_*},
\ec{concrete\_haetae\_reference\_*wf},
\ec{haetae\_reference\_keygen\_public\_key\_ref\_roundtrip} &
2966--3019. Machine-checked fact. &
Reference-named formal public keys are well-formed and roundtrip through the
reference-style packer. &
Imports earlier reference packer lemmas.  Used by concrete/reference
verification-matrix facts.  Still a formal-object theorem. \\
\midrule
\ec{public\_key\_of\_secret\_relation},
\ec{public\_key\_of\_secret\_equation\_holds} &
3022--3037. Machine-checked fact. &
The public key derived from a secret seed satisfies the formal public-key
relation/equation. &
Definitional proof.  Used by transcript validity, Rejection signing-attempt
relations, and Security correctness; it is not a hardness theorem. \\
\midrule
Deterministic signing-field lemmas:
\ec{deterministic\_*wf}, \ec{deterministic\_*unit},
\ec{commitment\_*wf}, \ec{commitment\_lowbits\_entropy\_token},
\ec{challenge\_from\_seed\_*}, \ec{challenge\_hash\_*},
\ec{response\_*}, \ec{signing\_sample\_*}, \ec{hint\_*},
\ec{message\_hash\_size}, \ec{challenge\_hash\_size} &
3171--3371. Machine-checked fact. &
Deterministic signer fields are well-formed/unit/sparse and have expected
hash lengths; lowbits coefficient zero equals the signing entropy token. &
Imports list construction facts and mode dimensions.  Used by Rejection
field-consistency lemmas, HopGames simulation rewrites, and counted-ROM
prequery/token arguments. \\
\midrule
\ec{keygen\_internal\_valid}, \ec{public\_key\_of\_keygen},
keygen equation lemmas, reference-keygen equation lemmas,
\ec{haetae\_reference\_keygen\_seed\_flow\_wf\_ok},
packed-key and verification-matrix identity lemmas through
\ec{public\_verification\_matrix\_packedE} &
3407--3605. Machine-checked fact. &
Honest formal keygen returns a valid keypair, and the various public-key and
reference-keygen views unfold to the same formal equations. &
Some lemmas require Mode2/3 premises or seed-size premises.  Used by Security
correctness and by matrix/public-key bridge lemmas in transcript games. \\
\midrule
\ec{packed\_haetae\_keygen\_verification\_matrixE},
\ec{honest\_public\_verification\_matrix\_eq\_packed\_keygen},
\ec{concrete\_*verification\_matrix\_*correct},
\ec{concrete\_*verification\_matrix\_*size},
\ec{concrete\_*verification\_matrix\_*rows\_wf} &
3607--3758. Machine-checked fact. &
The verification matrix obtained from honest keygen or packed/unpacked formal
keys agrees with the corresponding formal concrete/reference matrix and has
the expected shape. &
Premises include \ec{size sd = seedbytes} for keygen concrete matrix
roundtrip.  Used wherever HopGames or Security move between public keys,
encoded keys, and verification matrices. \\
\midrule
\ec{valid\_signature\_sign\_internal},
\ec{verify\_challenge\_consistent\_sign\_internal} &
3760--3775. Machine-checked fact. &
The formal signer emits a signature satisfying the formal validity predicate
and challenge consistency under \ec{public\_key\_of\_secret}. &
Depends on deterministic-field \ec{\_wf}, message/challenge size, and
challenge/hash definitions.  Used directly in Rejection signing-attempt
simulation and Security correctness. \\
\midrule
\ec{secret\_norm\_ok\_current},
\ec{response\_norm\_ok\_current},
\ec{verify\_norm\_ok\_current},
\ec{verify\_internal\_sign\_internal} &
3819--3883. Machine-checked fact. &
Honest formal secret/signature values satisfy norm bounds, and verification of
\ec{sign\_internal} under \ec{public\_key\_of\_secret} succeeds. &
Imports mode bounds and deterministic unit-vector lemmas.  This is the local
formal correctness theorem cited by \file{HAETAE\_Security.ec}; it is not
EUF-CMA security and not implementation equivalence. \\
\bottomrule
\end{longtable}
}

\paragraph{Algebra import and boundary ledger.}
\begin{longtable}{p{0.24\linewidth}p{0.66\linewidth}}
\toprule
\textbf{Category} & \textbf{Status inside \file{HAETAE\_Algebra.ec}} \\
\midrule
Machine-checked facts &
All displayed \ec{lemma} statements in this file, including shape/range
preservation, packing roundtrips, keygen equations, public verification-matrix
facts, norm facts, and \ec{verify\_internal\_sign\_internal}. \\
\midrule
Imported/library facts &
\ec{AllCore}, \ec{IntDiv}, and \ec{List} supply list equality, list indexing,
sequence sizes, integer division/modulus, and SMT-dischargeable arithmetic.
\file{HAETAE\_Params.ec} supplies positivity and mode constants.
\file{HAETAE\_FIPS202.ec} supplies SHAKE size/rate/domain facts such as
\ec{shake256\_size}, \ec{shake256\_squeeze\_size},
\ec{shake256\_rate\_bytesE}, and \ec{shake256\_domain\_separatorE}. \\
\midrule
External cryptographic assumptions &
None are introduced locally.  This file does not state MLWE, Module-SIS, ROM,
Fiat-Shamir, or rejection-sampling security assumptions. \\
\midrule
Conditional premises &
No global cryptographic premise is assumed locally.  Individual lemmas still
have mathematical hypotheses, for example \ec{size sd = seedbytes},
\ec{md = Mode2 \/ md = Mode3}, \ec{polyveck\_wf md b}, or
\ec{haetae\_polyveck\_polyq\_coeffs\_wf md b}; these hypotheses are theorem
premises, not established facts about arbitrary inputs. \\
\midrule
Explanatory prose boundary &
Names containing ``reference'' or ``concrete'' identify formal EasyCrypt
objects in this file.  They do not by themselves certify a C/Jasmin
implementation, the HAETAE standard document, or SHAKE pseudorandomness. \\
\bottomrule
\end{longtable}

\paragraph{Final correctness dependency chain.}
The local correctness theorem decomposes exactly according to the conjunction
in \ec{verify\_internal}.  A reader replaying the proof should check these
four obligations rather than treating \ec{verify\_internal\_sign\_internal} as
an opaque correctness fact.

\begin{longtable}{p{0.20\linewidth}p{0.33\linewidth}p{0.37\linewidth}}
\toprule
\textbf{\ec{verify\_internal} conjunct} & \textbf{Discharged by} &
\textbf{Proof-script dependency reading} \\
\midrule
\ec{valid\_signature md sig} &
\ec{valid\_signature\_sign\_internal}. &
Uses well-formedness of commitment highbits, lowbits, message hash, challenge,
response, auxiliary response, and hint.  These depend on deterministic
constructor \ec{\_wf} lemmas and imported list-size facts. \\
\midrule
\ec{verify\_challenge\_consistent md pk m ctx sig} &
\ec{verify\_challenge\_consistent\_sign\_internal}. &
Pure definitional unfolding of \ec{sign\_internal},
\ec{commitment\_challenge}, and \ec{challenge\_hash}; no hardness or
probability reasoning enters. \\
\midrule
\ec{verify\_public\_equation md pk sig} &
Direct rewrite inside \ec{verify\_internal\_sign\_internal}. &
The proof expands \ec{verify\_public\_equation},
\ec{verify\_reconstructed\_highbits}, \ec{sign\_internal}, and
\ec{commitment\_highbits}, then observes that the signature's auxiliary
response field is the same value used to build the commitment highbits. \\
\midrule
\ec{verify\_norm\_ok md sig} &
\ec{verify\_norm\_ok\_current}. &
Uses the unit-vector lemmas for deterministic response, auxiliary response,
and hint, then mode case analysis against the constants in
\file{HAETAE\_Params.ec}. \\
\bottomrule
\end{longtable}

\paragraph{Packing and roundtrip obligation map.}
\begin{longtable}{p{0.18\linewidth}p{0.31\linewidth}p{0.41\linewidth}}
\toprule
\textbf{Range} & \textbf{Main obligations} & \textbf{Premises and boundary} \\
\midrule
Mode2/Mode3 reference packer &
The offset lemmas for \ec{ofs0} through \ec{ofs7} prove that each 15-byte
block decodes eight 15-bit coefficients, then
\ec{haetae\_unpack\_poly\_q\_ref\_pack\_mode23\_coeff} and vector-body lemmas
lift coefficient equality to polynomials and vectors. &
Requires \ec{md = Mode2 \/ md = Mode3} and
\ec{haetae\_polyq\_coeffs\_wf}.  Arithmetic is discharged by imported integer
SMT plus byte-normalization lemmas. \\
\midrule
Mode5 reference packer &
\ec{haetae\_unpack\_poly\_q\_ref\_pack\_mode5\_coeff},
\ec{haetae\_unpack\_poly\_q\_ref\_pack\_mode5\_at}, and vector-body lemmas
prove the 16-bit two-byte decoding roundtrip. &
Requires \ec{haetae\_polyq\_coeffs\_wf Mode5}.  The checked bound is formal
poly-\(q\) well-formedness, not an implementation parser theorem. \\
\midrule
Simple model packer &
\ec{haetae\_unpack\_poly\_q\_pack\_at},
\ec{haetae\_unpack\_polyveck\_body\_pack\_at}, and
\ec{haetae\_public\_key\_roundtrip}. &
Requires \ec{poly\_wf} or \ec{polyveck\_wf}.  This packer is a model
roundtrip surface distinct from the reference-style byte-packing surface. \\
\midrule
Public-key recovery &
\ec{haetae\_keygen\_public\_key\_roundtrip},
\ec{haetae\_keygen\_public\_key\_ref\_roundtrip}, and verification-matrix
correctness lemmas. &
Requires seed-size or unpacked-key well-formedness premises.  These lemmas
justify downstream formal verification matrices, but not external
implementation equivalence. \\
\bottomrule
\end{longtable}

\section{Script companion coverage for HAETAE Distributions}

\file{HAETAE\_Distributions.ec} defines the distributions used by key
generation, random-oracle outputs, signing coins, and signing-attempt samples.
Its most important boundary is the distinction between a checked bounded
carrier and a paper-faithful exact HAETAE hyperball sampler.

\begin{longtable}{p{0.27\linewidth}p{0.38\linewidth}p{0.25\linewidth}}
\toprule
\textbf{Script object family} & \textbf{Mathematical reading} & \textbf{Downstream role} \\
\midrule
Byte, seed, CRH, coefficient, polynomial, matrix, and challenge distributions &
\ec{dbyte}, \ec{dseed}, \ec{dcrh}, \ec{dcoeff}, \ec{dpoly},
\ec{dpolyveck}, \ec{dpolyvecl}, \ec{dmatrix}, and \ec{dchallenge} are product
or mapped distributions over the formal carriers. &
\ec{HAETAE\_ROM.ec} uses these as random-oracle output components.
\file{HAETAE\_Scheme.ec} samples \ec{dseed} and \ec{drandom\_coins}. \\
\midrule
Signing randomness token family &
\ec{signing\_entropy\_token\_cardinality},
\ec{signing\_entropy\_token\_to\_coins},
\ec{signing\_randomness\_domain}, \ec{dsigning\_entropy\_token},
\ec{drandom\_coins}, and \ec{signing\_randomness\_point\_bound} encode a
58-bit entropy token embedded into signing coins. &
Feeds the sampler-prequery and ROM min-entropy bounds in
\file{HAETAE\_HopGames.ec} and \file{HAETAE\_ROM\_Programming.ec}. \\
\midrule
Signing sample pairs &
\ec{signing\_sample\_pair}, projections, \ec{signing\_sample\_pair\_wf},
\ec{signing\_sample\_pair\_unit}, \ec{signing\_sample\_pair\_bounded}, and
\ec{signing\_sample\_pair\_sample\_ok} define the shape of the two-vector
sample used inside signing attempts. &
\file{HAETAE\_Rejection.ec} lifts these samples into full signing-attempt
states and rejection predicates. \\
\midrule
Structural and bounded sample distributions &
\ec{dstructural\_signing\_sample\_pair},
\ec{dbounded\_unit\_signing\_sample\_pair}, and
\ec{dchecked\_hyperball\_signing\_sample\_pair} provide deterministic,
unit-product, and bounded-product sampling surfaces. &
Used by hop games to relate real signing, paper simulation, and exact-hyperball
sampling surfaces. \\
\midrule
Hyperball interface &
\ec{hyperball\_coeff\_ok}, \ec{hyperball\_sample\_ok},
\ec{dhyperball\_*}, point-bound operations, and
\ec{dexact\_hyperball\_signing\_sample\_pair} expose the sampler interface
needed by rejection-sampling lifts. &
The current checked implementation sets the exact hyperball interface to a
bounded product distribution.  The comments in the script mark the paper
sphere/hyperball law as replacement/correspondence work. \\
\bottomrule
\end{longtable}

\begin{longtable}{p{0.31\linewidth}p{0.43\linewidth}p{0.16\linewidth}}
\toprule
\textbf{Lemma family} & \textbf{Checked fact} & \textbf{Boundary} \\
\midrule
Losslessness lemmas &
\ec{byte\_distribution\_lossless}, \ec{seed\_distribution\_lossless},
\ec{signing\_coin\_distribution\_lossless}, and hyperball/sample-pair
losslessness lemmas show that each formal distribution has total mass \(1\). &
Machine-checked. \\
\midrule
Support/domain lemmas &
\ec{signing\_coin\_distribution\_domain},
\ec{signing\_sample\_pair\_of\_coins\_sample\_ok},
unit/bounded support lemmas, and hyperball support \ec{P} lemmas connect
sampling to shape predicates. &
Machine-checked. \\
\midrule
Point-bound lemmas &
\ec{signing\_coin\_distribution\_token\_point\_bound},
\ec{dlist\_point\_bound}, \ec{hyperball\_poly\_point\_boundE},
\ec{hyperball\_polyvecl\_point\_boundE}, and
\ec{exact\_hyperball\_signing\_sample\_pair\_point\_bound} provide the
probability terms consumed by sampler-prequery bounds. &
Machine-checked for the formal distributions. \\
\midrule
Checked/exact equality &
\ec{exact\_hyperball\_signing\_sample\_pair\_checkedE} proves that the current
\ec{dexact\_hyperball\_signing\_sample\_pair} equals the checked bounded
carrier. &
Machine-checked equality, but not a paper-sampler equivalence theorem. \\
\bottomrule
\end{longtable}

Boundary: the file checks distribution properties for the formal carriers.
It does not by itself prove that the formal signing randomness and bounded
sample law are the final HAETAE paper distributions.

\subsection{Lemma-ledger companion for HAETAE Distributions}

\paragraph{Lemma-ledger criticism for HAETAE Distributions.}
The previous distribution companion made the main boundary visible, but it was
still too coarse for reconstructing proof obligations.  A reader needs to know
which names are merely product-distribution definitions, which support and
point-probability facts are machine-checked, where the deterministic
\ec{signing\_sample\_pair\_of\_coins} surface differs from an entropy-bearing
sample surface, and where the word ``hyperball'' is only the formal interface
name of the current bounded-product carrier.  Without a ledger, it is easy to
overclaim that \file{HAETAE\_Distributions.ec} proves paper-sampler
correctness, SHAKE-to-ROM behavior, or the sampling-distance premise used in
later hop games.  It does not.  It proves losslessness, support, domain, and
point-bound facts for the formal distributions below; the paper-sampler and
ROM correspondence obligations remain separate conditional or external
boundaries unless a later checked lemma discharges them.

\paragraph{Imported-interface and assumption ledger for HAETAE Distributions.}
{\footnotesize
\begin{longtable}{p{0.23\linewidth}p{0.16\linewidth}p{0.28\linewidth}p{0.23\linewidth}}
\toprule
\textbf{EasyCrypt object} & \textbf{Lines / category} &
\textbf{Mathematical statement paraphrase} &
\textbf{Local premises, imports, downstream role} \\
\midrule
\ec{AllCore}, \ec{Distr}, \ec{DInterval}, \ec{DList}, \ec{IntDiv},
\ec{List}, \ec{Real}, \ec{StdOrder} &
1. Imported/library fact. &
Provides distributions, intervals, product/list distributions, real
probabilities, order reasoning, integer division, lists, \ec{mu}, support, and
losslessness lemmas. &
No HAETAE-specific premise.  Used by every distribution proof and inherited by
Rejection and HopGames sampler arguments. \\
\midrule
\file{HAETAE\_Params.ec}, \file{HAETAE\_Algebra.ec},
\ec{HAETAE\_Params}, \ec{HAETAE\_Algebra}, \ec{RealOrder} &
2, 6--8. Imported/library fact. &
Imports the formal mode constants, carrier types, well-formedness predicates,
signing-field constructors, challenge constructor, and real-order tactics. &
All support lemmas are relative to these formal carriers and parameter
definitions.  \file{HAETAE\_Algebra.ec} itself does not import this file. \\
\midrule
Paper sampler and ROM correspondence &
61--67, 142--202. External cryptographic assumption. &
Comments explicitly mark the challenge carrier and hyperball carrier as
structural/formal surfaces rather than paper-faithful correspondence theorems. &
Later \file{HAETAE\_Rejection.ec} and \file{HAETAE\_HopGames.ec} can use these
formal distributions, but equality to the HAETAE paper sampler or SHAKE/ROM
semantics must be supplied separately. \\
\midrule
Structural-to-exact sampling-loss premises &
Not defined in this file. Conditional premise. &
The later \ec{structural\_to\_exact\_hyperball} loss obligations compare
attempt-state distributions built from the sample-pair distributions defined
here. &
The distributions file supplies support and point-bound ingredients, not the
full sampling-distance theorem consumed by \file{HAETAE\_HopGames.ec} and
\file{HAETAE\_Security.ec}. \\
\bottomrule
\end{longtable}
}

\paragraph{Definition ledger for HAETAE Distributions.}
{\footnotesize
\begin{longtable}{p{0.25\linewidth}p{0.14\linewidth}p{0.28\linewidth}p{0.23\linewidth}}
\toprule
\textbf{EasyCrypt object(s)} & \textbf{Lines / category} &
\textbf{Mathematical statement paraphrase} &
\textbf{Local premises, imports, downstream role} \\
\midrule
\ec{signing\_randomness\_source}, \ec{dbyte}, \ec{dseed}, \ec{dcrh},
\ec{signing\_entropy\_token\_cardinality},
\ec{signing\_entropy\_token\_in\_range},
\ec{signing\_randomness\_byte\_ok},
\ec{signing\_entropy\_token\_to\_coins},
\ec{signing\_randomness\_domain}, \ec{dsigning\_entropy\_token},
\ec{dsigning\_randomness\_source},
\ec{signing\_randomness\_source\_valid},
\ec{signing\_randomness\_from\_source},
\ec{dsigning\_randomness}, \ec{drandom\_coins},
\ec{signing\_randomness\_point\_bound},
\ec{signing\_randomness\_token\_spread},
\ec{signing\_randomness\_distribution\_ok} &
10--53. Machine-checked fact. &
Defines byte, seed, CRH, and signing-coin distributions, with a finite entropy
token embedded into formal signing coins and a point-bound predicate for token
events. &
Imports parameter lengths and algebraic coin-token projection.  Used by
\file{HAETAE\_Scheme.ec} signing, \file{HAETAE\_Rejection.ec} attempt-state
sampling, \file{HAETAE\_ROM\_Programming.ec} min-entropy bounds, and
\file{HAETAE\_HopGames.ec} sampler-prequery bounds. \\
\midrule
\ec{dcoeff}, \ec{dpoly}, \ec{dpolyveck}, \ec{dpolyvecl}, \ec{dmatrix} &
55--59. Machine-checked fact. &
Defines uniform coefficient and product list distributions over the formal
polynomial/vector/matrix carriers. &
Imports \(q\), \(n\), \ec{mode\_k}, and \ec{mode\_l}.  Used as primitive
distribution infrastructure; direct security hops mainly consume the signing
sample and attempt-state distributions built later. \\
\midrule
\ec{dchallenge} &
61--66. Machine-checked fact. &
Maps \ec{dcrh} through \ec{challenge\_from\_seed md} to obtain a structural
prefix-sparse challenge carrier. &
Imports the algebraic challenge constructor.  The checked support is local to
that constructor; equality to a specification challenge sampler is not proved
here. \\
\midrule
\ec{signing\_sample\_pair},
\ec{signing\_sample\_pair\_y1},
\ec{signing\_sample\_pair\_y2},
\ec{signing\_sample\_pair\_wf},
\ec{signing\_sample\_pair\_unit},
\ec{signing\_sample\_bound},
\ec{coeff\_bounded\_by}, \ec{poly\_bounded\_by},
\ec{polyvecl\_bounded\_by}, \ec{polyveck\_bounded\_by},
\ec{signing\_sample\_pair\_bounded},
\ec{signing\_sample\_pair\_sample\_ok},
\ec{signing\_sample\_pair\_of\_coins} &
68--115. Machine-checked fact. &
Defines the two-vector signing sample interface, its well-formed/unit/bounded
predicates, and the deterministic sample extracted from formal signing coins. &
Imports algebraic \ec{signing\_sample\_y1} and
\ec{signing\_sample\_y2}.  This is the interface lifted into full signing
attempt states in \file{HAETAE\_Rejection.ec}. \\
\midrule
\ec{dstructural\_signing\_sample\_pair} &
117--120. Machine-checked fact. &
Defines a degenerate distribution concentrated on the deterministic sample
pair from coins. &
Used by structural attempt-state distributions and by loss lemmas comparing
structural samples to exact-hyperball samples. \\
\midrule
\ec{dsigning\_unit\_coeff}, \ec{dsigning\_unit\_poly},
\ec{dsigning\_unit\_polyvecl}, \ec{dsigning\_unit\_polyveck},
\ec{signing\_bounded\_coeff\_cardinality},
\ec{signing\_bounded\_coeff\_point\_bound},
\ec{dsigning\_bounded\_coeff}, \ec{dsigning\_bounded\_poly},
\ec{dsigning\_bounded\_polyvecl}, \ec{dsigning\_bounded\_polyveck},
\ec{dbounded\_unit\_signing\_sample\_pair} &
122--149. Machine-checked fact. &
Defines unit and bounded product sample carriers and a bounded-unit pair
sampler. &
Imports \ec{mode\_ln}, \ec{mode\_k}, \ec{mode\_l}, and \(n\).  Used by
Rejection as an intermediate sampler surface; it is not the paper hyperball
law. \\
\midrule
\ec{dchecked\_hyperball\_signing\_sample\_pair},
\ec{hyperball\_coeff\_ok}, \ec{hyperball\_poly\_ok},
\ec{hyperball\_polyvecl\_ok}, \ec{hyperball\_polyveck\_ok},
\ec{hyperball\_sample\_ok}, \ec{dhyperball\_coeff},
\ec{dhyperball\_poly}, \ec{dhyperball\_polyvecl},
\ec{dhyperball\_polyveck}, hyperball point-bound operations,
\ec{dexact\_hyperball\_signing\_sample\_pair} &
156--203. Machine-checked fact. &
Defines the checked bounded-product hyperball-shaped carrier, its support
predicates, component point bounds, and the current exact interface as a
product of bounded components. &
Used by \file{HAETAE\_Rejection.ec} exact-hyperball attempt states and by
\file{HAETAE\_HopGames.ec} exact-hyperball sampling hops.  The name
``exact'' is formal here; paper-sphere equivalence remains outside this file. \\
\midrule
\ec{signing\_sample\_pair\_distribution\_ok} &
205--209. Machine-checked fact. &
Names the contract that a sample-pair distribution is lossless and supports
only formally sample-ok pairs. &
Used by sampler-lifting lemmas in \file{HAETAE\_Rejection.ec}. \\
\bottomrule
\end{longtable}
}

\paragraph{Lemma ledger for HAETAE Distributions.}
{\footnotesize
\begin{longtable}{p{0.25\linewidth}p{0.14\linewidth}p{0.28\linewidth}p{0.23\linewidth}}
\toprule
\textbf{EasyCrypt lemma(s)} & \textbf{Lines / category} &
\textbf{Mathematical statement paraphrase} &
\textbf{Local premises, imports, downstream role} \\
\midrule
\ec{byte\_distribution\_lossless},
\ec{seed\_distribution\_lossless},
\ec{crh\_distribution\_lossless} &
211--223. Machine-checked fact. &
The byte, seed, and CRH product distributions are lossless. &
Imports \ec{dinter\_ll}, \ec{dlist\_ll}, and parameter lengths.  Used by
\file{HAETAE\_Scheme.ec}, \file{HAETAE\_ROM.ec}, and challenge-distribution
losslessness. \\
\midrule
\ec{signing\_entropy\_token\_cardinality\_positive},
\ec{signing\_entropy\_token\_distribution\_lossless},
\ec{signing\_randomness\_source\_distribution\_lossless},
\ec{signing\_coin\_distribution\_lossless} &
226--248. Machine-checked fact. &
The entropy-token and signing-coin distributions are well-defined and
lossless. &
Imports interval and map losslessness.  Used by signing procedures,
attempt-state samplers, and HopGames probability calculations over coins. \\
\midrule
\ec{signing\_entropy\_token\_to\_coins\_tokenK},
\ec{signing\_randomness\_from\_source\_tokenK},
\ec{signing\_entropy\_token\_to\_coins\_size},
\ec{signing\_entropy\_token\_to\_coins\_byte\_domain},
\ec{signing\_entropy\_token\_to\_coins\_domain},
\ec{signing\_randomness\_from\_source\_domain} &
251--311. Machine-checked fact. &
Encoding a token into coins recovers the same token, has seed length, and
satisfies the byte/domain predicates under the token-range premise. &
Local premise for byte/domain lemmas is
\ec{signing\_entropy\_token\_in\_range x}.  Used to prove coin support and
ROM prequery/min-entropy facts. \\
\midrule
\ec{signing\_entropy\_token\_distribution\_point\_bound},
\ec{signing\_randomness\_source\_distribution\_point\_bound},
\ec{signing\_coin\_distribution\_token\_point\_bound},
\ec{signing\_coin\_distribution\_token\_spread},
\ec{signing\_coin\_distribution\_domain},
\ec{signing\_coin\_distribution\_domain\_full},
\ec{signing\_coin\_distribution\_ok} &
314--387. Machine-checked fact. &
Coin sampling has the stated token point bound, support domain, full-domain
mass, and distribution-ok contract. &
Imports \ec{mu}, \ec{dmapE}, interval support, and real-order facts.  Used by
\file{HAETAE\_ROM\_Programming.ec} and HopGames sampler-prequery bounds. \\
\midrule
\ec{coeff\_distribution\_lossless},
\ec{poly\_distribution\_lossless},
\ec{polyveck\_distribution\_lossless},
\ec{polyvecl\_distribution\_lossless},
\ec{matrix\_distribution\_lossless},
\ec{challenge\_distribution\_lossless} &
390--420. Machine-checked fact. &
Primitive coefficient, polynomial, vector, matrix, and structural challenge
distributions are lossless. &
Imports \ec{dinter\_ll}, \ec{dlist\_ll}, \ec{dmap\_ll}, and CRH
losslessness.  These facts justify well-formed random-oracle output carriers,
not SHAKE security. \\
\midrule
\ec{signing\_sample\_pair\_of\_coins\_wf},
\ec{signing\_sample\_pair\_of\_coins\_unit},
\ec{signing\_sample\_bound\_gt0},
\ec{unit\_coeff\_bounded\_by\_signing\_sample\_bound},
\ec{poly\_unit\_bounded\_by\_signing\_sample\_bound},
\ec{polyvecl\_unit\_bounded\_by\_signing\_sample\_bound},
\ec{polyveck\_unit\_bounded\_by\_signing\_sample\_bound},
\ec{signing\_sample\_pair\_unit\_bounded},
\ec{signing\_sample\_pair\_of\_coins\_sample\_ok} &
423--513. Machine-checked fact. &
The deterministic sample pair derived from coins is well-formed, unit, and
bounded by the formal signing-sample bound. &
Imports algebraic signer-field \ec{\_wf}/\ec{\_unit} lemmas and
\ec{mode\_ln\_gt0}.  Used by structural sample support and current-signing
attempt lemmas. \\
\midrule
\ec{structural\_signing\_sample\_pair\_lossless},
\ec{structural\_signing\_sample\_pair\_support},
\ec{structural\_signing\_sample\_pair\_sample\_ok},
\ec{structural\_signing\_sample\_pair\_distribution\_ok} &
516--543. Machine-checked fact. &
The deterministic structural sample-pair distribution is lossless and supports
sample-ok pairs. &
Imports \ec{dunit\_ll} and deterministic sample lemmas.  Used by
\file{HAETAE\_Rejection.ec} structural attempt-state samplers. \\
\midrule
\ec{signing\_unit\_coeff\_lossless},
\ec{signing\_unit\_coeff\_support},
\ec{signing\_bounded\_coeff\_cardinality\_positive},
\ec{signing\_bounded\_coeff\_lossless},
\ec{signing\_bounded\_coeff\_support},
\ec{signing\_bounded\_coeff\_supportP},
\ec{signing\_bounded\_coeff\_point\_bound},
\ec{signing\_bounded\_coeff\_point\_bound\_nonnegative},
\ec{dlist\_size\_point\_bound}, \ec{dlist\_point\_bound},
\ec{hyperball\_coeff\_point\_bound\_nonnegative},
\ec{hyperball\_coeff\_point\_boundE} &
553--674. Machine-checked fact. &
Unit and bounded coefficient distributions have the stated support,
losslessness, nonnegative point bounds, and list-point-bound lifting lemmas. &
Imports \ec{duniform}, \ec{dinter}, list-distribution, and real
multiplication/order facts.  Used by hyperball polynomial/vector point bounds
and sampling-loss bookkeeping. \\
\midrule
\ec{signing\_unit\_poly\_lossless},
\ec{signing\_unit\_poly\_support},
\ec{signing\_bounded\_poly\_lossless},
\ec{signing\_bounded\_poly\_support},
\ec{signing\_bounded\_poly\_supportP},
\ec{hyperball\_poly\_point\_bound\_nonnegative},
\ec{hyperball\_poly\_lossless},
\ec{hyperball\_poly\_support},
\ec{hyperball\_poly\_supportP},
\ec{hyperball\_poly\_point\_boundE} &
676--777. Machine-checked fact. &
Polynomial unit/bounded/hyperball distributions are lossless, have the
intended support predicates, and satisfy the product point bound. &
Premises include list sizes governed by \(n\).  Used by vector sampler support
and exact-hyperball pair point-bound proofs. \\
\midrule
Vector support, losslessness, and point-bound families for
\ec{dsigning\_unit\_polyvecl}, \ec{dsigning\_unit\_polyveck},
\ec{dsigning\_bounded\_polyvecl}, \ec{dsigning\_bounded\_polyveck},
\ec{dhyperball\_polyvecl}, and \ec{dhyperball\_polyveck} &
778--1027. Machine-checked fact. &
Product vector samplers are lossless, support the corresponding unit/bounded
or hyperball predicates, and have nonnegative point bounds. &
Imports \ec{mode\_l\_gt0}, \ec{mode\_k\_gt0}, and polynomial support/point
lemmas.  Used by checked and exact sample-pair support and by Rejection
attempt-state support equivalences. \\
\midrule
\ec{bounded\_unit\_signing\_sample\_pair\_lossless},
\ec{bounded\_unit\_signing\_sample\_pair\_support},
\ec{bounded\_unit\_signing\_sample\_pair\_sample\_ok},
\ec{bounded\_unit\_signing\_sample\_pair\_distribution\_ok},
\ec{checked\_hyperball\_signing\_sample\_pair\_lossless},
\ec{checked\_hyperball\_signing\_sample\_pair\_support},
\ec{checked\_hyperball\_signing\_sample\_pair\_distribution\_ok},
\ec{hyperball\_sample\_ok\_sample\_ok},
\ec{exact\_hyperball\_signing\_sample\_pair\_lossless},
\ec{exact\_hyperball\_signing\_sample\_pair\_support},
\ec{exact\_hyperball\_signing\_sample\_pair\_supportP},
\ec{exact\_hyperball\_signing\_sample\_pair\_sample\_ok},
\ec{exact\_hyperball\_signing\_sample\_pair\_distribution\_ok} &
1029--1182. Machine-checked fact. &
The bounded-unit, checked-hyperball, and current exact-hyperball pair
distributions are lossless and satisfy the sample-pair distribution contract. &
Imports component vector support/losslessness.  Used directly by
\file{HAETAE\_Rejection.ec} sampler-ok lemmas and by HopGames exact-hyperball
sampling procedures. \\
\midrule
\ec{exact\_hyperball\_signing\_sample\_pair\_point\_bound},
\ec{exact\_hyperball\_signing\_sample\_pair\_checkedE} &
1192--1208. Machine-checked fact. &
The current exact-hyperball pair distribution has the product point bound and
is equal to the checked bounded-product carrier. &
Imports product-distribution point-bound facts.  This is the strongest local
theorem boundary for sampling: it identifies two formal carriers, not a
paper-sampler equivalence or a rejection-loop correspondence proof. \\
\bottomrule
\end{longtable}
}

\section{Script companion coverage for HAETAE Rejection}

\file{HAETAE\_Rejection.ec} is the bridge between deterministic signing fields,
attempt-state sampling, transcript logs, and rejection-sampling loss terms.
It is where many later hop premises get their vocabulary.

\begin{longtable}{p{0.26\linewidth}p{0.39\linewidth}p{0.25\linewidth}}
\toprule
\textbf{Script object family} & \textbf{Mathematical reading} & \textbf{Downstream role} \\
\midrule
Attempt-state records &
\ec{signing\_transcript\_record}, \ec{signing\_attempt\_sample}, and
\ec{signing\_attempt\_state} package coins, sample vectors, commitment
fields, challenge, response vectors, hint, and signature. &
\file{HAETAE\_HopGames.ec} samples and transforms these states in
paper-simulation and exact-hyperball games. \\
\midrule
Field reconstruction predicates &
\ec{signing\_attempt\_state\_programmed\_lowbits},
\ec{signing\_attempt\_state\_lowbits\_consistent},
\ec{signing\_attempt\_state\_programmed\_challenge},
\ec{signing\_attempt\_state\_public\_equation\_consistent}, and
\ec{signing\_attempt\_state\_fields\_match\_signature}. &
Used to prove that a sampled attempt can be read as a valid transcript and as
a paper-simulation signature. \\
\midrule
Rejection predicates &
\ec{signing\_attempt\_state\_reject1\_*},
\ec{signing\_attempt\_state\_reject2\_*},
\ec{signing\_attempt\_state\_pack\_*},
\ec{signing\_attempt\_state\_reference\_rejection\_*},
\ec{signing\_attempt\_state\_rejection\_*}, and
\ec{signing\_attempt\_state\_aborts} split the rejection event into named
components. &
Hop games charge or erase these events when moving from real signing to
paper-simulation and exact-hyperball sampling games. \\
\midrule
Attempt distributions &
\ec{dsigning\_attempt\_state},
\ec{dstructural\_signing\_attempt\_state},
\ec{dbounded\_unit\_signing\_attempt\_state},
\ec{dchecked\_hyperball\_signing\_attempt\_state}, and
\ec{dexact\_hyperball\_signing\_attempt\_state} lift sample-pair
distributions to complete attempt states. &
Security/hop lemmas compare success probabilities under these distributions. \\
\midrule
Signing and transcript relations &
\ec{signing\_attempt\_relation}, \ec{signing\_simulation\_relation},
\ec{signing\_transcript\_relation}, \ec{signing\_record\_relation},
\ec{signing\_record\_log\_sound}, and \ec{transcript\_log\_valid} connect
attempts to transcripts. &
Used by transcript-log invariants in \file{HAETAE\_HopGames.ec} and by ROM
programming no-failure predicates. \\
\midrule
Loss obligations &
\ec{structural\_to\_exact\_hyperball\_sample\_pair\_*},
\ec{structural\_to\_exact\_hyperball\_coin\_sample\_pair\_*},
\ec{structural\_to\_exact\_hyperball\_attempt\_loss\_obligation}, and
\ec{rejection\_sampling\_bound\_obligation} name the loss assumptions and
nonnegativity premises. &
Some are conditional premises for hop lemmas.  The nonnegativity obligation is
proved from \file{HAETAE\_Reductions.ec} terms. \\
\bottomrule
\end{longtable}

\begin{longtable}{p{0.31\linewidth}p{0.42\linewidth}p{0.17\linewidth}}
\toprule
\textbf{Lemma family} & \textbf{Checked fact} & \textbf{Boundary} \\
\midrule
Current-signing field equations &
\ec{signing\_attempt\_state\_of\_coins\_signatureE},
\ec{signing\_attempt\_state\_of\_coins\_lowbitsE},
\ec{signing\_attempt\_state\_of\_coins\_challengeE},
\ec{signing\_attempt\_state\_of\_coins\_field\_accepts},
\ec{signing\_attempt\_state\_of\_coins\_accepts\_current}, and
\ec{signing\_attempt\_state\_of\_coins\_rejection\_accepts\_current} prove
that the formal current signer does not abort. &
Machine-checked for the formal signer. \\
\midrule
Sampler lifting and support &
\ec{dsigning\_attempt\_state\_from\_sample\_pair\_sampler\_lossless},
\ec{dexact\_hyperball\_signing\_attempt\_state\_supportP}, and distribution
support lemmas connect samples, coins, and full states. &
Machine-checked for supplied sampler surfaces. \\
\midrule
Abort decomposition and zero facts &
\ec{signing\_attempt\_state\_reference\_rejection\_aborts\_componentE},
union-bound lemmas, \ec{dexact\_hyperball\_signing\_attempt\_state\_reject1\_abort\_zero},
\ec{dexact\_hyperball\_signing\_attempt\_state\_reject2\_abort\_zero}, and
related zero-abort lemmas show which rejection components vanish under the
checked exact-hyperball attempt state. &
Machine-checked; relies on the current exact-hyperball carrier definition. \\
\midrule
Loss-obligation propagation &
\ec{structural\_to\_exact\_hyperball\_sample\_pair\_loss\_from\_point\_loss},
\ec{structural\_to\_exact\_hyperball\_coin\_sample\_pair\_loss\_from\_sample\_pair\_loss},
and \ec{structural\_to\_exact\_hyperball\_attempt\_loss\_from\_coin\_sample\_pair\_loss}
lift a point/sample loss premise to an attempt-state loss premise. &
Conditional unless the initial point-loss obligation is supplied. \\
\midrule
Record-log soundness &
\ec{sign\_internal\_record\_relation},
\ec{signing\_record\_relation\_sound},
\ec{signing\_record\_log\_sound\_cons}, and
\ec{signing\_record\_log\_sound\_transcripts\_valid} turn signing records into
valid transcript logs. &
Machine-checked; consumed by HopGames transcript invariants. \\
\bottomrule
\end{longtable}

Boundary: this file is the precise place where rejection-sampling claims can
become overclaimed.  The checked scripts prove many zero-abort and propagation
facts for the formal carrier.  A paper-faithful sampler replacement still
needs its own checked correspondence or must remain a conditional premise.

\subsection{Line-by-line proof-script companion for HAETAE Rejection}

Focused criticism for this pass: the previous report named the rejection
object families, but it did not let a reader reconstruct the proof boundary of
\file{HAETAE\_Rejection.ec}.  It underexplained that many no-abort theorems are
consequences of the current formal signer setting branch bytes to zero and of
the current exact-hyperball support being lifted through
\ec{signing\_attempt\_state\_of\_sample\_pair}; it also did not expose the
conditional loss-obligation chain from point loss to full attempt-state loss.
The line-range companion below records the proof obligations that a reader
must track when replaying the EasyCrypt script.

\begin{longtable}{p{0.13\linewidth}p{0.19\linewidth}p{0.21\linewidth}p{0.37\linewidth}}
\toprule
\textbf{Lines} & \textbf{Script content} & \textbf{Classification} &
\textbf{Companion reading and downstream role} \\
\midrule
1--14 &
Imports, theory wrapper, and real-order import. &
Imported/library facts. &
\ec{Distr}, \ec{Real}, \ec{StdOrder}, and \ec{RealSeries} provide
probability mass \ec{mu}, support, \ec{dmap}, \ec{dlet}, union bounds,
summability, and nonnegative real arithmetic.  \file{HAETAE\_Algebra.ec},
\file{HAETAE\_Distributions.ec}, \file{HAETAE\_Events.ec},
\file{HAETAE\_Reductions.ec}, and \file{HAETAE\_Transcript.ec} provide the
domain-specific predicates used below. \\
\midrule
16--96 &
Attempt records, samples, state tuple layout, projections, rejection-branch
byte/bit, and reject2 balance vectors. &
Machine-checked definitions. &
These declarations define the state shape that hop games manipulate.  The
branch bit is read from signature component 6; because \ec{sign\_internal}
sets that component to zero, later current-signing and sample-pair no-abort
lemmas can prove reject2 is clear. \\
\midrule
102--163 &
Programmed lowbits, lowbits consistency, programmed challenge, reconstructed
highbits, public-equation consistency, signature reconstruction, and field
acceptance. &
Machine-checked definitions. &
This block states the field-equality obligations needed to recognize a
sampled attempt as a valid formal signature/transcript.  Downstream
\file{HAETAE\_HopGames.ec} invariants use these predicates when comparing
real signing, simulated signing, and ROM-programmed transcripts. \\
\midrule
166--199 &
\ec{signing\_attempt\_state\_of\_coins},
\ec{signing\_attempt\_sample\_of\_pair}, and
\ec{signing\_attempt\_state\_of\_sample\_pair}. &
Machine-checked definitions. &
The first constructor builds the current formal signing attempt from coins.
The second and third constructors replace the sampled pair while keeping the
deterministic signing fields from \file{HAETAE\_Algebra.ec}.  This is the
bridge used by structural, bounded, checked-hyperball, and exact-hyperball
attempt distributions. \\
\midrule
203--339 &
Acceptance predicates, reject1/reject2/pack/reference-rejection predicates,
verification acceptance, full rejection acceptance, abort predicates, and
\ec{signing\_accepts}. &
Machine-checked definitions; explanatory rejection boundary. &
The script decomposes rejection into named events.  These definitions are not
probability bounds by themselves; bounds appear only in later lemmas using
\ec{mu}.  The names are the event vocabulary consumed by hop-game failure
terms. \\
\midrule
345--421 &
Attempt-state distributions and sample-pair shape predicates. &
Machine-checked definitions plus imported distribution interface. &
\ec{dsigning\_attempt\_state} maps \ec{drandom\_coins}; the other
distributions use \ec{dlet} over coins and then map a sample-pair
distribution into a full attempt.  Losslessness and support are proved later,
conditional on sampler-ok predicates supplied by \file{HAETAE\_Distributions.ec}. \\
\midrule
438--488 &
Signing relation, simulation relation, transcript relation, record
projections, record log soundness, and transcript-log validity. &
Machine-checked definitions. &
These predicates connect rejection attempts to transcript reasoning.  The
security games use them to say that a logged signing query corresponds to a
valid formal transcript under the public key. \\
\midrule
491--523 &
Point, sample-pair, coin-sample-pair, and attempt-state loss obligations, plus
nonnegativity obligation. &
Conditional premises and machine-checked definitions. &
\ec{structural\_to\_exact\_hyperball\_sample\_pair\_point\_loss\_obligation}
is not proved in this file.  It is a conditional premise used to derive the
larger distribution-comparison obligations.  By contrast,
\ec{rejection\_sampling\_bound\_obligation} is later proved from nonnegativity
lemmas imported from \file{HAETAE\_Reductions.ec}. \\
\midrule
525--569 &
\ec{signing\_attempt\_verifies},
\ec{signing\_simulation\_verifies},
\ec{signing\_attempt\_simulates}, and
\ec{sign\_internal\_attempt\_relation}. &
Machine-checked facts. &
These lemmas lift \file{HAETAE\_Algebra.ec}'s
\ec{verify\_internal\_sign\_internal} into the rejection relation vocabulary.
The proof obligation is mostly unfolding \ec{valid\_keypair} and applying
formal correctness and norm lemmas. \\
\midrule
571--623 &
Current sample-pair equality and sampler-ok lemmas for structural, bounded,
checked-hyperball, and exact-hyperball samplers. &
Machine-checked facts relying on imported distribution facts. &
The sampler-ok lemmas are adapters: they expose the distribution-ok facts from
\file{HAETAE\_Distributions.ec} in the shape required by
\ec{dsigning\_attempt\_state\_from\_sample\_pair\_sampler}. \\
\midrule
625--747 &
Losslessness and support/sample-ok lemmas for lifted attempt distributions. &
Machine-checked facts with imported distribution rules. &
The proofs use \ec{dlet\_ll}, \ec{dmap\_ll}, \ec{supp\_dlet}, and
\ec{supp\_dmap}.  They show that if the underlying sampler is ok, then the
full attempt distribution is lossless and supported on well-formed bounded
sample pairs. \\
\midrule
749--799 &
Exact-hyperball support equivalence, one-bound, and pullback identity. &
Machine-checked facts; model boundary. &
\ec{dexact\_hyperball\_signing\_attempt\_state\_supportP} expands support
into an existential over coins and a hyperball sample.  This is the key lemma
used repeatedly to turn a probabilistic zero-abort goal into a pointwise
no-abort proof.  It inherits the current formal definition of
\ec{dexact\_hyperball\_signing\_sample\_pair}; it does not prove paper-sampler
faithfulness. \\
\midrule
801--850 &
Reference-rejection abort decomposition and union bound. &
Machine-checked facts with imported probability facts. &
\ec{signing\_attempt\_state\_reference\_rejection\_aborts\_componentE}
rewrites the negation of a conjunction into a disjunction.  The union-bound
lemma uses \ec{mu\_or} and \ec{mu\_bounded}; these are library probability
facts, not HAETAE-specific theorems. \\
\midrule
852--964 &
Pack no-abort for sample-pair states and equality of structural/exact
coin-sample-pair mappings. &
Machine-checked facts. &
Pack acceptance follows from \ec{valid\_signature\_sign\_internal} and
\ec{mode\_crypto\_bytes\_gt0}.  Distribution equalities are functorial
rewrites using \ec{dmap\_dunit}, \ec{dlet\_dunit}, \ec{dmap\_dlet}, and
\ec{dmap\_comp}. \\
\midrule
966--1053 &
\ec{mu\_dlet\_le\_add\_all} and the loss-obligation propagation chain. &
Machine-checked facts conditional on the earlier point/sample loss premise. &
The probabilistic lemma proves that a per-coin distribution inequality with
additive \(\epsilon\) lifts through \ec{dlet}.  The three propagation lemmas
then derive sample-pair loss, coin-sample-pair loss, and full attempt-state
loss.  The initial point-loss premise remains conditional unless separately
proved. \\
\midrule
1055--1270 &
Field-equality lemmas for \ec{signing\_attempt\_state\_of\_coins} and current
field acceptance/no-abort. &
Machine-checked facts. &
The proofs are definitional rewrites showing that every projection of the
attempt state equals the corresponding deterministic field from
\file{HAETAE\_Algebra.ec}.  These lemmas are the backbone of current-signing
transcript validity and current no-abort facts. \\
\midrule
1272--1513 &
Reject1 bound lift, reject2 concrete rewrites, branch-clear proofs,
sample-pair reject1/reject2/pack/hint no-abort, and exact-hyperball
reference-rejection abort-zero chain. &
Machine-checked facts with imported probability facts. &
The pointwise no-abort lemmas rely on \ec{sign\_internal}'s branch byte being
zero and on formal current norm bounds.  The distribution zero lemmas use
\ec{dexact\_hyperball\_signing\_attempt\_state\_supportP}, \ec{mu\_le},
\ec{mu0}, and \ec{mu\_bounded}.  The conclusion is zero abort in the current
formal exact-hyperball attempt model, not a paper-level rejection-rate theorem. \\
\midrule
1515--1672 &
Current-coin reject2 branch clearing, balance identities, reference-rejection
acceptance, verification acceptance, full rejection acceptance, and full
current no-abort. &
Machine-checked facts. &
These lemmas show that attempts generated directly from current signing coins
pass all formal acceptance predicates.  They are used later to prove zero
abort probability for \ec{dsigning\_attempt\_state}. \\
\midrule
1674--1724 &
\ec{dsigning\_attempt\_state} losslessness, entropy-token spread, and current
reference/full rejection abort-zero probabilities. &
Machine-checked facts with imported distribution facts and a distribution
point-bound imported from \file{HAETAE\_Distributions.ec}. &
The token-spread lemma maps the signing-coin token bound through the attempt
state.  This feeds sampler-prequery and ROM bad-event arguments in
\file{HAETAE\_HopGames.ec}. \\
\midrule
1726--1838 &
Simulation relation, transcript validity, transcript relation, record
relation, record-log soundness, and query/transcript list constructors. &
Machine-checked facts. &
These lemmas connect accepted signing attempts to transcript logs.  They are
used by hop games and ROM programming to maintain that logged signing
transcripts are valid and field-consistent. \\
\midrule
1840--1854 &
Nonnegativity facts for rejection and Fiat-Shamir-with-aborts loss terms. &
Machine-checked facts relying on imported reduction-term lemmas. &
\ec{rejection\_sampling\_bound\_obligation\_holds} is checked here, but the
meaning of the loss terms comes from \file{HAETAE\_Reductions.ec}.  It proves
nonnegativity only, not smallness or cryptographic tightness. \\
\midrule
1856--2130 &
Field-equality and no-abort lemmas for
\ec{signing\_attempt\_state\_of\_sample\_pair}. &
Machine-checked facts. &
This is the sample-pair analogue of the current-coins field-equality block.
Because the deterministic signature fields still come from
\ec{sign\_internal}, the proofs again reduce to definitional rewrites and the
formal correctness/norm lemmas from \file{HAETAE\_Algebra.ec}. \\
\midrule
2132--2202 &
Exact-hyperball support implies reference rejection acceptance, full rejection
acceptance, no-abort predicates, and zero full-rejection abort probability. &
Machine-checked facts with model boundary. &
The proof repeatedly expands exact-hyperball support into coins/sample
witnesses and rewrites the state to the sample-pair constructor.  The final
zero-probability theorem is therefore conditional on the current formal
exact-hyperball support definition; replacing that sampler requires replaying
or reestablishing these support-to-no-abort arguments. \\
\bottomrule
\end{longtable}

For downstream use, the theorem boundary is as follows.  Machine-checked:
state constructors, field predicates, support/lossless facts, current-formal
no-abort facts, record-log soundness, and the conditional propagation from
sample-pair loss to attempt-state loss.  Conditional: the initial
structural-to-exact hyperball point-loss obligation unless separately proved.
Imported: probability algebra, list algebra, real summability, and the
nonnegativity of loss terms from \file{HAETAE\_Reductions.ec}.  External or
explanatory: the claim that the exact-hyperball carrier is the paper HAETAE
sampler, the cryptographic interpretation of the loss terms, and any
implementation-level rejection-sampling equivalence.

\paragraph{Rejection obligation sheet.}
The four loss obligations in lines 491--519 are not interchangeable.  The
script proves implications between them, but the first one remains the
entry-point premise unless another file supplies it.  To keep the formulas
readable, write \(D_s^{str}\) for
\ec{dstructural\_signing\_sample\_pair}, \(D_s^{ex}\) for
\ec{dexact\_hyperball\_signing\_sample\_pair}, \(D_c^{str}\) and
\(D_c^{ex}\) for the structural and exact coin-sample-pair distributions,
\(D_a^{str}\) for \ec{dsigning\_attempt\_state}, \(D_a^{ex}\) for
\ec{dexact\_hyperball\_signing\_attempt\_state}, \(S\) for
\ec{signing\_sample\_pair\_of\_coins}, and \(L_{rej}\) for
\ec{rejection\_sampling\_loss\_term}.  With that shorthand, the EasyCrypt
obligations are:

\[
\begin{array}{ll}
O_{\mathrm{point}} &
  \forall md,sk,m,ctx,coins.\; 1 \le
  \mu(D_s^{ex}(md))(\mathrm{pred1}(S(md,sk,m,ctx,coins))) + L_{rej},\\[0.35em]
O_{\mathrm{sample}} &
  \forall md,sk,m,ctx,coins,p.\;
  \mu(D_s^{str}(md,sk,m,ctx,coins))(p) \le
  \mu(D_s^{ex}(md))(p) + L_{rej},\\[0.35em]
O_{\mathrm{coin}} &
  \forall md,sk,m,ctx,p.\;
  \mu(D_c^{str}(md,sk,m,ctx))(p) \le
  \mu(D_c^{ex}(md,sk,m,ctx))(p) + L_{rej},\\[0.35em]
O_{\mathrm{attempt}} &
  \forall md,sk,m,ctx,p.\;
  \mu(D_a^{str}(md,sk,m,ctx))(p) \le
  \mu(D_a^{ex}(md,sk,m,ctx))(p) + L_{rej}.
\end{array}
\]

\begin{longtable}{p{0.23\linewidth}p{0.28\linewidth}p{0.39\linewidth}}
\toprule
\textbf{Obligation} & \textbf{Status in \file{HAETAE\_Rejection.ec}} &
\textbf{Downstream reading} \\
\midrule
\(O_{\mathrm{point}}\) &
Conditional premise. &
This is the carrier-level point-mass statement from which the script can
derive larger loss comparisons.  It must not be reported as proved by
\file{HAETAE\_Rejection.ec}. \\
\midrule
\(O_{\mathrm{sample}}\) &
Machine-checked implication from \(O_{\mathrm{point}}\). &
\ec{structural\_to\_exact\_hyperball\_sample\_pair\_loss\_from\_point\_loss}
uses the structural sample's \ec{dunit} form and monotonicity of \ec{mu}. \\
\midrule
\(O_{\mathrm{coin}}\) &
Machine-checked implication from \(O_{\mathrm{sample}}\). &
\ec{structural\_to\_exact\_hyperball\_coin\_sample\_pair\_loss\_from\_sample\_pair\_loss}
uses \ec{mu\_dlet\_le\_add\_all} to lift the bound through coins. \\
\midrule
\(O_{\mathrm{attempt}}\) &
Machine-checked implication from \(O_{\mathrm{coin}}\). &
\ec{structural\_to\_exact\_hyperball\_attempt\_loss\_from\_coin\_sample\_pair\_loss}
uses distribution equalities that map coin/sample pairs into full attempt
states. \\
\midrule
\ec{rejection\_sampling\_bound\_obligation} &
Machine-checked locally from imported nonnegativity lemmas. &
This proves only nonnegativity of the three named loss terms.  It does not
prove that the losses are small or that the exact-hyperball law is paper
faithful. \\
\bottomrule
\end{longtable}

\paragraph{Carrier-specific theorem boundary.}
\begin{longtable}{p{0.24\linewidth}p{0.31\linewidth}p{0.35\linewidth}}
\toprule
\textbf{Carrier or distribution} & \textbf{Checked theorem surface} &
\textbf{Boundary} \\
\midrule
\ec{signing\_attempt\_state\_of\_coins} &
All current-field equalities, current acceptance, current rejection no-abort,
zero abort probability for \ec{dsigning\_attempt\_state}, and token-spread
bound. &
Machine-checked for the formal current signer and formal
\ec{drandom\_coins}.  Does not by itself prove paper rejection-sampling
statistics. \\
\midrule
\ec{signing\_attempt\_state\_of\_sample\_pair} &
Field equalities, reference rejection acceptance, verification acceptance,
full rejection acceptance, and no-abort predicates for any supplied sample
pair. &
Machine-checked because the deterministic signature fields still come from
\ec{sign\_internal}.  If a future sampler changes those fields, this boundary
must be revisited. \\
\midrule
\ec{dexact\_hyperball\_signing\_attempt\_state} &
Support characterization, pack/reject1/reject2/hint zero-abort probabilities,
reference rejection zero abort, and full rejection zero abort. &
Machine-checked for the current formal exact-hyperball carrier.  The paper
hyperball/sphere sampler correspondence is outside this file unless separately
proved. \\
\midrule
Paper HAETAE sampler &
No direct theorem in this file. &
External or conditional.  The report must not claim that these zero-abort
lemmas certify the paper sampler without a checked replacement/correspondence
argument. \\
\bottomrule
\end{longtable}

\paragraph{Two nontrivial proof blocks.}
For \ec{mu\_dlet\_le\_add\_all}, the mathematical content is an additive-loss
lifting rule: if every branch distribution \(F_1(x)\) is bounded by
\(F_2(x)+\epsilon\), and \(\epsilon\ge 0\), then the \ec{dlet}-mixtures obey
the same additive bound.  The proof expands \ec{dletE}, proves summability of
the weighted series using imported \ec{summable\_mu1\_wght},
\ec{summableD}, and \ec{summableZ}, applies \ec{RealSeries.ler\_sum}, and
uses \ec{weightE} plus \ec{mu\_bounded} for the final real-arithmetic side
condition.

For the abort-zero chain, the pattern is pointwise-to-probabilistic.  The
script first proves that sample-pair states cannot trigger the relevant abort
predicate.  Then, under exact-hyperball support, every sampled state is
rewritten to such a sample-pair state.  Finally \ec{mu\_le} bounds the abort
event by \ec{pred0}, \ec{mu0} rewrites the probability of \ec{pred0} to zero,
and \ec{mu\_bounded}/SMT closes the equality.  This pattern proves zero
probability for the current formal carrier; it is not a paper-level rejection
rate proof.

\paragraph{Lemma-ledger criticism for HAETAE Rejection.}
Even after the line-range companion above, the report plus
\file{HAETAE\_Rejection.ec} still required too much reverse engineering for
proof-obligation reconstruction.  The missing layer was a ledger that says,
for each security-relevant definition or lemma surface, what the checked
statement is, what premises or imported facts it depends on, and how it feeds
\file{HAETAE\_HopGames.ec} or \file{HAETAE\_Security.ec}.  The ledger below
adds that layer.  Pure tuple projections and repeated field-equality siblings
are grouped when they share the same mathematical proof obligation; the row
lists the relevant EasyCrypt names and line ranges so the source remains the
authority for exact syntax and tactic order.  In particular, the ledger names
the token/prequery bridge, the abort-channel predicates, the current-formal
zero-abort theorems, the transcript/log bridge theorems, and the
exact-hyperball support-to-no-abort interface, because these are the surfaces
later games can actually cite.

\paragraph{Definition ledger for HAETAE Rejection.}
{\footnotesize
\begin{longtable}{p{0.24\linewidth}p{0.15\linewidth}p{0.27\linewidth}p{0.24\linewidth}}
\toprule
\textbf{EasyCrypt object(s)} & \textbf{Lines / category} &
\textbf{Statement paraphrase} & \textbf{Premises, imports, downstream role} \\
\midrule
\ec{signing\_transcript\_record}, \ec{signing\_attempt\_sample},
\ec{signing\_attempt\_state} &
16--20. Machine-checked definitions. &
Tuple carriers for logged signing transcripts, sampled vectors, and full
signing-attempt state. &
No premises.  Imported carriers from \file{HAETAE\_Algebra.ec}.  HopGames uses
these as the state space for rejection and transcript-signing games. \\
\midrule
\ec{signing\_attempt\_state\_*} projections,
\ec{signing\_attempt\_sample\_*},
\ec{signature\_rejection\_branch\_byte} &
22--70. Machine-checked definitions. &
Projection vocabulary for coins, sample components, commitment fields,
challenge, response, hint, signature, and reject2 branch bit. &
No premises.  Tuple projection only.  The reject2 branch reads signature field
6; this matters because \ec{sign\_internal} writes zero there. \\
\midrule
\ec{polyvecl\_double},
\ec{signing\_attempt\_state\_reject2\_sampled\_*},
\ec{signing\_attempt\_state\_reject2\_balance\_*} &
73--91. Machine-checked definitions. &
Defines the formal reject2 balance vectors \(2z-y\) and \(2z_{\rm aux}-y_2\). &
Imports formal vector arithmetic from \file{HAETAE\_Algebra.ec}.  HopGames
uses these predicates when charging or erasing reject2 aborts. \\
\midrule
\ec{signing\_attempt\_state\_lowbits\_carrier},
\ec{signing\_attempt\_state\_token},
\ec{signing\_attempt\_state\_programmed\_lowbits},
\ec{signing\_attempt\_state\_lowbits\_consistent} &
93--114. Machine-checked definitions. &
Names the lowbits carrier and signing-token triple, then reconstructs lowbits
by placing the signing entropy token in coefficient zero and requiring equality
with the attempt state's lowbits. &
Depends on \ec{signing\_entropy\_token\_of\_coins}.  Used by transcript and ROM
programming invariants and the token point-bound lemma that bind sampler coins
to challenge input. \\
\midrule
\ec{signing\_attempt\_state\_programmed\_challenge},
\ec{signing\_attempt\_state\_challenge\_consistent} &
116--128. Machine-checked definitions. &
Recomputes the challenge from response auxiliary data, lowbits, and message
hash under the public key. &
Imports \ec{challenge\_hash} and \ec{message\_hash}.  Used to show sampled
attempts match the same challenge predicate as verification. \\
\midrule
\ec{signing\_attempt\_state\_reconstructed\_highbits},
\ec{signing\_attempt\_state\_public\_equation\_consistent} &
130--139. Machine-checked definitions. &
Rebuilds highbits from auxiliary response and public-key challenge term, then
requires equality with the attempt commitment. &
Imports public-key challenge operations from \file{HAETAE\_Algebra.ec}.  Feeds
transcript validity and simulated-signing soundness. \\
\midrule
\ec{signing\_attempt\_state\_signature\_of\_fields},
\ec{signing\_attempt\_state\_fields\_match\_signature},
\ec{signing\_attempt\_state\_field\_accepts} &
141--163. Machine-checked definitions. &
Assembles the signature tuple from attempt fields and combines lowbits,
challenge, public-equation, and signature-field consistency. &
Imports signature layout from \file{HAETAE\_Algebra.ec}.  HopGames uses this
as the field-consistency side condition for accepted sampled transcripts. \\
\midrule
\ec{signing\_attempt\_state\_of\_coins},
\ec{signing\_attempt\_sample\_of\_pair},
\ec{signing\_attempt\_state\_of\_sample\_pair} &
166--199. Machine-checked definitions. &
Constructs current signing attempts from coins, or attempts that replace only
the sampled pair while keeping deterministic signer fields. &
Imports \ec{sign\_internal}, \ec{commitment\_*}, \ec{response\_*}, and
\ec{hint\_vector}.  These constructors are the bridge between real signing,
structural attempts, and exact-hyperball attempts. \\
\midrule
\ec{signing\_attempt\_state\_valid\_signature\_accepts},
\ec{signing\_attempt\_state\_response\_norm\_accepts},
\ec{signing\_attempt\_state\_hint\_norm\_accepts},
\ec{signing\_attempt\_state\_accepts} &
203--218. Machine-checked definitions. &
Defines the non-rejection validity and norm acceptance components. &
Imports \ec{valid\_signature}, \ec{response\_norm\_ok}, and
\ec{verify\_norm\_ok}.  Used by current no-abort and transcript soundness
lemmas. \\
\midrule
\ec{signing\_attempt\_reject1\_bound},
\ec{signing\_attempt\_reject2\_bound},
\ec{signing\_attempt\_reject2\_balance\_norm\_sq},
\ec{signing\_attempt\_reject2\_concrete\_aborts} &
221--238. Machine-checked definitions. &
Defines reject1/reject2 thresholds and the concrete reject2 event
\(branch \land \|2z-y\|^2 < B_0 \ell_n^2\). &
Imports mode bounds from \file{HAETAE\_Params.ec} and norm functions from
\file{HAETAE\_Algebra.ec}.  Used in the rejection-abort decomposition in
HopGames. \\
\midrule
\ec{signing\_attempt\_state\_reject1\_*},
\ec{signing\_attempt\_state\_reject2\_*},
\ec{signing\_attempt\_state\_pack\_*},
\ec{signing\_attempt\_state\_reference\_rejection\_*},
\ec{signing\_attempt\_state\_rejection\_*},
\ec{signing\_attempt\_state\_response\_norm\_aborts},
\ec{signing\_attempt\_state\_hint\_norm\_aborts},
\ec{signing\_attempt\_state\_rejection\_sampling\_aborts},
\ec{signing\_attempt\_state\_aborts},
\ec{signing\_accepts} &
240--343. Machine-checked definitions. &
Names the reject1, reject2, pack, reference-rejection, verification, full
rejection, concrete abort-channel, coarse abort, and signing-acceptance
predicates. &
No cryptographic assumption locally.  These event predicates are the exact
bad-event and loss vocabulary used in HopGames; the report must not treat them
as probability bounds until later \ec{mu} lemmas are applied. \\
\midrule
\ec{dsigning\_attempt\_state},
\ec{dsigning\_attempt\_state\_from\_sample\_pair\_sampler},
\ec{dstructural\_*}, \ec{dbounded\_*}, \ec{dchecked\_*}, \ec{dexact\_*} &
345--419. Machine-checked definitions. &
Defines distributions over attempts by mapping signing coins and sample-pair
distributions into full states. &
Imports \ec{drandom\_coins} and sample-pair distributions from
\file{HAETAE\_Distributions.ec}.  Used by hop probability comparisons and
sampler-prequery arguments. \\
\midrule
\ec{signing\_attempt\_state\_sample\_pair\_wf},
\ec{signing\_attempt\_state\_sample\_pair\_unit},
\ec{signing\_attempt\_state\_sample\_pair\_bounded} &
421--435. Machine-checked definitions. &
Support predicates for sampled \(y_1,y_2\) components. &
Imports \ec{polyvecl\_wf}, \ec{polyveck\_wf}, unit and bounded predicates.
Used by support lemmas and sampler-correctness checks. \\
\midrule
\ec{signing\_attempt\_relation},
\ec{signing\_simulation\_relation},
\ec{signing\_transcript\_relation} &
438--456. Machine-checked definitions. &
Relates a signature to an honest attempt, wraps it as a verified simulation,
and then as a valid transcript. &
Imports \ec{valid\_keypair}, \ec{verify\_internal}, and transcript predicates.
Used by NMA transcript construction and signing-log soundness. \\
\midrule
\ec{signing\_record\_*},
\ec{signing\_record\_log\_sound},
\ec{transcript\_log\_valid} &
459--488. Machine-checked definitions. &
Projects record fields, lists record queries/transcripts, and states log
soundness/validity. &
Imports list \ec{map} and transcript validity.  HopGames consumes these to
preserve transcript logs across signing-oracle games. \\
\midrule
\ec{structural\_to\_exact\_hyperball\_*obligation},
\ec{rejection\_sampling\_bound\_obligation} &
491--523. Conditional premises plus checked definitions. &
States the point, sample-pair, coin-sample-pair, attempt-state loss
obligations and the nonnegativity obligation. &
\(O_{\mathrm{point}}\) is not proved here.  Nonnegativity is proved later from
\file{HAETAE\_Reductions.ec}.  These obligations feed the exact-hyperball
transition in HopGames and the final security bound. \\
\bottomrule
\end{longtable}
}

\paragraph{Lemma ledger for HAETAE Rejection.}
{\footnotesize
\begin{longtable}{p{0.25\linewidth}p{0.14\linewidth}p{0.27\linewidth}p{0.24\linewidth}}
\toprule
\textbf{EasyCrypt lemma(s)} & \textbf{Lines / category} &
\textbf{Statement paraphrase} & \textbf{Premises, imports, downstream role} \\
\midrule
\ec{signing\_attempt\_verifies},
\ec{signing\_simulation\_verifies},
\ec{signing\_attempt\_simulates},
\ec{sign\_internal\_attempt\_relation} &
525--569. Machine-checked facts. &
Honest signing attempts verify, simulation relations verify, attempts lift to
simulation relations, and \ec{sign\_internal} forms an honest attempt. &
Imports \ec{verify\_internal\_sign\_internal},
\ec{valid\_signature\_sign\_internal}, \ec{response\_norm\_ok\_current}, and
\ec{verify\_norm\_ok\_current}.  Used to justify signing-oracle transcripts. \\
\midrule
\ec{signing\_attempt\_sample\_of\_pair\_currentE},
\ec{signing\_attempt\_state\_of\_sample\_pair\_currentE} &
571--591. Machine-checked facts. &
The structural sample pair built from the current coins reconstructs the same
attempt as \ec{signing\_attempt\_state\_of\_coins}. &
Definitional rewrites using \file{HAETAE\_Distributions.ec}.  Used to identify
structural attempts with current signing attempts. \\
\midrule
\ec{structural\_signing\_sample\_pair\_sampler\_ok},
\ec{bounded\_unit\_signing\_sample\_pair\_sampler\_ok},
\ec{checked\_hyperball\_signing\_sample\_pair\_sampler\_ok},
\ec{exact\_hyperball\_signing\_sample\_pair\_sampler\_ok} &
593--623. Machine-checked facts relying on imports. &
Each named sample-pair sampler satisfies the generic sampler-ok interface. &
Imports distribution-ok lemmas from \file{HAETAE\_Distributions.ec}.  Used to
derive attempt-distribution losslessness and support properties. \\
\midrule
\ec{dsigning\_attempt\_state\_from\_sample\_pair\_sampler\_lossless},
\ec{dsigning\_attempt\_state\_from\_sample\_pair\_sampler\_sample\_ok} &
625--676. Machine-checked facts with conditional premise. &
If a supplied sample-pair sampler is ok, the lifted attempt distribution is
lossless and supported on well-formed bounded samples. &
Premise: \ec{signing\_sample\_pair\_sampler\_ok}.  Imports \ec{dlet\_ll},
\ec{dmap\_ll}, \ec{supp\_dlet}, \ec{supp\_dmap}.  Used by all sampler variants
in HopGames. \\
\midrule
\ec{dstructural\_*lossless},
\ec{dbounded\_*lossless},
\ec{dchecked\_*lossless},
\ec{dexact\_*lossless},
\ec{dbounded\_*sample\_ok},
\ec{dchecked\_*sample\_ok},
\ec{dexact\_*sample\_ok} &
678--747. Machine-checked facts with imported sampler facts. &
Instantiates the generic lifting lemmas for structural, bounded-unit,
checked-hyperball, and exact-hyperball attempt distributions. &
Imports sampler-ok lemmas above and \file{HAETAE\_Distributions.ec}.  Used as
side conditions in probabilistic hop proofs. \\
\midrule
\ec{dexact\_hyperball\_signing\_attempt\_state\_supportP},
\ec{dexact\_hyperball\_signing\_attempt\_state\_reference\_rejection\_abort\_bound1},
\ec{dexact\_hyperball\_signing\_attempt\_state\_mu\_pullback} &
749--799. Machine-checked facts with model boundary. &
Characterizes exact-hyperball attempt support, bounds any event by one, and
rewrites the exact-hyperball attempt distribution as the underlying \ec{dlet}. &
Imports \ec{exact\_hyperball\_signing\_sample\_pair\_supportP} and
\ec{mu\_bounded}.  Used by all exact-hyperball abort-zero proofs. \\
\midrule
\ec{signing\_attempt\_state\_reference\_rejection\_aborts\_componentE},
\ec{dexact\_hyperball\_signing\_attempt\_state\_reference\_rejection\_abort\_split},
\ec{dexact\_hyperball\_signing\_attempt\_state\_reference\_rejection\_abort\_union\_bound} &
801--850. Machine-checked facts with imported probability facts. &
Decomposes reference rejection aborts into reject1/reject2/pack/hint events
and upper-bounds their probability by the sum of component probabilities. &
Imports \ec{predU}, \ec{mu\_eq}, \ec{mu\_or}, \ec{mu\_bounded}.  Used when
charging rejection components in hop bounds. \\
\midrule
\ec{signing\_attempt\_state\_of\_sample\_pair\_signatureE},
\ec{signing\_attempt\_state\_of\_sample\_pair\_pack\_accepts},
\ec{signing\_attempt\_state\_of\_sample\_pair\_pack\_no\_abort},
\ec{dexact\_hyperball\_signing\_attempt\_state\_pack\_abort\_zero} &
852--902. Machine-checked facts. &
Sample-pair attempts use the current \ec{sign\_internal} signature, so pack
accepts and exact-hyperball pack abort has probability zero. &
Imports \ec{valid\_signature\_sign\_internal}, \ec{mode\_crypto\_bytes\_gt0},
support characterization, and \ec{mu0}/\ec{mu\_le}.  Used to remove the pack
component of rejection loss. \\
\midrule
\ec{dstructural\_signing\_attempt\_state\_from\_samplerE},
\ec{dstructural\_signing\_attempt\_stateE},
\ec{dstructural\_signing\_attempt\_state\_from\_sampler\_coin\_sample\_pairE},
\ec{dexact\_hyperball\_signing\_attempt\_state\_coin\_sample\_pairE} &
904--964. Machine-checked facts with imported distribution functor laws. &
Equates alternate distribution presentations so loss lemmas can move between
coin/sample-pair and full-state distributions. &
Imports \ec{dmap\_dunit}, \ec{dlet\_dunit}, \ec{dmap\_dlet}, \ec{dmap\_comp}.
Used by the loss-obligation propagation chain. \\
\midrule
\ec{mu\_dlet\_le\_add\_all} &
966--1002. Machine-checked imported-library style fact. &
Lifts a per-branch additive probability inequality through \ec{dlet}. &
Premises: \(\epsilon\ge 0\) and per-branch inequality.  Imports
\ec{RealSeries.ler\_sum}, summability lemmas, \ec{dletE}, \ec{weightE}.  Used
to lift sample-pair loss through signing coins. \\
\midrule
\ec{structural\_to\_exact\_hyperball\_sample\_pair\_loss\_from\_point\_loss},
\ec{structural\_to\_exact\_hyperball\_coin\_sample\_pair\_loss\_from\_sample\_pair\_loss},
\ec{structural\_to\_exact\_hyperball\_attempt\_loss\_from\_coin\_sample\_pair\_loss} &
1004--1053. Machine-checked implications from conditional premises. &
Derives sample-pair, coin-sample-pair, and attempt-state loss obligations from
the preceding obligation in the chain. &
Conditional until \(O_{\mathrm{point}}\) is supplied.  Used by HopGames to
charge the structural-to-exact-hyperball transition. \\
\midrule
\ec{signing\_attempt\_state\_of\_coins\_*E},
\ec{signing\_attempt\_state\_of\_coins\_*consistent},
\ec{signing\_attempt\_state\_of\_coins\_field\_accepts},
\ec{signing\_attempt\_state\_of\_coins\_accepts\_current} &
1055--1270. Machine-checked facts. &
Every projected field of a current-coins attempt equals the corresponding
deterministic signing field; therefore current attempts satisfy field and norm
acceptance predicates. &
Imports \ec{sign\_internal}, \ec{commitment\_*}, \ec{signature\_token},
\ec{valid\_signature\_sign\_internal}, and norm lemmas.  Used for current
signing soundness and zero-abort probabilities. \\
\midrule
\ec{signing\_attempt\_reject1\_bound\_ge\_response\_bound},
\ec{signing\_attempt\_state\_reject1\_accepts\_from\_response\_norm\_ok} &
1272--1287. Machine-checked facts. &
Response-norm acceptance implies reject1 acceptance because the reject1 bound
dominates the response bound. &
Imports mode constants from \file{HAETAE\_Params.ec}.  Used in current and
sample-pair reference-rejection acceptance proofs. \\
\midrule
\ec{signing\_attempt\_state\_balance\_norm\_sq\_concreteE},
\ec{signing\_attempt\_state\_reject2\_under\_bound\_concreteE},
\ec{signing\_attempt\_state\_reject2\_aborts\_concreteE},
\ec{signing\_attempt\_state\_reject2\_accepts\_from\_branch\_clear} &
1289--1330. Machine-checked facts. &
Unfolds reject2 state predicates and proves that clearing the branch bit
implies reject2 acceptance. &
Definitional dependencies only.  Used by current and sample-pair reject2
no-abort lemmas. \\
\midrule
\ec{signing\_attempt\_state\_of\_sample\_pair\_reject2\_branch\_clear},
\ec{signing\_attempt\_state\_of\_sample\_pair\_reject2\_accepts},
\ec{signing\_attempt\_state\_of\_sample\_pair\_reject2\_no\_abort},
\ec{dexact\_hyperball\_signing\_attempt\_state\_reject2\_abort\_zero} &
1332--1382. Machine-checked facts with model boundary. &
Because \ec{sign\_internal} writes zero in the branch byte, sample-pair states
do not reject2-abort; exact-hyperball reject2 abort has probability zero. &
Imports exact-hyperball support and probability zero pattern.  Used to remove
the reject2 component from exact-hyperball rejection loss. \\
\midrule
\ec{dexact\_hyperball\_signing\_attempt\_state\_reference\_rejection\_abort\_reduced\_bound},
\ec{dexact\_hyperball\_signing\_attempt\_state\_reference\_rejection\_abort\_hint\_bound},
\ec{dexact\_hyperball\_signing\_attempt\_state\_reference\_rejection\_abort\_zero} &
1384--1513. Machine-checked facts with model boundary. &
Successively removes reject2, pack, reject1, and hint abort components until
reference rejection abort probability is zero. &
Imports component zero lemmas and \ec{mu\_bounded}.  Used by HopGames when the
exact-hyperball formal carrier is in force. \\
\midrule
\ec{signing\_attempt\_state\_of\_sample\_pair\_response\_norm\_accepts},
\ec{signing\_attempt\_state\_of\_sample\_pair\_reject1\_accepts},
\ec{signing\_attempt\_state\_of\_sample\_pair\_reject1\_no\_abort},
\ec{dexact\_hyperball\_signing\_attempt\_state\_reject1\_abort\_zero},
\ec{signing\_attempt\_state\_of\_sample\_pair\_hint\_norm\_accepts},
\ec{signing\_attempt\_state\_of\_sample\_pair\_hint\_norm\_no\_abort},
\ec{dexact\_hyperball\_signing\_attempt\_state\_hint\_norm\_abort\_zero} &
1401--1500. Machine-checked facts. &
Sample-pair attempts inherit response and hint norm acceptance from current
\ec{sign\_internal}; exact-hyperball reject1 and hint aborts are zero. &
Imports \ec{response\_norm\_ok\_current}, \ec{verify\_norm\_ok\_current}, and
support-to-zero pattern.  Used in the reference-rejection zero chain. \\
\midrule
\ec{signing\_attempt\_state\_of\_coins\_reject2\_*},
\ec{signing\_attempt\_state\_of\_coins\_balance\_norm\_sq\_current},
\ec{signing\_attempt\_state\_of\_coins\_reject2\_accepts\_current} &
1515--1587. Machine-checked facts. &
Current-coins attempts have clear reject2 branch and the expected concrete
balance norm decomposition. &
Imports projection equalities and branch-byte zero fact.  Used by current
reference-rejection acceptance. \\
\midrule
\ec{signing\_attempt\_state\_pack\_accepts\_from\_valid},
\ec{signing\_attempt\_state\_reference\_rejection\_accepts\_from\_accepts},
\ec{signing\_attempt\_state\_of\_coins\_reference\_rejection\_accepts\_current},
\ec{signing\_attempt\_state\_of\_coins\_verification\_accepts\_current},
\ec{signing\_attempt\_state\_of\_coins\_rejection\_accepts\_current},
\ec{signing\_attempt\_state\_of\_coins\_rejection\_no\_abort\_current},
\ec{signing\_attempt\_state\_of\_coins\_no\_abort\_current} &
1589--1672. Machine-checked facts. &
Current-coins attempts pass pack, reference rejection, verification, full
rejection, and coarse abort predicates. &
Imports current field acceptance, norm acceptance, reject2 acceptance, and
\ec{mode\_crypto\_bytes\_gt0}.  Used by current distribution zero-abort
probability lemmas. \\
\midrule
\ec{signing\_attempt\_state\_distribution\_lossless},
\ec{signing\_attempt\_state\_distribution\_token\_spread},
\ec{signing\_attempt\_state\_distribution\_reference\_rejection\_abort\_zero\_current},
\ec{signing\_attempt\_state\_distribution\_rejection\_abort\_zero\_current} &
1674--1724. Machine-checked facts. &
The current attempt distribution is lossless, preserves the signing-token
point bound, and has zero current reference/full rejection abort probability. &
Imports \ec{signing\_coin\_distribution\_lossless},
\ec{signing\_coin\_distribution\_token\_point\_bound}, \ec{dmapE}, and
\ec{mu0}.  Used by sampler prequery and current-signing hop steps. \\
\midrule
\ec{sign\_internal\_simulation\_relation},
\ec{signing\_attempt\_transcript\_valid},
\ec{signing\_attempt\_transcript\_relation},
\ec{sign\_internal\_transcript\_relation},
\ec{sign\_internal\_record\_relation} &
1726--1790. Machine-checked facts. &
Turns honest formal signing into simulation and transcript/record relations. &
Imports \ec{transcript\_of\_signature}, \ec{transcript\_valid}, and
\ec{verify\_internal\_sign\_internal}.  Used by transcript logging games and
NMA lifts. \\
\midrule
\ec{signing\_record\_relation\_sound},
\ec{signing\_record\_log\_sound\_cons},
\ec{signing\_record\_log\_sound\_transcripts\_valid},
\ec{signing\_record\_queries\_cons},
\ec{signing\_record\_transcripts\_cons} &
1792--1838. Machine-checked facts. &
Sound record logs imply simulated signatures and valid transcript lists; cons
lemmas expose list structure. &
Imports list induction and transcript predicates.  Used by HopGames transcript
log invariants and query-budget bookkeeping. \\
\midrule
\ec{rejection\_bound\_parts\_nonnegative},
\ec{rejection\_sampling\_bound\_obligation\_holds} &
1840--1854. Machine-checked facts relying on imports. &
Shows the named rejection/Fiat-Shamir-with-aborts loss terms are nonnegative. &
Imports nonnegativity lemmas from \file{HAETAE\_Reductions.ec}.  This proves
nonnegativity only; smallness and cryptographic interpretation remain outside
this file. \\
\midrule
\ec{signing\_attempt\_state\_of\_sample\_pair\_*E},
\ec{signing\_attempt\_state\_of\_sample\_pair\_*consistent},
\ec{signing\_attempt\_state\_of\_sample\_pair\_field\_accepts},
\ec{signing\_attempt\_state\_of\_sample\_pair\_valid\_signature\_accepts},
\ec{signing\_attempt\_state\_of\_sample\_pair\_accepts\_current} &
1856--2073. Machine-checked facts. &
Sample-pair attempts have the same deterministic signature fields as the
current signer and therefore satisfy field, validity, response, and hint
acceptance. &
Imports projection equalities, \ec{sign\_internal}, formal correctness, and
norm lemmas.  Used by exact-hyperball support-to-acceptance proofs. \\
\midrule
\ec{signing\_attempt\_state\_of\_sample\_pair\_verification\_accepts\_current},
\ec{signing\_attempt\_state\_of\_sample\_pair\_reference\_rejection\_accepts},
\ec{signing\_attempt\_state\_of\_sample\_pair\_reference\_rejection\_no\_abort},
\ec{signing\_attempt\_state\_of\_sample\_pair\_rejection\_accepts\_current},
\ec{signing\_attempt\_state\_of\_sample\_pair\_rejection\_no\_abort\_current} &
2075--2130. Machine-checked facts. &
Sample-pair states pass verification, reference rejection, full rejection, and
the corresponding no-abort predicates. &
Imports field acceptance, reference-rejection acceptance, and deterministic
signer facts.  Used directly by exact-hyperball no-abort lemmas. \\
\midrule
\ec{dexact\_hyperball\_signing\_attempt\_state\_reference\_rejection\_accepts},
\ec{dexact\_hyperball\_signing\_attempt\_state\_reference\_rejection\_no\_abort},
\ec{dexact\_hyperball\_signing\_attempt\_state\_rejection\_accepts},
\ec{dexact\_hyperball\_signing\_attempt\_state\_rejection\_no\_abort},
\ec{dexact\_hyperball\_signing\_attempt\_state\_rejection\_abort\_zero} &
2132--2202. Machine-checked facts with model boundary. &
Any state in the current exact-hyperball attempt distribution passes reference
and full rejection; full rejection abort probability is zero. &
Imports support characterization and sample-pair acceptance lemmas.  Used by
exact-hyperball hop games.  This remains a formal-carrier theorem, not a
paper-sampler correspondence theorem. \\
\bottomrule
\end{longtable}
}

\section{Script companion coverage for HAETAE ROM}

\file{HAETAE\_ROM.ec} instantiates the generic random-oracle library with
HAETAE-specific query and output types.  This file is small but central:
all ROM programming, counted-query, and sampler-prequery events refer to its
constructors.

\begin{longtable}{p{0.27\linewidth}p{0.38\linewidth}p{0.25\linewidth}}
\toprule
\textbf{Script object family} & \textbf{Mathematical reading} & \textbf{Downstream role} \\
\midrule
\ec{ro\_query} constructors &
\ec{MessageHashQuery}, \ec{ChallengeHashQuery}, \ec{MatrixExpandQuery}, and
\ec{SamplerExpandQuery} partition random-oracle inputs by protocol role. &
Bad-event flags and query-budget predicates count only the relevant
constructor families, especially challenge and sampler queries. \\
\midrule
\ec{ro\_output} tuple &
The random-oracle output is a tuple \(\mathsf{crh} \times
\mathsf{challenge} \times \mathsf{matrix} \times \mathsf{random\_coins}\).
Projection functions select the component appropriate to the query. &
\file{HAETAE\_Scheme.ec} uses \ec{ro\_message\_hash},
\ec{ro\_challenge\_hash}, \ec{ro\_matrix}, and \ec{ro\_signing\_coins}. \\
\midrule
\ec{dro\_output} &
The output distribution dispatches by query kind: message hash samples
\ec{dcrh}, challenge hash samples \ec{dchallenge}, matrix expansion samples
\ec{dmatrix}, and sampler expansion samples \ec{drandom\_coins}. &
The \ec{FullRO} clone uses this as the oracle's output law.  Losslessness is
needed for game proofs that call the oracle. \\
\midrule
\ec{clone import FullRO as HAETAE\_RO} &
The generic full random oracle is instantiated with HAETAE query/output types
and \ec{dro\_output}. &
HopGames and ROM-programming modules reason about \ec{HAETAE\_RO.FRO.m} and
oracle methods through this clone. \\
\midrule
\ec{sampler\_expand\_query\_dro\_output\_ro\_signing\_coins} &
For sampler-expansion queries, projecting \ec{ro\_signing\_coins} from
\ec{dro\_output} has the same law as \ec{drandom\_coins}. &
Used by sampler-freshness lemmas in \file{HAETAE\_HopGames.ec} to replace a
fresh ROM sampler query by the signing-coin distribution. \\
\bottomrule
\end{longtable}

Boundary: this is a random-oracle interface definition, not a theorem that any
concrete SHAKE invocation behaves randomly.  That modeling step remains the
ROM assumption boundary.

\subsection{Lemma-ledger companion for HAETAE ROM}

\paragraph{Lemma-ledger criticism for HAETAE ROM.}
The previous companion coverage described the ROM interface, but it was not yet
sufficient for a reader to reconstruct the proof obligations attached to the
query dictionary.  It did not identify which names are query constructors,
which names are output projections, which facts are distribution losslessness
lemmas, and which boundary is only the external interpretation of the formal
oracle as a random oracle for SHAKE-like calls.  This matters because
\file{HAETAE\_ROM\_Programming.ec}, \file{HAETAE\_Transcript.ec}, and
\file{HAETAE\_HopGames.ec} all count or program particular constructors, and a
single confused constructor changes the bad-event statement.  The ledger below
keeps the formal \ec{FullRO} instantiation separate from the external ROM
modeling assumption.

\paragraph{Imported-interface and assumption ledger for HAETAE ROM.}
{\footnotesize
\begin{longtable}{p{0.24\linewidth}p{0.15\linewidth}p{0.28\linewidth}p{0.23\linewidth}}
\toprule
\textbf{EasyCrypt object} & \textbf{Lines / category} &
\textbf{Mathematical statement paraphrase} &
\textbf{Local premises, imports, downstream role} \\
\midrule
\ec{AllCore}, \ec{Distr}, \ec{List}, \ec{PROM} &
1. Imported/library fact. &
Imports the core logic, probability distributions, lists, and generic
programmable random-oracle library. &
No HAETAE-specific premise.  The \ec{PROM.FullRO} interface is the formal
oracle machinery later used by HopGames and ROM programming. \\
\midrule
\file{HAETAE\_Params.ec}, \file{HAETAE\_Algebra.ec},
\file{HAETAE\_Distributions.ec} &
2, 6--8. Imported/library fact. &
Imports formal mode values, carrier types, algebraic challenge/matrix
constructors, and output distributions. &
Used to type queries and outputs and to prove \ec{dro\_output} lossless.  These
are formal-model imports, not SHAKE security. \\
\midrule
\ec{clone import FullRO as HAETAE\_RO} &
80--84. Imported/library fact. &
Instantiates the generic full random oracle with HAETAE query and output types
and \ec{dro\_output}. &
The clone supplies \ec{HAETAE\_RO.FRO} and its state/method vocabulary used by
\file{HAETAE\_ROM\_Programming.ec} and \file{HAETAE\_HopGames.ec}. \\
\midrule
Concrete SHAKE/FIPS-to-ROM interpretation &
1--111. External cryptographic assumption. &
The file does not prove that SHAKE or any implementation behaves as this
formal random oracle. &
That modeling step remains outside \file{HAETAE\_ROM.ec}; later theorems are
ROM theorems unless a separate standard-model hash argument is supplied. \\
\bottomrule
\end{longtable}
}

\paragraph{Definition ledger for HAETAE ROM.}
{\footnotesize
\begin{longtable}{p{0.25\linewidth}p{0.14\linewidth}p{0.28\linewidth}p{0.23\linewidth}}
\toprule
\textbf{EasyCrypt object(s)} & \textbf{Lines / category} &
\textbf{Mathematical statement paraphrase} &
\textbf{Local premises, imports, downstream role} \\
\midrule
\ec{ro\_query}, \ec{MessageHashQuery}, \ec{ChallengeHashQuery},
\ec{MatrixExpandQuery}, \ec{SamplerExpandQuery} &
10--14. Machine-checked fact. &
Defines the four disjoint formal ROM input classes: message hash, challenge
hash, matrix expansion, and sampler expansion. &
Used by \file{HAETAE\_Transcript.ec} challenge-site definitions,
\file{HAETAE\_ROM\_Programming.ec} counted-query predicates, and HopGames
sampler/challenge bad events. \\
\midrule
\ec{ro\_output} &
15. Machine-checked fact. &
Defines a ROM output as a tuple of CRH bytes, challenge, matrix, and signing
coins. &
Projection discipline is essential: each query class consumes only the
component relevant to that query. \\
\midrule
\ec{message\_hash\_query}, \ec{challenge\_hash\_query},
\ec{matrix\_expand\_query}, \ec{sampler\_expand\_query} &
17--24. Machine-checked fact. &
Names the canonical constructors used by the scheme and transcript scripts. &
Downstream proofs count challenge queries and sampler queries by recognizing
these constructors; they are not interchangeable. \\
\midrule
\ec{ro\_crh\_zero}, \ec{ro\_challenge\_zero}, \ec{ro\_matrix\_zero},
\ec{ro\_random\_coins\_zero}, \ec{ro\_output\_zero},
\ec{ro\_output\_to\_*}, \ec{ro\_output\_of\_*} &
27--45. Machine-checked fact. &
Defines zero/default output components, tuple projections, and embeddings of a
single component into the full output tuple. &
Used to state and prove output-distribution laws.  These are formal tuple
operations, not hash-domain-separation proofs. \\
\midrule
\ec{dmatrix\_for\_query}, \ec{dro\_output} &
48--57. Machine-checked fact. &
Defines the distribution sampled for each query class and embeds it into the
appropriate output component. &
Imports \ec{dcrh}, \ec{dchallenge}, \ec{dmatrix}, and
\ec{drandom\_coins}.  This is the output law passed to \ec{FullRO}. \\
\midrule
\ec{Oracle}, \ec{POracle} &
85--91. Imported/library fact. &
Re-exports full and public oracle module-type interfaces from \ec{FRO}. &
Used by \file{HAETAE\_Scheme.ec}, \file{HAETAE\_ROM\_Programming.ec}, and
HopGames modules to state what oracle methods are available. \\
\midrule
\ec{ro\_message\_hash}, \ec{ro\_challenge\_hash}, \ec{ro\_matrix},
\ec{ro\_signing\_coins} &
93--97. Machine-checked fact. &
Names the output projections used after a ROM call. &
\file{HAETAE\_Scheme.ec} uses these selectors; HopGames relies on their
projection equations when replacing sampler ROM calls by distribution facts. \\
\bottomrule
\end{longtable}
}

\paragraph{Lemma ledger for HAETAE ROM.}
{\footnotesize
\begin{longtable}{p{0.25\linewidth}p{0.14\linewidth}p{0.28\linewidth}p{0.23\linewidth}}
\toprule
\textbf{EasyCrypt lemma(s)} & \textbf{Lines / category} &
\textbf{Mathematical statement paraphrase} &
\textbf{Local premises, imports, downstream role} \\
\midrule
\ec{matrix\_query\_distribution\_lossless} &
59--65. Machine-checked fact. &
The matrix distribution selected by a \((mode, seed)\) query is lossless. &
Imports \ec{matrix\_distribution\_lossless}.  Used inside the full output
losslessness proof. \\
\midrule
\ec{ro\_output\_distribution\_lossless} &
67--77. Machine-checked fact. &
For every ROM query, \ec{dro\_output q} is lossless. &
Proof splits by query constructor and imports distribution losslessness for
CRH, challenge, matrix, and signing coins.  Required by the \ec{FullRO} clone
and all oracle-call game proofs. \\
\midrule
\ec{sampler\_expand\_query\_dro\_output\_ro\_signing\_coins} &
99--107. Machine-checked fact. &
Projecting signing coins from the output law of a sampler-expansion query has
exactly distribution \ec{drandom\_coins}. &
Used by sampler-freshness and sampler-prequery lemmas in HopGames and ROM
programming; it is a formal marginal law, not a SHAKE proof. \\
\bottomrule
\end{longtable}
}

\section{Script companion coverage for HAETAE Scheme}

\file{HAETAE\_Scheme.ec} clones the generic signature-game interface and
defines the scheme procedures used by \file{HAETAE\_Security.ec}.  It is the
interface contract between the algebraic model and generic EUF-CMA/NMA games.

\begin{longtable}{p{0.27\linewidth}p{0.38\linewidth}p{0.25\linewidth}}
\toprule
\textbf{Script object family} & \textbf{Mathematical reading} & \textbf{Downstream role} \\
\midrule
\ec{clone import Sig\_ROM as SIG} &
Instantiates generic signature-game types with HAETAE carriers and the HAETAE
random-oracle output distribution. &
All \ec{SIG.EUF\_CMA}, \ec{SIG.UF\_NMA}, adversary, and scheme module types
used in Security and HopGames come from this clone. \\
\midrule
\ec{haetae\_mode} &
An abstract operation selecting the parameter mode for the scheme instance. &
Final theorems are for the formal scheme at this mode; later premises such as
\ec{fs\_with\_aborts\_interfaces\_ready haetae\_mode} are mode-indexed. \\
\midrule
\ec{HAETAE(H).kg} &
Samples \ec{sd} from \ec{dseed}, derives \ec{rhoprime}, queries the ROM at
\ec{matrix\_expand\_query haetae\_mode rhoprime}, and returns
\ec{keygen\_internal haetae\_mode rhoprime}. &
Defines the key-generation behavior used by correctness and all security
games. \\
\midrule
\ec{HAETAE(H).sign} &
Samples \ec{drandom\_coins}, queries \ec{sampler\_expand\_query}, derives the
public key from the secret key, queries message hash and challenge hash, and
returns \ec{sign\_internal}. &
Fixes the query order and query kinds used by ROM programming and bad-event
proofs.  The signer ignores the challenge output except through the formal
\ec{sign\_internal} model. \\
\midrule
\ec{HAETAE(H).verify} &
Returns \ec{verify\_internal haetae\_mode pk m ctx sig}. &
Correctness and EUF-CMA success events use this predicate. \\
\bottomrule
\end{longtable}

Boundary: this file exposes the formal scheme to generic games.  It does not
prove security; it also does not prove byte-level implementation equivalence.
The checked facts about validity and correctness are imported from
\file{HAETAE\_Algebra.ec} and consumed later in \file{HAETAE\_Security.ec}.

\subsection{Lemma-ledger companion for HAETAE Scheme}

\paragraph{Lemma-ledger criticism for HAETAE Scheme.}
The earlier report explained the scheme at the object-family level, but it did
not expose enough proof-script detail to reconstruct the interface boundary.
The scheme file contains no lemmas; its proof significance is in the exact
clone instantiation and the exact program text of \ec{kg}, \ec{sign}, and
\ec{verify}.  In particular, the signing procedure fixes the ROM query order
that later bad-event arguments rely on.  It samples seed coins, queries
\ec{sampler\_expand\_query}, replaces the coins by \ec{ro\_signing\_coins},
then queries message hash and challenge hash before returning
\ec{sign\_internal}.  The ledger below records that wiring without treating it
as a proof of security, implementation equivalence, or SHAKE correctness.

\paragraph{Imported-interface and assumption ledger for HAETAE Scheme.}
{\footnotesize
\begin{longtable}{p{0.24\linewidth}p{0.15\linewidth}p{0.28\linewidth}p{0.23\linewidth}}
\toprule
\textbf{EasyCrypt object} & \textbf{Lines / category} &
\textbf{Mathematical statement paraphrase} &
\textbf{Local premises, imports, downstream role} \\
\midrule
\ec{AllCore}, \file{Sig\_ROM.ec}, \file{HAETAE\_Params.ec},
\file{HAETAE\_Algebra.ec}, \file{HAETAE\_Distributions.ec},
\file{HAETAE\_ROM.ec} &
1--10. Imported/library fact. &
Imports the generic signature-game skeleton plus the HAETAE formal types,
algorithms, distributions, and ROM dictionary. &
No local cryptographic premise.  The file only wires those imported surfaces
into one scheme module. \\
\midrule
\ec{clone import Sig\_ROM as SIG} &
12--20. Imported/library fact. &
Instantiates the abstract signature-game interface with HAETAE public keys,
secret keys, messages, contexts, signatures, ROM queries, ROM outputs, and
\ec{dro\_output}. &
All \ec{SIG.Correctness}, \ec{SIG.EUF\_CMA}, \ec{SIG.UF\_NMA}, adversary, and
oracle module types in HopGames and Security refer to this clone. \\
\midrule
Security theorem assumptions introduced locally &
1--66. External cryptographic assumption. &
None are introduced by \file{HAETAE\_Scheme.ec}. &
The file defines the formal scheme interface.  Hardness assumptions, ROM
modeling, sampler correspondence, and implementation equivalence enter
elsewhere or remain external. \\
\midrule
Direct theorem boundary &
1--66. Explanatory prose. &
\file{HAETAE\_Scheme.ec} contains no local lemmas or theorems. &
All theorem-level reasoning over \ec{HAETAE} occurs in \file{HAETAE\_Security.ec}
and \file{HAETAE\_HopGames.ec}, usually after inlining these procedures. \\
\bottomrule
\end{longtable}
}

\paragraph{Procedure ledger for HAETAE Scheme.}
{\footnotesize
\begin{longtable}{p{0.25\linewidth}p{0.14\linewidth}p{0.28\linewidth}p{0.23\linewidth}}
\toprule
\textbf{EasyCrypt object(s)} & \textbf{Lines / category} &
\textbf{Mathematical statement paraphrase} &
\textbf{Local premises, imports, downstream role} \\
\midrule
\ec{haetae\_mode} &
22. Conditional premise. &
An abstract formal mode parameter selects which parameter table the scheme
uses. &
All top-level bounds are relative to this mode.  Later obligations such as
\ec{fs\_with\_aborts\_interfaces\_ready haetae\_mode} are mode-indexed. \\
\midrule
\ec{module HAETAE(H : SIG.POracle)} &
24--65. Machine-checked fact. &
Defines a formal signature scheme parameterized by a public random-oracle
interface. &
This is the concrete scheme argument to \ec{SIG.Correctness},
\ec{SIG.EUF\_CMA}, and \ec{SIG.UF\_NMA}. \\
\midrule
\ec{HAETAE(H).kg} &
25--34. Machine-checked fact. &
Samples \ec{sd} from \ec{dseed}, derives \ec{rhoprime}, queries
\ec{matrix\_expand\_query haetae\_mode rhoprime}, and returns
\ec{keygen\_internal haetae\_mode rhoprime}. &
Imports seed distribution, keygen seed derivation, ROM query constructor, and
\ec{keygen\_internal}.  Used in correctness and game-equivalence proofs after
inlining. \\
\midrule
\ec{HAETAE(H).sign} &
36--57. Machine-checked fact. &
Samples \ec{drandom\_coins}, queries the sampler ROM site, replaces coins by
\ec{ro\_signing\_coins}, derives the public key, queries message hash and
challenge hash, then returns \ec{sign\_internal}. &
This fixes the formal query order and active-signing surface used by
\file{HAETAE\_ROM\_Programming.ec} and \file{HAETAE\_HopGames.ec}.  It is not
a proof that a concrete implementation makes the same calls. \\
\midrule
\ec{HAETAE(H).verify} &
59--64. Machine-checked fact. &
Verification is exactly \ec{verify\_internal haetae\_mode pk m ctx sig}. &
Correctness and EUF-CMA games use this verifier.  Algebra proves
\ec{verify\_internal\_sign\_internal}; this file only exposes the call. \\
\bottomrule
\end{longtable}
}

\section{Script companion coverage for HAETAE Events and hash support}

Three smaller manifest files complete the next dependency layer:
\file{HAETAE\_Events.ec}, \file{HAETAE\_FIPS202.ec}, and
\file{HAETAE\_Keccak1600.ec}.  They are not final security files, but they
matter because the proof imports their definitions and facts.

\begin{longtable}{p{0.20\linewidth}p{0.43\linewidth}p{0.27\linewidth}}
\toprule
\textbf{File} & \textbf{Definitions and checked lemma families} & \textbf{Boundary and downstream role} \\
\midrule
\file{HAETAE\_Events.ec} &
Defines \ec{keygen\_rejection\_failure},
\ec{signing\_rejection\_failure},
\ec{verification\_rejection\_failure}, and
\ec{any\_rejection\_failure}.  Lemmas such as
\ec{keygen\_rejection\_free\_current},
\ec{signing\_rejection\_free\_current},
\ec{verification\_rejection\_free\_current}, and
\ec{accepted\_transcript\_rejection\_free\_current} prove that the current
formal signer avoids these event predicates. &
Machine-checked event vocabulary.  Used by rejection reductions and by the
security proof's rejection-free branches.  It does not estimate rejection
probabilities for a separate implementation. \\
\midrule
\file{HAETAE\_FIPS202.ec} &
Defines byte-list \ec{fips202\_state}, SHAKE256 constants, the one-block
absorb/padding model, \ec{shake256\_squeeze}, and \ec{shake256}.  Lemmas prove
constant equalities and output-size facts. &
Imported by \file{HAETAE\_Algebra.ec} to model seed derivation.  It is checked
hash plumbing, not a ROM proof and not a full FIPS validation certificate. \\
\midrule
\file{HAETAE\_Keccak1600.ec} &
Defines lanes, byte/lane conversion, rho offsets, round constants, theta,
rho-pi, chi, iota, \ec{keccak\_round}, \ec{keccak\_f1600\_lanes}, and
\ec{keccak\_f1600\_bytes}.  Lemmas prove sizes, bit-level unfoldings, and
round-function projection equations. &
Imported by \file{HAETAE\_FIPS202.ec}.  These are machine-checked definitions
and local consistency lemmas for the formal byte-list permutation.  They do not
turn the final theorem into a standard-model SHAKE theorem. \\
\bottomrule
\end{longtable}

For a reimplementation, these support files should be handled before claiming
the algebra layer is reproduced.  Their imported facts are dependencies of the
formal model, even though the top-level security theorem remains a
random-oracle-model theorem.

\subsection{Lemma-ledger companion for HAETAE Events}

\paragraph{Lemma-ledger criticism for HAETAE Events.}
The previous report put \file{HAETAE\_Events.ec} in a combined support-file
table, which was useful for orientation but too weak for reconstructing the
event boundary.  The file is small, but it names the rejection-failure
predicates and proves that the current formal algorithms avoid them.  Those
facts must not be mistaken for a probability estimate for a paper rejection
loop or for an implementation-level no-abort theorem.  They are checked
boolean facts about \ec{keygen\_internal}, \ec{sign\_internal}, and the
formal norm predicates imported from \file{HAETAE\_Algebra.ec}.  Downstream,
\file{HAETAE\_Rejection.ec} imports this vocabulary directly; HopGames and
Security see it through the rejection layer rather than by expanding these
definitions.

\paragraph{Imported-interface and assumption ledger for HAETAE Events.}
{\footnotesize
\begin{longtable}{p{0.24\linewidth}p{0.15\linewidth}p{0.28\linewidth}p{0.23\linewidth}}
\toprule
\textbf{EasyCrypt object} & \textbf{Lines / category} &
\textbf{Mathematical statement paraphrase} &
\textbf{Local premises, imports, downstream role} \\
\midrule
\ec{AllCore}, \file{HAETAE\_Params.ec}, \file{HAETAE\_Algebra.ec},
\ec{HAETAE\_Params}, \ec{HAETAE\_Algebra} &
1--7. Imported/library fact. &
Imports the formal parameter, key/signature validity, norm, keygen, and signer
definitions used to state event predicates. &
\file{HAETAE\_Rejection.ec} imports \file{HAETAE\_Events.ec} directly; HopGames
and Security consume the resulting event vocabulary transitively. \\
\midrule
External rejection-sampling estimates &
1--57. External cryptographic assumption. &
No rejection probability or paper rejection-loop estimate is introduced in
this file. &
Such estimates are handled by \file{HAETAE\_Rejection.ec},
\file{HAETAE\_HopGames.ec}, \file{HAETAE\_Security.ec}, or remain external
conditional premises. \\
\bottomrule
\end{longtable}
}

\paragraph{Definition and lemma ledger for HAETAE Events.}
{\footnotesize
\begin{longtable}{p{0.25\linewidth}p{0.14\linewidth}p{0.28\linewidth}p{0.23\linewidth}}
\toprule
\textbf{EasyCrypt object(s)} & \textbf{Lines / category} &
\textbf{Mathematical statement paraphrase} &
\textbf{Local premises, imports, downstream role} \\
\midrule
\ec{keygen\_rejection\_failure} &
9--10. Machine-checked fact. &
Keygen failure is the disjunction of invalid keypair or failed secret-norm
predicate. &
Imports \ec{valid\_keypair} and \ec{secret\_norm\_ok}.  Used as named event
vocabulary for the rejection layer. \\
\midrule
\ec{signing\_rejection\_failure} &
12--13. Machine-checked fact. &
Signing failure is the disjunction of invalid signature or failed response
norm predicate. &
Imports \ec{valid\_signature} and \ec{response\_norm\_ok}.  Downstream
rejection proofs can refer to this event without restating the predicate. \\
\midrule
\ec{verification\_rejection\_failure} &
15--16. Machine-checked fact. &
Verification failure is the disjunction of invalid signature or failed
verification-norm predicate. &
Imports \ec{valid\_signature} and \ec{verify\_norm\_ok}.  Used only as formal
event vocabulary, not as a probability bound. \\
\midrule
\ec{any\_rejection\_failure} &
18--22. Machine-checked fact. &
Combines keygen, signing, and verification rejection failures by disjunction. &
This is the strongest combined event predicate exposed locally; later proofs
generally reason through \file{HAETAE\_Rejection.ec}. \\
\midrule
\ec{keygen\_rejection\_free\_current} &
24--30. Machine-checked fact. &
The current formal \ec{keygen\_internal} output does not satisfy
\ec{keygen\_rejection\_failure}. &
Proof unfolds current keygen, key validity, secret norm, and mode cases.  This
is formal-model no-failure, not an implementation theorem. \\
\midrule
\ec{signing\_rejection\_free\_current} &
32--38. Machine-checked fact. &
The current formal \ec{sign\_internal} output does not satisfy
\ec{signing\_rejection\_failure}. &
Imports \ec{valid\_signature\_sign\_internal} and
\ec{response\_norm\_ok\_current}.  Used to justify clean current-signing
branches. \\
\midrule
\ec{verification\_rejection\_free\_current} &
40--46. Machine-checked fact. &
The current formal \ec{sign\_internal} output does not satisfy
\ec{verification\_rejection\_failure}. &
Imports \ec{valid\_signature\_sign\_internal} and
\ec{verify\_norm\_ok\_current}.  Used as the verifier-side clean-event fact. \\
\midrule
\ec{accepted\_transcript\_rejection\_free\_current} &
48--55. Machine-checked fact. &
The keygen-plus-sign transcript produced by the current formal algorithms is
outside \ec{any\_rejection\_failure}. &
Depends on the three component rejection-free lemmas.  It is the local
combined theorem boundary for events; top-level security still relies on later
hop/reduction arguments and hardness premises. \\
\bottomrule
\end{longtable}
}

\section{Evidence required for a standalone claim}

If a future version of the documentation is advertised as sufficient by itself,
it must contain enough evidence for an independent reader to rebuild the proof
without opening the repository.  At minimum, it must satisfy the following
requirements.

\begin{longtable}{p{0.30\linewidth}p{0.42\linewidth}p{0.18\linewidth}}
\toprule
\textbf{Required evidence} & \textbf{Why it is necessary} & \textbf{Present here?} \\
\midrule
Literal EasyCrypt source for every manifest file & EasyCrypt checks exact
syntax, definitions, imports, and proof scripts, not prose descriptions. & No \\
Exact module restriction sets & Security games rely on separation between
adversaries, oracles, simulators, and hidden game state. & Representative only \\
Full helper-lemma statements and proofs & Top-level theorems depend on local
facts whose names alone do not determine their statements or proof obligations.
& No \\
Complete import and library version contract & Proof success can depend on
library names, theories, notations, and solver behavior. & Partial \\
Traceability matrix from prose to checked declarations & A reader needs to know
which checked object discharges each mathematical step. & Required, not
included in full \\
Transcript of a successful verification run & The claim is machine-checked only
after the checker has accepted the complete source tree. & Build command only \\
\bottomrule
\end{longtable}

This table is an acceptance gate.  The current document intentionally leaves
several entries as ``No'' or ``Partial'', so it must not be cited as a
self-contained proof archive.

\section{What this report deliberately does not claim}

The checked development is precise, but its scope is narrow.  Reimplementing
the proof from this report does not prove any of the following.

\begin{itemize}
\item A standard-model theorem replacing the random oracle by SHAKE.
\item Concrete MLWE or Module-SIS bit-security estimates for HAETAE parameter
sets.
\item Constant-time implementation security.
\item Byte-level parsing correctness for all concrete encodings.
\item A machine-checked equivalence theorem
\[
  \Spec \equiv \ECModel .
\]
\item A cost-semantics theorem bounding the running time of the reductions.
\end{itemize}

Those are boundary obligations.  The report explains where each boundary enters
so that a reimplementation does not silently promote the checked theorem into a
stronger end-to-end certification claim.

\section{Repository paths}

The proof gate lives under the project directory \file{haetae-security}.  The
active proof-management surface is:

\begin{lstlisting}[language=bash]
provable-security/
  proof-files.txt
  verify-provable-security.sh
  easycrypt/
  logs/
\end{lstlisting}

The manifest source files also exist at the project root.  The current proof
gate compiles the copies under \file{provable-security/easycrypt/}.  A
reimplementation may keep only one copy, but then the verification script and
import paths must be adjusted consistently.

\chapter{Verification Workflow}

\section{Prerequisites}

Install an EasyCrypt version capable of compiling the current scripts and make
sure the command is visible in the shell:

\begin{lstlisting}[language=bash]
easycrypt --version
\end{lstlisting}

The proof also imports supporting EasyCrypt libraries from
\file{kyber-security}.  A fresh reimplementation should preserve this include
path until all dependencies have been audited and replaced.

\section{The manifest}

The manifest is the list of files that define the proof gate.  It should contain
exactly these entries, in dependency-compatible order:

\begin{lstlisting}
HAETAE_Security.ec
Sig_ROM.ec
HAETAE_HopGames.ec
HAETAE_Reductions.ec
HAETAE_ROM.ec
HAETAE_ROM_Programming.ec
HAETAE_Scheme.ec
HAETAE_Transcript.ec
HAETAE_Events.ec
HAETAE_Assumptions.ec
HAETAE_Distributions.ec
HAETAE_Rejection.ec
HAETAE_Algebra.ec
HAETAE_Params.ec
HAETAE_FIPS202.ec
HAETAE_Keccak1600.ec
\end{lstlisting}

Generated known-answer-test files, archival experiments, and unrelated
implementation validation artifacts should not be added to this gate merely
because they compile.  The manifest is a proof-surface boundary, not a list of
all EasyCrypt files in the repository.

\section{Running the proof gate}

From the project root, run:

\begin{lstlisting}[language=bash]
sh provable-security/verify-provable-security.sh
\end{lstlisting}

The script compiles each manifest entry with include paths equivalent to:

\begin{lstlisting}[language=bash]
easycrypt compile provable-security/easycrypt/<file>.ec \
  -I provable-security/easycrypt \
  -I kyber-security
\end{lstlisting}

The run is successful only if the final line is:

\begin{lstlisting}
RESULT: PASS
\end{lstlisting}

\section{Independent audit commands}

A reimplementation should also provide small audit commands.  The following
checks are not substitutes for EasyCrypt, but they catch common proof-surface
drift.

\begin{lstlisting}[language=bash]
# Count declarations and confirm that the inventory is stable.
./dissertation/etc/count-easycrypt-declarations.sh

# Confirm that no local proof file uses unchecked proof holes.
rg -n "\b(axiom|admit\.|sorry\.)\b" provable-security/easycrypt *.ec

# Confirm that manifest copies match root copies, if both are retained.
rg -v "^#|^$" provable-security/proof-files.txt |
while read f; do
  cmp -s "$f" "provable-security/easycrypt/$f" || echo "$f differs"
done
\end{lstlisting}

For the current artifact, the declaration inventory is:

\begin{lstlisting}
TOTAL,22046,63,678,113,1088,0,105
\end{lstlisting}

The columns are file lines, type declarations, operation declarations, module
declarations, lemma declarations, axiom declarations, and Hoare/equivalence
judgments.

\chapter{Mathematical Theorem Boundary}

\section{Security game}

The informal target is the random-oracle EUF-CMA experiment.  An adversary
\(A\) receives a public key, random-oracle access, and signing-oracle access.
It outputs \((m^\star,ctx^\star,\sigma^\star)\).  It wins when
\[
  \mathsf{Verify}^{H}(pk,m^\star,ctx^\star,\sigma^\star)=1
\]
and \((m^\star,ctx^\star)\) was not previously asked of the signing oracle.
The formal checked advantage is the success probability of
\[
  \ec{SIG.EUF\_CMA(H, HAETAE, A).main}.
\]

\section{Final checked implication}

The strongest checked theorem emphasized by the current proof surface is the
following EasyCrypt implication.

\begin{theorem}[Checked counted-ROM implication, informal rendering]
For any EasyCrypt memory parameter, if the final exact-hyperball NMA hop
satisfies
\[
\begin{aligned}
&\Prb[\ec{UF\_NMA}(\ec{ROMInternalTranscriptPaperSimAsNMA}
  (A,\ec{ROExactHyperballPaperSimSampler}))]  \\
&\quad + \ec{rejection\_sampling\_loss\_term}
  \le
  \Prb[\ec{MLWE\_Game}(B_{\MLWE})]
  + \Prb[\ec{BimodalSelfTargetMSIS\_Game}(B_{\BMSIS})],
\end{aligned}
\]
and if
\[
  \Prb[\ec{MLWE\_Game}(B_{\MLWE})]
  \le \ec{mlwe\_hardness\_term},
\]
and
\[
  \Prb[\ec{ModuleSIS\_Game}
      (\ec{BimodalMSIS\_As\_ModuleSIS}(B_{\BMSIS}))]
  \le \ec{module\_sis\_hardness\_term},
\]
then
\[
  \Prb[\ec{SIG.EUF\_CMA(H, HAETAE, A).main}]
  \le
  \ec{haetae\_euf\_counted\_rom\_bound}.
\]
\end{theorem}

The EasyCrypt theorem name is:

\begin{lstlisting}[language=EasyCrypt]
euf_cma_security_from_checked_ro_sampling_hop_and_counted_bimodal
\end{lstlisting}

\section{Why the theorem is conditional}

The modules \(B_{\MLWE}\) and \(B_{\BMSIS}\) are declared reduction modules in
the final theorem surface.  The top-level composition lemma does not by itself
construct these solvers from \(A\).  It composes explicitly stated probability
premises.  A paper proof may describe constructed reductions, but a faithful
description of this checked theorem must say that the final composition consumes
declared reduction modules and antecedents.

\section{The counted-ROM bound}

The reduction file defines:

\begin{lstlisting}[language=EasyCrypt]
op mlwe_hardness_term : real.
op module_sis_hardness_term : real.
op bimodal_to_msis_reduction_loss_term : real = 0%r.

op signature_query_count : real = 2%r.
op hash_query_count : real = 2%r.
op challenge_support_cardinality_lower_bound : real =
  288230376151711744%r.

op counted_rom_collision_term =
  (rom_total_query_budget * rom_total_query_budget) /
    challenge_support_cardinality_lower_bound.

op counted_rom_prequery_term =
  (rom_signature_query_budget * (rom_hash_query_budget + 1%r)) /
    challenge_support_cardinality_lower_bound.

op counted_rom_min_entropy_term =
  (rom_signature_query_budget * (rom_hash_query_budget + 1%r)) /
    challenge_support_cardinality_lower_bound.
\end{lstlisting}

Thus
\[
  q_H=2,\qquad q_S=2,\qquad N=288230376151711744=2^{58}.
\]
The ROM loss is
\[
  \frac{(2+2+1)^2}{2^{58}}
  +\frac{2(2+1)}{2^{58}}
  +\frac{2(2+1)}{2^{58}}
  =
  \frac{25+6+6}{2^{58}}
  =
  \frac{37}{2^{58}}.
\]

\section{The vacuity warning}

The proof also defines a paper-style rejection/Fiat-Shamir-with-aborts loss:

\begin{lstlisting}[language=EasyCrypt]
op rejection_sampling_loss_term =
  fs_with_aborts_reprogramming_term + fs_with_aborts_min_entropy_term.
\end{lstlisting}

The checked constants imply:

\begin{lstlisting}[language=EasyCrypt]
lemma fs_with_aborts_min_entropy_termE :
  fs_with_aborts_min_entropy_term = 12%r.

lemma rejection_sampling_loss_term_ge1 :
  1%r <= rejection_sampling_loss_term.
\end{lstlisting}

Therefore a premise of the form
\[
  \Prb[G] + \ec{rejection\_sampling\_loss\_term}
  \le
  \Prb[\MLWE] + \Prb[\BMSIS]
\]
is not a useful concrete-security premise in the current constants.  The right
side is at most \(2\), while the extra loss term is already at least \(1\) and
has a min-entropy component equal to \(12\).  The checked theorem is still a
valid logical implication, but it should be cited with this warning attached.

\chapter{Pen-and-Paper Reimplementation}

\section{Proof skeleton}

A paper proof reconstructed from the checked files should be written as the
following chain.

\begin{enumerate}
\item Define the real EUF-CMA game \(G_0\).
\item Replace real signing by a simulator while preserving the visible
distribution or charging an explicit signing-hop loss.
\item Expose transcript state so that commitments, challenges, responses, and
oracle queries become explicit.
\item Move from adaptive signing access to an NMA-style internal transcript
surface.
\item Program the random oracle and charge collision, prequery, and
min-entropy bad events.
\item Replace structural or paper sampling views by the exact hyperball sampler
used by the formal model.
\item Apply special soundness to compatible accepting transcript pairs.
\item Convert the extracted bimodal self-target relation to the Module-SIS
hardness game.
\item Compose MLWE, Module-SIS, bimodal-to-MSIS, and counted-ROM terms.
\end{enumerate}

Every step must be labelled either as an exact equivalence, a bad-event bound,
a sampling-distance bound, a reduction step, or a real-arithmetic
normalization.  Do not add rejection, exact-sampling, or active-signing losses
as final summands if the EasyCrypt theorem already internalizes them into hop
premises.  Doing so double-counts proof obligations.

\section{Game notation}

Use the following names in the paper proof.

\begin{longtable}{p{0.28\textwidth}p{0.62\textwidth}}
\toprule
\textbf{Paper object} & \textbf{EasyCrypt counterpart} \\
\midrule
\endfirsthead
\toprule
\textbf{Paper object} & \textbf{EasyCrypt counterpart} \\
\midrule
\endhead
Real EUF-CMA game &
\ec{SIG.EUF\_CMA(H, HAETAE, A).main}. \\
\midrule
No-message game &
\ec{SIG.UF\_NMA(H, HAETAE, A).main}. \\
\midrule
Random oracle &
\ec{HAETAE\_RO.FRO} cloned from the generic full random oracle. \\
\midrule
Signing oracle &
Modules named \ec{O.sign} inside transcript and budgeted NMA games. \\
\midrule
MLWE reduction game &
\ec{MLWE\_Game(Bmlwe).main}. \\
\midrule
Bimodal self-target MSIS game &
\ec{BimodalSelfTargetMSIS\_Game(Bbimodal).main}. \\
\midrule
Module-SIS game &
\ec{ModuleSIS\_Game(BimodalMSIS\_As\_ModuleSIS(Bbimodal)).main}. \\
\bottomrule
\end{longtable}

\section{Bad events}

The random-oracle and sampler bad events should be separated.

\begin{itemize}
\item The counted-ROM collision term accounts for collisions among programmed
or observed ROM sites.
\item The counted-ROM prequery term accounts for programmed challenge sites
that were already queried in the counted transcript ledger.
\item The counted-ROM min-entropy term accounts for the challenge entropy bad
event in the counted transcript ledger.
\item The sampler prequery event \(\ec{sampler\_bad\_prequery}\) is a local
active-signing event.  It is bounded by a different expression involving the
adversary hash log and the sampler self-log.
\end{itemize}

The paper proof must not identify \(\ec{sampler\_bad\_prequery}\) with
\(\ec{counted\_rom\_prequery\_term}\).  They are related proof ideas, but they
are different events in the checked source.

\section{Special soundness and extraction}

The extraction step is not a conclusion from a single valid transcript.  It
requires a forking pair: two accepting transcripts sharing the appropriate
public data and carrying distinct challenges.  The paper proof should state the
obligation as:
\[
  \mathsf{ValidForkingPair}(tr_1,tr_2)
  \Longrightarrow
  \mathsf{Solve}_{\BMSIS/\MSIS}(tr_1,tr_2).
\]

The checked anchors are:

\begin{lstlisting}[language=EasyCrypt]
extract_bimodal_forking_some
extract_bimodal_forking_nonempty
valid_forking_pair_public_equations
bimodal_to_module_sis_valid
bimodal_solver_lift_exact
bimodal_mode_solver_lifted_distribution_exact
bimodal_msis_as_module_sis_exact
\end{lstlisting}

\chapter{EasyCrypt Reimplementation Plan}

\section{Development order}

Reimplement the files from leaves to root.  Do not begin with the final theorem.
Use the following order.

\begin{enumerate}
\item \file{HAETAE_Params.ec}
\item \file{HAETAE_Keccak1600.ec}
\item \file{HAETAE_FIPS202.ec}
\item \file{HAETAE_Algebra.ec}
\item \file{HAETAE_Distributions.ec}
\item \file{HAETAE_Events.ec}
\item \file{HAETAE_Assumptions.ec}
\item \file{Sig_ROM.ec}
\item \file{HAETAE_ROM.ec}
\item \file{HAETAE_Scheme.ec}
\item \file{HAETAE_Transcript.ec}
\item \file{HAETAE_Reductions.ec}
\item \file{HAETAE_Rejection.ec}
\item \file{HAETAE_ROM_Programming.ec}
\item \file{HAETAE_HopGames.ec}
\item \file{HAETAE_Security.ec}
\end{enumerate}

The manifest compiles root-first because each file is checked independently
with imports resolved by EasyCrypt.  Human reimplementation should proceed
leaf-first because each later file uses definitions from earlier ones.

\section{Global proof idioms}

Use four recurring EasyCrypt idioms.

\paragraph{Typed interfaces.}
Use module types for adversaries, schemes, oracles, and hardness solvers.  This
prevents games from depending on implementation details that should remain
abstract.

\paragraph{State exposure.}
When a paper proof says ``the simulator records queries'', make a module with
explicit state variables for logs, counters, and bad flags.

\paragraph{Local lemmas before global theorem.}
First prove preservation lemmas, support lemmas, losslessness facts, and
pointwise arithmetic bounds.  Then assemble them in top-level security lemmas.

\paragraph{No silent axioms.}
Do not use \ec{axiom}, \ec{admit.}, or \ec{sorry.}.  If a fact is an external
assumption, represent it as a declared operation or module interface and make
the final theorem conditional on it.

\section{Placeholder convention}

The snippets in the next chapter are intentionally schematic where the full
source would obscure the proof architecture.  In particular, code fragments
containing \ec{...} are not EasyCrypt programs.  They mark locations where the
complete project source must be consulted or rederived.  This convention is a
guardrail against a false kind of self-containment: copying the snippets alone
will not produce a compiling proof, and the report should not be cited as a
literal replacement for the checked \file{.ec} files.

For a genuinely standalone manual, every occurrence of \ec{...} would have to be
expanded into exact EasyCrypt source, every helper lemma would need its full
statement and proof, and every module declaration would need its exact
restriction set.  That is a different artifact from this roadmap.

Before claiming that a documentation-derived proof tree is machine checked, run
a placeholder audit over the generated files.  The audit is deliberately simple:
any remaining ellipsis or proof-hole marker must be explained as a comment or
removed from the proof surface.

\begin{lstlisting}[language=bash]
rg -n "\.\.\.|TODO|FIXME|admit\.|sorry\.|\baxiom\b" \
  provable-security/easycrypt
\end{lstlisting}

The expected result for the active proof surface is no unchecked proof hole.  A
hit on \ec{...} in executable EasyCrypt means the roadmap has not yet been
expanded into source.

\chapter{File-by-File Blueprint}

\section{HAETAE\_Params.ec}

Define the fixed HAETAE parameter vocabulary.

\begin{lstlisting}[language=EasyCrypt]
type mode = [ Mode2 | Mode3 | Mode5 ].

op n : int = 256.
op q : int = 64513.
op dq : int = 2 * q.

op mode_k (md : mode) : int = ...
op mode_l (md : mode) : int = ...
op mode_tau (md : mode) : int = ...
\end{lstlisting}

Required lemmas:

\begin{itemize}
\item positivity of byte sizes, \(n\), \(q\), and all mode-dependent bounds;
\item \(dq=2q=129026\);
\item \(\tau\le n\) for every mode;
\item positivity of public-key and signature-size constants.
\end{itemize}

These lemmas are simple case splits, but they are essential because later
probability and support facts need positive denominators and nonempty domains.

\section{HAETAE\_Keccak1600.ec and HAETAE\_FIPS202.ec}

These files provide modeled hash-interface support.  They do not replace the
random oracle in the security theorem.  A reimplementation should define:

\begin{itemize}
\item byte and state list types;
\item state-size, rate, capacity, domain-separator, and padding constants;
\item deterministic absorb, squeeze, and permutation operations;
\item length and well-formedness lemmas needed by the algebra layer.
\end{itemize}

Keep the cryptographic boundary explicit: compiling these files does not prove
standard-model SHAKE security.

\section{HAETAE\_Algebra.ec}

This is the algebraic model of HAETAE objects.  Define carrier types for:

\begin{itemize}
\item bytes, seeds, CRHs, random coins, messages, and contexts;
\item coefficients, polynomials, challenges, polynomial vectors, and matrices;
\item public keys, secret keys, and signatures.
\end{itemize}

Then define total operations for:

\begin{itemize}
\item modular coefficient arithmetic;
\item polynomial and vector operations;
\item key generation internals;
\item signing internals;
\item verification internals;
\item transcript commitment and challenge components;
\item norm and validity predicates.
\end{itemize}

The most important final facts are:

\begin{lstlisting}[language=EasyCrypt]
lemma verify_internal_sign_internal ...
lemma valid_signature_sign_internal ...
lemma secret_norm_ok_current ...
lemma response_norm_ok_current ...
lemma verify_norm_ok_current ...
\end{lstlisting}

The proof model stores squared norm thresholds.  The validity reading is:
\[
  \|z\|^2 \le B^2
  \quad\text{implemented as}\quad
  \ec{normsq}(z)\le \ec{mode\_b*sq}(md).
\]
Do not replace this by a strict inequality unless the corresponding rejection
predicate actually uses a strict bound.

\section{HAETAE\_Distributions.ec}

Define the distributions used by key generation, signing, challenges, and
structural samplers.

Essential distribution properties:

\begin{itemize}
\item losslessness of byte, seed, CRH, random-coin, coefficient, polynomial,
vector, matrix, and challenge distributions;
\item support facts showing sampled objects satisfy the predicates expected by
the algebra layer;
\item projection facts connecting structural signing samples to response,
challenge, and randomness components.
\end{itemize}

Every sampler used by a program that should terminate with probability one must
have a corresponding losslessness lemma or a local proof of losslessness.

\section{HAETAE\_Events.ec}

Define named event predicates rather than spelling conditions repeatedly.
Events should include rejection, bad prequery, bad min-entropy, collision, and
validity conditions needed by the game hops.

This file should contain no heavy proof.  Its job is vocabulary alignment:
later files can reason about the same event name instead of duplicating
Boolean formulas.

\section{HAETAE\_Assumptions.ec}

Introduce the external hardness games and reduction interfaces.

\begin{lstlisting}[language=EasyCrypt]
module type MLWE_Reduction = {
  proc main() : bool
}.

module type BimodalSelfTargetMSIS_Reduction = {
  proc main() : bool
}.

module type ModuleSIS_Reduction = {
  proc main() : bool
}.
\end{lstlisting}

Define game wrappers that call these modules:

\begin{lstlisting}[language=EasyCrypt]
module MLWE_Game(B : MLWE_Reduction) = {
  proc main() : bool = {
    var b : bool;
    b <@ B.main();
    return b;
  }
}.
\end{lstlisting}

Also define the bimodal-to-Module-SIS adapter and prove exactness lemmas:

\begin{lstlisting}[language=EasyCrypt]
bimodal_to_module_sis_valid
bimodal_solver_lift_exact
bimodal_mode_solver_lifted_distribution_exact
bimodal_msis_as_module_sis_exact
\end{lstlisting}

These are not concrete cryptanalytic estimates.  They are formal assumption
surfaces used by the final reduction theorem.

\section{Sig\_ROM.ec}

Define the generic random-oracle signature game independent of HAETAE.

\begin{lstlisting}[language=EasyCrypt]
type pkey.
type skey.
type message.
type context.
type signature.
type ro_query.
type ro_output.

op [lossless] dro_output : ro_query -> ro_output distr.
\end{lstlisting}

Define module types:

\begin{lstlisting}[language=EasyCrypt]
module type Oracle = { include FRO [init, get] }.
module type POracle = { include FRO [get] }.
module type Scheme(O : POracle) = { ... }.
module type SignOracle = { proc sign(m : message, ctx : context) : signature }.
module type Adversary(H : POracle, O : SignOracle) = { ... }.
module type NMA_Adversary(H : POracle) = { ... }.
\end{lstlisting}

Define:

\begin{itemize}
\item correctness game;
\item EUF-CMA game;
\item UF-NMA game;
\item NMA-as-CMA adapter;
\item freshness predicates for signed message-context pairs.
\end{itemize}

The generic game should not contain HAETAE-specific facts.  It must be reusable.

\section{HAETAE\_ROM.ec}

Instantiate random-oracle query and output types for HAETAE.  Use typed query
constructors for distinct hash domains:

\begin{itemize}
\item matrix expansion;
\item message hashing;
\item sampler expansion;
\item challenge hashing.
\end{itemize}

The reason for typed constructors is domain separation.  A proof should not be
able to confuse a sampler-expansion query with a challenge query.

\section{HAETAE\_Scheme.ec}

Clone the generic signature interface and instantiate it with HAETAE types and
the HAETAE random-oracle output distribution:

\begin{lstlisting}[language=EasyCrypt]
clone import Sig_ROM as SIG with
  type pkey <- pkey,
  type skey <- skey,
  ...
  op dro_output <- dro_output.
\end{lstlisting}

Define the formal scheme module:

\begin{lstlisting}[language=EasyCrypt]
module HAETAE(H : SIG.POracle) = {
  proc kg() : pkey * skey = { ... }
  proc sign(sk : skey, m : message, ctx : context) : signature = { ... }
  proc verify(pk : pkey, m : message, ctx : context,
              sig : signature) : bool = { ... }
}.
\end{lstlisting}

The current wrapper makes ROM calls visible around deterministic internal
operations.  In particular, signing samples typed random coins, asks the sampler
ROM query, interprets the output as signing coins, computes message and
challenge ROM queries, and finally calls \ec{sign\_internal}.

\section{HAETAE\_Transcript.ec}

Define transcript records and validity predicates.  A transcript should contain
at least the public key, message, context, signature, commitment components,
challenge query, and challenge value needed by extraction.

Key proof obligations:

\begin{lstlisting}[language=EasyCrypt]
transcript_valid_signature_challenge_query
transcript_valid_signature_challenge_matches
transcript_valid_signature_challengeE
extract_bimodal_forking_some
extract_bimodal_forking_nonempty
valid_forking_pair_public_equations
\end{lstlisting}

The extraction lemmas should require a valid forking pair.  Do not prove or
state special soundness from one accepting transcript alone.

\section{HAETAE\_Reductions.ec}

Define all real-valued loss terms and prove their arithmetic properties.

Required groups:

\begin{itemize}
\item abstract hardness terms:
\(\ec{mlwe\_hardness\_term}\),
\(\ec{module\_sis\_hardness\_term}\);
\item bimodal-to-MSIS loss slot, currently zero;
\item fixed query counts \(q_H=q_S=2\);
\item challenge support lower bound \(N=2^{58}\);
\item counted-ROM collision, prequery, and min-entropy terms;
\item paper-style rejection/Fiat-Shamir-with-aborts terms;
\item grouped final bounds
\(\ec{haetae\_euf\_bound}\) and
\(\ec{haetae\_euf\_counted\_rom\_bound}\).
\end{itemize}

Arithmetic lemmas to reprove:

\begin{lstlisting}[language=EasyCrypt]
haetae_euf_counted_rom_bound_groupedE
counted_rom_collision_termE
counted_rom_prequery_termE
counted_rom_min_entropy_termE
counted_rom_programming_loss_termE
counted_rom_programming_loss_term_concreteE
counted_rom_programming_loss_term_subunit
fs_with_aborts_min_entropy_termE
rejection_sampling_loss_term_ge1
bimodal_to_msis_reduction_loss_termE
\end{lstlisting}

This file is where most accidental theorem inflation can be caught.  The final
counted-ROM bound is not the same as the older paper-style rejection bound.

\section{HAETAE\_Rejection.ec}

Define signing-attempt state, rejection predicates, and the relationship
between reference rejection and the modeled signing state.

Important predicates:

\begin{itemize}
\item \(\ec{signing\_attempt\_state\_valid\_signature\_accepts}\);
\item \(\ec{signing\_attempt\_state\_reject1\_accepts}\);
\item \(\ec{signing\_attempt\_state\_reject2\_accepts}\);
\item \(\ec{signing\_attempt\_state\_reference\_rejection\_accepts}\);
\item \(\ec{signing\_attempt\_state\_rejection\_accepts}\).
\end{itemize}

Important obligations:

\begin{itemize}
\item bound rejection probabilities by named terms;
\item show exact-hyperball signing attempt states satisfy the reference
rejection predicates;
\item prove zero-abort facts for the exact paths currently used by the model;
\item expose coarse fallback bounds honestly when they are used.
\end{itemize}

\section{HAETAE\_ROM\_Programming.ec}

Define the counted random-oracle programming surface.

This file should model:

\begin{itemize}
\item lazy ROM maps;
\item observed query logs;
\item programmed query sites;
\item collision and prequery bad events;
\item preservation of good state under oracle operations;
\item probability bounds for programming failures.
\end{itemize}

Use explicit counters and support lower bounds.  Do not write a prose
``fresh with high probability'' argument without a program variable tracking the
queries that make the statement meaningful.

\section{HAETAE\_HopGames.ec}

This is the largest file because it contains the concrete hybrid machinery.
Reimplement it in layers.

\subsection*{Layer 1: concrete lazy-ROM laws}

Prove sampler expansion freshness:

\begin{lstlisting}[language=EasyCrypt]
concrete_lazy_rom_sampler_expand_query_fresh_le
concrete_lazy_rom_sampler_expand_query_fresh_eq
\end{lstlisting}

Mathematical reading:
if \(\ec{sampler\_expand\_query}(r)\) is not in the lazy-ROM table, then reading
the ROM at that point yields a signing-coin distribution bounded by the ideal
coin distribution.

\subsection*{Layer 2: signing-oracle clean invariants}

Expose the active signing oracle state and prove that clean conditions are
preserved:

\begin{lstlisting}[language=EasyCrypt]
concrete_budgeted_ro_signing_attempt_o_sign_clean_sampler_fresh
\end{lstlisting}

The clean condition should include sampler-ROM coverage, absence of the
current sampler prequery bad event, budget discipline, public-key consistency,
and secret-key consistency.

\subsection*{Layer 3: one-call lifting}

Prove the clean one-call bound:

\begin{lstlisting}[language=EasyCrypt]
concrete_budgeted_o_sign_clean_one_call_loss_from_self_logged_surface
\end{lstlisting}

Its mathematical shape is:
\[
\begin{aligned}
&\Prb[O_{\mathsf{attempt}}.\mathsf{sign}:p(res)\wedge
  \neg\mathsf{sampler\_bad\_prequery}]\\
&\qquad\le
  \Prb[O_{\mathsf{exact}}.\mathsf{sign}:p(res)]
  + \ec{rejection\_sampling\_loss\_term}.
\end{aligned}
\]

\subsection*{Layer 4: NMA and transcript game lifting}

Lift the one-call facts through NMA-style transcript games:

\begin{lstlisting}[language=EasyCrypt]
rom_internal_budgeted_nma_ro_signing_attempt_ro_exact_hyperball_loss_from_paper_sample_lifting
rom_internal_budgeted_nma_ro_signing_attempt_ro_exact_hyperball_clean_conditioned_lifting
\end{lstlisting}

This is where the current development still contains coarse fallback bounds.
A faithful reimplementation should keep those fallbacks visible and should not
describe the final result as non-vacuous.

\section{HAETAE\_Security.ec}

Assemble the final theorem family.

First prove correctness:

\begin{lstlisting}[language=EasyCrypt]
op correctness_obligation = ...
lemma correctness_obligation_holds : correctness_obligation.
lemma correctness_perfect &m :
  Pr[SIG.Correctness(H, HAETAE, A).main() @ &m : res] = 0%r.
\end{lstlisting}

Then assemble progressively stronger security lemmas:

\begin{itemize}
\item generic EUF-CMA security from hop interfaces;
\item ROM-transcript security;
\item bimodal-to-Module-SIS lift;
\item counted-ROM grouped bound;
\item paper-sampling and exact-hyperball bridges;
\item final checked counted-ROM implication.
\end{itemize}

The final theorem should look like:

\begin{lstlisting}[language=EasyCrypt]
lemma euf_cma_security_from_checked_ro_sampling_hop_and_counted_bimodal &m :
  Pr[SIG.UF_NMA(H, HAETAE,
       ROMInternalTranscriptPaperSimAsNMA
         (A, ROExactHyperballPaperSimSampler(H))).main() @ &m : res] +
    rejection_sampling_loss_term <=
    Pr[MLWE_Game(Bmlwe).main() @ &m : res] +
    Pr[BimodalSelfTargetMSIS_Game(Bbimodal).main() @ &m : res] =>
  Pr[MLWE_Game(Bmlwe).main() @ &m : res] <= mlwe_hardness_term =>
  Pr[ModuleSIS_Game(BimodalMSIS_As_ModuleSIS(Bbimodal)).main()
       @ &m : res] <= module_sis_hardness_term =>
  Pr[SIG.EUF_CMA(H, HAETAE, A).main() @ &m : res] <=
    haetae_euf_counted_rom_bound.
\end{lstlisting}

The proof should be mostly composition: apply the previous bridge lemma, pass
the hop premise, pass the MLWE premise, and pass the Module-SIS premise.

\chapter{Reimplementation Checklists}

\section{Pen-and-paper checklist}

Before writing EasyCrypt, verify the paper proof has these properties.

\begin{enumerate}
\item The starting game is the same EUF-CMA game later encoded in
\ec{Sig\_ROM.ec}.
\item The source game has no syntactic query cap; fixed counts belong to the
wrapper and bound ledger.
\item Every hop is classified as exact, bad-event, sampling-distance, or
reduction.
\item Random-oracle programming freshness is stated in terms of explicit query
sets.
\item Sampler-prequery events are separated from counted-ROM challenge
prequery events.
\item Forking/extraction uses two compatible accepting transcripts with
distinct challenges.
\item MLWE and Module-SIS hardness are external assumptions represented by
abstract games.
\item The final counted-ROM bound is not mixed with the older
\ec{haetae\_euf\_bound} rejection ledger.
\item The final theorem is marked conditional and carries the vacuity warning.
\end{enumerate}

\section{EasyCrypt checklist}

Before accepting a reimplementation, run these checks.

\begin{enumerate}
\item The reproduction mode is declared: audit mode with companion sources,
literal standalone mode with full source in the document, or independent
reimplementation mode with a traceability matrix.
\item If \file{.ec} scripts are supplied in script companion mode, every
security-relevant script block has a corresponding mathematical explanation in
the report.
\item The script-companion coverage ledger is complete: no manifest file,
security-critical declaration, nontrivial lemma, tactic block, module
restriction, or external assumption remains without a prose explanation.
\item Every manifest file compiles.
\item There are no \ec{axiom}, \ec{admit.}, or \ec{sorry.} tokens in the proof
surface.
\item There are no executable \ec{...}, \ec{TODO}, or \ec{FIXME} placeholders in
the proof surface.
\item Each declared module in \file{HAETAE_Security.ec} has the right
restriction set, so the adversary cannot call hidden game internals.
\item All distributions sampled by lossless procedures have losslessness
lemmas.
\item All denominators in real bounds have positivity lemmas.
\item The counted-ROM concrete arithmetic rewrites to \(37/2^{58}\).
\item The rejection min-entropy term rewrites to \(12\).
\item The final theorem consumes declared reduction modules and explicit
premises rather than pretending to construct all solvers inside the final
composition.
\item Every theorem, loss term, and declared module used by the final theorem is
traceable either to the companion source or to a checked object in the new
reimplementation.
\end{enumerate}

\section{Module restriction-set caveat}

Module restriction sets are part of the proof, not formatting.  They prevent an
adversary or oracle module from calling game internals that should be hidden.
The checklist item above therefore requires exact source-level comparison, not a
verbal approximation.

For example, the final mechanized ROM-transcript/NMA section begins with
declarations of this form:

\begin{lstlisting}[language=EasyCrypt]
section MechanizedROMTranscriptConstructedNMALift.

declare module H <: SIG.Oracle {-SIG.EUF_CMA,
                                 -EUF_CMA_InternalTranscriptSign,
                                 -EUF_CMA_ROMInternalTranscriptSign,
                                 -ROMInternalTranscriptAsNMA,
                                 -ROMInternalTranscriptPublicSimAsNMA,
                                 -ROMInternalTranscriptPaperSimAsNMA,
                                 -ROMInternalTranscriptBudgetedPaperSimAsNMA,
                                 -ROMPaperSimSigningSampler,
                                 -RealSigningPaperSimSampler,
                                 -ROSigningAttemptPaperSimSampler,
                                 -HAETAERejectionPaperSimSampler,
                                 -ExactHyperballHAETAERejectionPaperSimSampler,
                                 -ExactHyperballPaperSimSampler,
                                 -ROExactHyperballPaperSimSampler}.
declare module A <: SIG.Adversary {-H, -HAETAE, -SIG.EUF_CMA,
                                    -EUF_CMA_InternalTranscriptSign,
                                    -EUF_CMA_ROMInternalTranscriptSign,
                                    -ROMInternalTranscriptAsNMA,
                                    -ROMInternalTranscriptPublicSimAsNMA,
                                    -ROMInternalTranscriptPaperSimAsNMA,
                                    -ROMInternalTranscriptBudgetedPaperSimAsNMA,
                                    -ROMPaperSimSigningSampler,
                                    -RealSigningPaperSimSampler,
                                    -ROSigningAttemptPaperSimSampler,
                                    -HAETAERejectionPaperSimSampler,
                                    -ExactHyperballHAETAERejectionPaperSimSampler,
                                    -ExactHyperballPaperSimSampler,
                                    -ROExactHyperballPaperSimSampler}.
declare module Samp <: PaperSimSigningSampler {-H, -A,
                                               -ROMInternalTranscriptPaperSimAsNMA}.
declare module Bmlwe <: MLWE_Reduction.
declare module Bbimodal <: BimodalSelfTargetMSIS_Reduction.
\end{lstlisting}

This listing is representative, not exhaustive.  A complete reimplementation
must audit every \ec{declare module} line in \file{HAETAE\_Security.ec} and in
the imported hop-game files.  Omitting one restriction can turn a meaningful
game proof into a circular proof in which an adversary has access to hidden
state or procedures.

\section{Common failure modes}

\paragraph{Overstating the theorem.}
Do not write ``the EasyCrypt proof proves HAETAE EUF-CMA security'' without the
modifiers ``formal model'', ``random oracle'', ``conditional'', and ``fixed
ledger''.

\paragraph{Counting a loss twice.}
If a loss is already internalized into a hop premise or bridge lemma, do not add
it as a separate final summand.

\paragraph{Erasing the query-bound boundary.}
The constants \(q_H=q_S=2\) are fixed in the current ledger.  They are not a
symbolic arbitrary-query theorem.

\paragraph{Confusing SHAKE support with ROM security.}
\file{HAETAE_FIPS202.ec} and \file{HAETAE_Keccak1600.ec} are modeled hash
support files.  They do not prove that SHAKE instantiates the random oracle.

\paragraph{Using proof holes.}
If a fact is not proved, make it a premise or an explicitly documented
assumption.  Do not hide it behind an axiom in the manifest proof surface.

\chapter{Minimal Build and Verification Artifacts}

\section{Verification script template}

A reimplementation can use this shell skeleton.

\begin{lstlisting}[language=bash]
#!/usr/bin/env sh
set -eu

EASYCRYPT=${EASYCRYPT:-easycrypt}
ROOT=$(pwd)
PROOF_DIR="$ROOT/provable-security/easycrypt"
MANIFEST="$ROOT/provable-security/proof-files.txt"
LOG_DIR="$ROOT/provable-security/logs"

mkdir -p "$LOG_DIR"

failed=0
while IFS= read -r file; do
  case "$file" in
    ""|\#*) continue ;;
  esac

  echo "=== $file ==="
  if "$EASYCRYPT" compile "$PROOF_DIR/$file" \
      -I "$PROOF_DIR" \
      -I "$ROOT/kyber-security" \
      > "$LOG_DIR/$file.log" 2>&1; then
    echo "PASS $file"
  else
    echo "FAIL $file"
    failed=1
  fi
done < "$MANIFEST"

if [ "$failed" -eq 0 ]; then
  echo "RESULT: PASS"
else
  echo "RESULT: FAIL"
  exit 1
fi
\end{lstlisting}

\section{Documentation build}

This report can be built from the \file{report/} directory with:

\begin{lstlisting}[language=bash]
make
\end{lstlisting}

or directly with:

\begin{lstlisting}[language=bash]
pdflatex -interaction=nonstopmode -halt-on-error \
  haetae-provable-security-reimplementation-guide
pdflatex -interaction=nonstopmode -halt-on-error \
  haetae-provable-security-reimplementation-guide
\end{lstlisting}

\chapter{Final Acceptance Criteria}

A reimplementation should not be considered complete until all criteria below
hold.

\begin{checkpoint}
The reimplementation is checked against the companion EasyCrypt source files or
against an independently expanded document that contains complete source, full
lemma statements, proof scripts, and exact module restriction sets.  This report
alone is not treated as sufficient source material.
\end{checkpoint}

\begin{checkpoint}
The report user records which reproduction mode is being claimed.  A claim of
literal standalone reproduction is rejected unless the documentation has been
expanded beyond this roadmap into a complete source-and-proof archive.
\end{checkpoint}

\begin{checkpoint}
Every schematic snippet copied from this report has been replaced by checked
EasyCrypt source.  The active proof surface contains no executable ellipsis,
\ec{TODO}, \ec{FIXME}, \ec{admit.}, \ec{sorry.}, or hidden \ec{axiom}.
\end{checkpoint}

\begin{checkpoint}
The reimplementation includes a traceability matrix mapping each final-theorem
dependency to either an exact companion-source declaration or an independently
checked replacement.
\end{checkpoint}

\begin{checkpoint}
The manifest-driven EasyCrypt proof gate terminates with \ec{RESULT: PASS}.
\end{checkpoint}

\begin{checkpoint}
The declaration inventory has no axioms and no admitted proofs in the active
proof surface.
\end{checkpoint}

\begin{checkpoint}
The final theorem statement has the same logical shape as
\ec{euf\_cma\_security\_from\_checked\_ro\_sampling\_hop\_and\_counted\_bimodal}.
\end{checkpoint}

\begin{checkpoint}
The documentation states that the theorem is conditional, fixed-ledger,
formal-model, and random-oracle-model.
\end{checkpoint}

\begin{checkpoint}
The documentation distinguishes understanding the checked formal-model theorem
from understanding or certifying end-to-end security of the HAETAE digital
signature algorithm.
\end{checkpoint}

\begin{checkpoint}
The documentation gives a clear negative answer to any claim that the report
alone suffices both to understand HAETAE algorithm security and to reproduce the
machine-checked EasyCrypt verification.
\end{checkpoint}

\begin{checkpoint}
Before any completed script-companion claim is made, the documentation is
complete enough for mathematicians and cryptologists to understand the
mathematical content of the provided \file{.ec} scripts from the report:
declarations, state variables, module restrictions, lemma statements,
proof-script blocks, tactic roles, assumptions, and global security use are all
explained.
\end{checkpoint}

\begin{checkpoint}
Before any completed script-companion claim is made, the script-companion
coverage ledger has no open rows.  Each manifest file and each
security-relevant declaration, module, lemma, proof block, tactic pattern,
module restriction, and external assumption is linked to a mathematical
explanation in the report.
\end{checkpoint}

\begin{checkpoint}
The documentation states that the current final-hop premise is not a
non-vacuous concrete-security premise because of
\ec{rejection\_sampling\_loss\_term}.
\end{checkpoint}

\begin{checkpoint}
Any claim of standalone documentation sufficiency is backed by literal
EasyCrypt source, complete proof scripts, exact module restriction sets, full
helper-lemma statements, an import/library contract, and a successful
verification transcript.
\end{checkpoint}

\begin{checkpoint}
The model/specification bridge is recorded as an audited or assumed boundary,
not as a checked equivalence theorem.
\end{checkpoint}

\begin{checkpoint}
The proof does not claim concrete SHAKE instantiation, implementation
correctness, concrete lattice bit security, or reduction runtime bounds.
\end{checkpoint}

\chapter{Summary}

The HAETAE EasyCrypt proof is best reimplemented as a layered proof surface,
using this report as the map and the checked EasyCrypt files as the authoritative
source.  First define parameters, algebra, distributions, events, and assumption
games.  Then define generic ROM signature games and instantiate them with the
HAETAE scheme model.  Next expose transcript state, rejection state,
random-oracle programming state, and budgeted signing-oracle invariants.
Finally assemble the top-level theorem by composing checked hop premises,
declared hardness bounds, and the counted-ROM arithmetic ledger.

The central discipline is not merely to make EasyCrypt accept a long file.  The
theorem must say exactly what has been proved.  For the current artifact, that
is a conditional probability implication for the formal HAETAE random-oracle
model, with a fixed counted-ROM ledger and an explicit vacuity warning on the
final rejection-sampling premise.  When the document-only question is asked, the
answer is no: this report is not, by itself, a proof of deployed algorithm
security, a SHAKE-instantiation theorem, or a standalone machine-checkable
archive.  A faithful reimplementation preserves that boundary while making
every program, event, loss, and reduction interface visible to the checker.  The
report is therefore an entry point for understanding and auditing the proof; the
proof itself remains the checked EasyCrypt source plus its manifest-driven
verification run.

If the \file{.ec} scripts are distributed together with this document, the
documentation burden increases.  The report must then serve as a self-contained
mathematical reading guide to those scripts: every security-relevant definition,
game, invariant, module restriction, lemma, tactic block, and assumption must
have a documented mathematical interpretation.  That still does not make the
prose a substitute for machine checking, but it should make the scripts
intelligible to a mathematician or cryptologist without relying on unwritten
project context.  The present version records this obligation and supplies a
partial guide; it becomes a completed script companion only after the coverage
ledger has no open rows.

\end{document}
